% PDF/A-1 è PDF 1.4. Con TeX Live 2026 pdfx non riesce più a impostarlo da sé
% (il PDF uscirebbe 1.7): va detto prima della classe, come chiede pdfx.
\pdfminorversion=4
\documentclass{theme/uniprthesis}

% PDF/A-1b non ammette gruppi di trasparenza, e pdftex.def (dal 2024) ne mette
% uno su ogni PDF incluso. L'unico PDF incluso è il logo, senza trasparenza.
\makeatletter
\def\Gread@pdf#1{\let\Gin@decode\@empty\Gin@interpolatefalse\Gread@transgroupfalse\Gread@pdftex{#1}}
\makeatother

\usepackage[italian]{babel}
\usepackage{todonotes}
\usepackage{ifdraft}
\usepackage{afterpage}
\usepackage{amsmath}
\usepackage{amssymb}
\usepackage{amsfonts}
\usepackage{bbold}
\usepackage{footmisc}
\usepackage{xcolor}
\usepackage{ragged2e}
\usepackage{listings}
\usepackage{multicol}
\usepackage{verbatim}
\usepackage{xurl}
\usepackage{booktabs}
\usepackage{array}
\usepackage{longtable}
\usepackage{pgfplots}
\usepackage{pgfplotstable}
\pgfplotsset{compat=1.18}
\usepgfplotslibrary{groupplots}
\usetikzlibrary{arrows.meta,positioning,shapes.geometric,fit,calc,backgrounds,patterns}
\pgfplotsset{tesi/.style={
  tick label style={font=\footnotesize}, /pgf/number format/use comma, /pgf/number format/1000 sep={.}, label style={font=\footnotesize},
  legend style={font=\footnotesize, draw=none, fill=none},
  grid=major, grid style={line width=.2pt, draw=gray!35}, every axis plot/.append style={line width=.9pt}}}

\emergencystretch=1.5em
\usepackage{csquotes}
\usepackage[htt]{hyphenat}
\usepackage[backend=biber, style=numeric-comp, sorting=nyt, giveninits=true, maxbibnames=10]{biblatex}
\addbibresource{bibliografia.bib}

\lstset{
  basicstyle=\ttfamily\small,
  keywordstyle=\bfseries\color{blue},
  commentstyle=\color{gray},
  stringstyle=\color{red},
  showstringspaces=false,
  breaklines=true,
  frame=single,
  numbers=left,
  numberstyle=\tiny\color{gray},
}

\title{Recupero ibrido lessicale e vettoriale in un sistema di gestione documentale}
\author{Martin Trajkovski}
\advisor{Prof. Alessandro Dal Palù}
\college{Dipartimento di Scienze Matematiche, Fisiche e Informatiche}
\degree{Corso di Laurea Magistrale in Scienze Informatiche}
\degreeyears{2026--2027}
\serialNumber{397464}

\newcommand{\coadvisor}{}
% Segnaposto dei numeri che arrivano con le known-item umane complete.
\newcommand{\dacompletare}{\textcolor{red}{\textbf{??}}}

\newcommand\blankpage{
    \null
    \thispagestyle{empty}
    \addtocounter{page}{-1}
    \newpage
}

\begin{document}

\maketitle
\blankpage

\newpage
\thispagestyle{empty}
\null\vspace{\stretch{1}}
\begin{flushright}
	\textit{The first principle is that you must not fool yourself - and you are the easiest person to fool.}\\[1em]
	Richard P. Feynman, \textit{Cargo Cult Science}, 1974
\end{flushright}
\vspace{\stretch{3}}\null
\newpage
\blankpage

\newpage
\thispagestyle{empty}
\null\vspace{\stretch{1}}
\begin{flushright}
	\textit{Dedica da scrivere.}
\end{flushright}
\vspace{\stretch{3}}\null
\newpage
\blankpage

\pagenumbering{arabic}
\tableofcontents
\listoffigures

\chapter{Introduzione}\label{chapter:introduzione}

\section{Ritrovare un atto in un documentale}\label{sec:ritrovare}

Un sistema documentale è l'archivio di un'organizzazione: determine, delibere,
contratti, verbali, archiviati con un nome, qualche etichetta e un riassunto. Chi
lo usa di rado sfoglia: cerca. E cerca quasi sempre qualcosa che sa esistere,
l'ordinanza di cui ricorda il numero, la determina che affidava un lavoro di cui
ricorda l'argomento. Se la ricerca non la trova, per quella persona l'atto non
c'è.

Le ricerche di un documentale sono di due specie. Alcune sono identificative:
\emph{ordinanza 187}, \emph{determina 1223 Crispiano}, dove conta un numero e
sbagliarlo di una cifra vuol dire un altro atto. Altre sono in lingua naturale:
\emph{manutenzione strade}, \emph{contributi alle famiglie per l'infanzia}, dove
non c'è una sola risposta giusta e le parole della query possono non stare
nell'atto. Un motore di ricerca che va bene per le une non va bene, per forza,
per le altre, e questa tesi misura quanto.

\section{Koskidex e Documentale}\label{sec:due-sistemi}

\textbf{Documentale} è un sistema documentale sviluppato da Dieffetech, con un
backend Laravel, MySQL ed Elasticsearch per la ricerca. \textbf{Koskidex} è un
motore di ricerca full-text scritto in Go dall'autore di questa tesi, nato per la
barra di ricerca di un piccolo negozio online: un binario solo, senza servizi
esterni, con una dipendenza di codice. La tesi mette Koskidex al posto di
Elasticsearch dentro Documentale, e chiede cosa costa.

L'innesto in sé non è il contributo, e sarebbe solo lavoro di integrazione. Il
contributo è quello che l'innesto obbliga a guardare. Per sostituire
Elasticsearch bisogna prima restituire gli stessi risultati, e arrivarci ha
richiesto di trovare e riprodurre ogni differenza di matching fra i due motori.
Da lì, ogni correzione si misura contro due riferimenti: il codice di Koskidex
com'era il 23 settembre 2026, congelato in un test, ed Elasticsearch come lo
interroga Documentale in produzione.

\section{La domanda di ricerca}\label{sec:domanda}

\begin{quote}
  \emph{Quali difetti di recupero ha un motore di ricerca piccolo e senza
  servizi esterni quando prende il posto di Elasticsearch in un sistema documentale,
  quanto costa ciascuno in pertinenza e in tempo, e quanto se ne recupera
  correggendoli uno alla volta?}
\end{quote}

La domanda si articola in quattro:
\begin{enumerate}
  \item quanto pesano, misurati, i difetti del recupero di Koskidex: i tre che la
        lettura del codice fa sospettare, l'assenza di IDF, un vettore che
        riordina senza recuperare e una fusione che somma scale incomparabili, e
        il recupero congiuntivo, che la prima misura ha messo davanti a tutti (capitoli~\ref{chapter:lessicale} e~\ref{chapter:ibrido});
  \item quali difetti ha la ricerca di Documentale in produzione, che l'innesto
        porta alla luce (capitolo~\ref{chapter:produzione});
  \item se esiste una configurazione che vada bene per le ricerche identificative
        e per quelle in lingua naturale, o se vada scelta per tipo di query
        (sezioni~\ref{sec:nessuna-unica} e~\ref{sec:scelta});
  \item come si confrontano, alla fine, Koskidex corretto ed Elasticsearch di
        produzione, sullo stesso corpus e con le stesse query
        (capitolo~\ref{chapter:confronto}).
\end{enumerate}

\section{Contributi}\label{sec:contributi}

\begin{itemize}
  \item \textbf{Un impianto di valutazione riproducibile}: le stesse metriche,
        calcolate dallo stesso codice, per Koskidex in memoria e per i due motori
        interrogati attraverso Documentale; un archivio in cui ogni numero ha un
        file, un commit e un albero pulito; e le previsioni di ogni misura che ne ha
        committate prima della misura (capitolo~\ref{chapter:valutazione},
        appendici~\ref{chapter:appendice-riproducibilita}
        e~\ref{chapter:appendice-ipotesi}).
  \item \textbf{Quattro difetti di Koskidex, il primo trovato strada facendo,
        misurati e corretti}, ciascuno dietro un'impostazione che lascia il
        comportamento di partenza come default. Il primo, il recupero
        congiuntivo, non era fra i tre sospettati: è stato trovato misurando.
        Su SciFact il recupero lessicale passa da 0,0246 a 0,6757 di nDCG@10,
        il 99,5\% del riferimento BM25 di Lucene.
  \item \textbf{Quattro difetti della ricerca di Documentale in produzione}:
        il numero d'atto, le elisioni, la combinazione di numero e comune e le
        date in formati diversi. Su 300 ricerche \texttt{<numero> <comune>}
        l'atto giusto passa dall'1,3\% al 91,7\% dei casi al primo posto, con
        due impostazioni; una data scritta in un formato diverso da quello
        dell'atto passa dallo 0-8\% al 98-100\% di atti trovati entro i primi
        dieci.
  \item \textbf{Un innesto con gli stessi risultati}: Koskidex dentro
        Documentale, scelto da una variabile d'ambiente, che restituisce gli
        stessi insiemi di Elasticsearch su tutte le query del confronto
        (capitolo~\ref{chapter:sistemi}), con una mediana di 2,7~ms per ricerca
        contro 22, su due motori che non girano alla pari.
  \item \textbf{Un corpus pubblico di dominio}: 10.018 atti degli albi pretori
        comunali, riscaricabili da chiunque, con 300 query known-item generate e
        104 scritte da tre persone con un protocollo fissato prima della
        raccolta.
  \item \textbf{Un confronto su query scritte da persone}: sulle 104
        known-item umane Koskidex innestato con il profilo consigliato fa
        0,432 di MRR@10 contro 0,326 di Elasticsearch di produzione, e lascia a
        vuoto il 30\% delle ricerche contro il 46\%. Con le stesse correzioni anche
        Elasticsearch arriva a 0,449: sulle ricerche per contenuto il
        vantaggio sta nelle impostazioni, non nel motore. Dall'app, il profilo che tiene insieme le ricerche per
        contenuto e quelle per numero, con il recupero congiuntivo e i vettori,
        fa 0,607 e 0,956; lo stesso motore fuori dall'app, in recupero
        disgiuntivo e con i vettori, arriva a 0,613
        (capitolo~\ref{chapter:confronto}).
  \item \textbf{Una misura di come usare il vettore}: con il recupero
        congiuntivo di Documentale l'unione dei candidati vale +0,62 di nDCG@10
        su SciFact, e il peso giusto del vettore dipende dal tipo di query, tanto
        che sceglierlo con una regola di una riga vale +0,023 sulla migliore
        configurazione fissa.
\end{itemize}

\section{Il metodo: ipotesi prima della misura}\label{sec:metodo}

Ogni esperimento di questa tesi con una previsione segue lo stesso ordine; i
quattro senza previsioni numeriche sono dichiarati come tali
nell'appendice~\ref{chapter:appendice-ipotesi}. Prima una domanda e un
metodo; poi le previsioni, con una soglia numerica, committate prima di lanciare
la misura e spesso prima di scrivere il codice; poi la misura, da codice
committato, in un file che non si sovrascrive; poi l'esito, previsione per
previsione. Quello che si aggiunge dopo aver visto un numero lo si dichiara.

Serve a una cosa: che un risultato non sia stato scelto dopo averlo visto. E ha un
effetto collaterale che non avevo previsto: le previsioni sbagliate sono state le
più utili. Che le stopword valessero più dello stemmer, che le espansioni per
prefisso non fossero rare nemmeno sulle query in inglese, che la costante della
fusione non dipendesse dalla lunghezza delle query ma dal tipo: ognuna di queste
ha smentito una previsione, ed è scritta come tale.

\section{Struttura della tesi}\label{sec:struttura}

Il capitolo~\ref{chapter:fondamenti} richiama le nozioni di recupero e di
valutazione che servono. Il capitolo~\ref{chapter:sistemi} descrive Koskidex e
Documentale com'erano, e l'innesto che restituisce gli stessi risultati di
Elasticsearch. Il capitolo~\ref{chapter:valutazione} descrive l'impianto di
valutazione: regole, collezioni, corpus, query e strumenti. I
capitoli~\ref{chapter:lessicale} e~\ref{chapter:ibrido} misurano e correggono i
difetti di Koskidex, lessicali e vettoriali. Il capitolo~\ref{chapter:produzione}
misura i difetti della ricerca di Documentale in produzione. Il
capitolo~\ref{chapter:confronto} confronta Koskidex corretto ed Elasticsearch, e
il capitolo~\ref{chapter:conclusione} chiude con le risposte, i limiti e gli
sviluppi. Le appendici dicono come rifare ogni numero, riportano le istruzioni di
annotazione, il registro di tutte le previsioni e le impostazioni aggiunte a
Koskidex.


\pagestyle{fancy}
\chapter{Fondamenti e lavori correlati}\label{chapter:fondamenti}

Questo capitolo raccoglie le nozioni che i capitoli seguenti usano, nella forma in
cui le usano. Per una trattazione completa si rimanda a~\cite{manning2008ir}.

\section{Indice invertito e recupero booleano, congiuntivo e disgiuntivo}\label{sec:indice}

Un motore di ricerca testuale tiene un \textbf{indice invertito}: per ogni
termine, la lista dei documenti che lo contengono, con le posizioni o almeno il
numero delle occorrenze. Una query si spezza negli stessi termini, e i documenti
candidati si ottengono combinando le liste.

Le due combinazioni elementari sono l'intersezione e l'unione. Nel
\textbf{recupero congiuntivo} un documento è candidato solo se contiene tutti i
termini della query, come in un AND booleano; nel \textbf{recupero disgiuntivo}
basta che ne contenga uno, e l'ordine lo decide il punteggio. Il primo dà insiemi
piccoli e precisi quando le query sono di due o tre parole, ed è la scelta
naturale per la barra di ricerca di un negozio; il secondo è quello dei modelli
probabilistici, che assumono che ogni termine della query porti un'evidenza
parziale. Fra i due stanno le vie di mezzo: \texttt{minimum\_should\_match} di
Elasticsearch chiede un numero minimo di termini, eventualmente in funzione della
lunghezza della query; il fattore di \emph{coordinazione} della similarità
classica di Lucene moltiplica il punteggio per la quota di termini presenti, ed è
stato tolto in Lucene~7, dopo il passaggio a BM25.

Un dettaglio che conta nei documenti a più campi: in Elasticsearch un
\texttt{multi\_match} di tipo \texttt{best\_fields}~\cite{es-multimatch} valuta la
query campo per campo e tiene il campo migliore, e con \texttt{operator: and}
chiede che tutti i termini stiano \emph{nello stesso campo}. È una forma di
recupero congiuntivo più stretta di quella sul documento intero.

\section{Il punteggio: dall'euristico a BM25}\label{sec:bm25}

Un modello di punteggio assegna a ogni candidato un numero che ne stima la
pertinenza. Tre ingredienti ricorrono nei modelli classici:
\begin{itemize}
  \item la \textbf{frequenza del termine} nel documento, $\text{tf}(t, D)$: un
        documento che nomina un termine più volte ne parla di più, ma con un
        guadagno che si satura;
  \item la \textbf{frequenza documentale inversa}~\cite{sparckjones1972idf}: un
        termine che compare in pochi documenti distingue di più di uno che
        compare in tutti;
  \item la \textbf{normalizzazione sulla lunghezza}: un documento lungo contiene
        più termini per il solo fatto di essere lungo.
\end{itemize}

BM25~\cite{robertson2009bm25} li combina così. Per una query $Q$ e un documento
$D$,
\[
  \text{BM25}(D, Q) = \sum_{t \in Q} \text{idf}(t) \cdot
  \frac{\text{tf}(t, D) \cdot (k_1 + 1)}
       {\text{tf}(t, D) + k_1 \cdot \left(1 - b + b \cdot \frac{|D|}{\text{avgdl}}\right)},
\]
\[
  \text{idf}(t) = \ln\left(1 + \frac{N - \text{df}(t) + 0{,}5}{\text{df}(t) + 0{,}5}\right),
\]
dove $N$ è il numero dei documenti, $\text{df}(t)$ quelli che contengono $t$,
$|D|$ la lunghezza del documento in termini e avgdl la lunghezza media. $k_1$
regola la saturazione della frequenza, $b$ quanto pesa la lunghezza; i valori
predefiniti di Lucene sono $k_1 = 1{,}2$ e $b = 0{,}75$, e l'IDF qui sopra è la
variante di Lucene, sempre positiva. BM25 è la similarità predefinita di Lucene dalla
versione~6, e di Elasticsearch dalla~5.

Un punteggio \textbf{euristico}, come quello di Koskidex di partenza
(sezione~\ref{sec:koskidex}), non ha nessuno dei tre ingredienti: somma un
valore fisso per ogni termine trovato, con un premio per le corrispondenze esatte
e il peso del campo. Con il recupero congiuntivo la differenza si sente poco,
perché i candidati contengono tutti gli stessi termini; con il recupero
disgiuntivo diventa il difetto principale, perché un documento con molti termini
comuni batte quello con l'unico termine raro che decide la query.

\section{Prefisso, refusi e frequenza mescolata}\label{sec:espansioni-teoria}

Chi cerca non sempre scrive la parola che sta nel documento. Due meccanismi
allargano la corrispondenza, \emph{espandendo} un termine della query in più
termini dell'indice:
\begin{itemize}
  \item la \textbf{ricerca per prefisso}: \texttt{manut} trova
        \texttt{manutenzione}; utile mentre si scrive, o per parole lunghe
        scritte a metà;
  \item la \textbf{tolleranza ai refusi}: un termine trova quelli a distanza di
        modifica limitata, di solito la distanza di Damerau~\cite{damerau1964},
        che conta inserimenti, cancellazioni, sostituzioni e scambi di due
        caratteri vicini. La distanza ammessa cresce con la lunghezza del
        termine: in Elasticsearch, con \texttt{fuzziness: AUTO}, nessun refuso
        fino a due caratteri, uno da tre, due da sei; in
        Meilisearch~\cite{meili-typo} uno da cinque e due da nove.
\end{itemize}

Un'espansione pone un problema di punteggio: con quale frequenza documentale si
pesa un termine trovato per espansione? Se con la sua, un termine raro trovato per
prefisso o per refuso ha un IDF altissimo e può battere il termine esatto, più
comune. Lucene risolve il problema in due modi diversi. Le ricerche per prefisso
le calcola, per default, a punteggio costante, senza IDF. Le ricerche con refusi
le riscrive con \texttt{TopTermsBlendedFreqScoringRewrite}, che \emph{mescola} le
frequenze documentali dei termini espansi, così che nessuna espansione pesi più
del più comune dei suoi fratelli. Il capitolo~\ref{chapter:lessicale} ritrova
esattamente questo difetto in Koskidex, e lo corregge allo stesso modo.

\section{Analisi lessicale}\label{sec:analisi}

Prima di entrare nell'indice un testo passa da un \textbf{analizzatore}: una
catena di stadi che lo trasforma in termini. Gli stadi usati in questa tesi sono
quattro.
\begin{itemize}
  \item La \textbf{tokenizzazione} spezza il testo in parole. La scelta più
        semplice spezza su ogni carattere che non è una lettera o una cifra; il
        tokenizer standard di Lucene segue le regole di segmentazione di
        Unicode~\cite{uax29}, che tengono insieme alcune parole con punteggiatura
        interna, fra cui \emph{dell'infanzia}.
  \item Le \textbf{stopword} sono le parole funzionali (articoli, preposizioni,
        congiunzioni) che si tolgono perché non distinguono i documenti.
  \item Lo \textbf{stemming} riduce le varianti morfologiche di una parola a una
        radice comune: in inglese l'algoritmo di Porter~\cite{porter1980}, in
        italiano lo stemmer leggero di Savoy~\cite{savoy2002light}, che toglie
        soprattutto le desinenze di genere e numero.
  \item Il filtro delle \textbf{elisioni} toglie gli articoli elisi davanti a una
        parola (\emph{l'}, \emph{dell'}, \emph{un'}): in italiano, con un
        tokenizer che non spezza sull'apostrofo, senza questo filtro
        \emph{dell'infanzia} resta un termine diverso da \emph{infanzia}.
\end{itemize}
L'analizzatore \texttt{italian} di Elasticsearch~\cite{es-italian} usa il
tokenizer standard, il filtro delle elisioni, le stopword italiane di Snowball e
lo stemmer leggero. L'analizzatore predefinito, che Documentale usa, ha il solo
tokenizer standard e il minuscolo.

\section{Recupero denso e ibrido, fusione dei punteggi}\label{sec:denso}

Il recupero lessicale trova i documenti che condividono termini con la query. Il
\textbf{recupero denso} rappresenta query e documenti come vettori in uno spazio
di qualche centinaio o migliaio di dimensioni, calcolati da un modello neurale
(\emph{embedding}), e trova i documenti più vicini per prodotto scalare o
similarità del coseno~\cite{karpukhin2020dpr}. Trova documenti pertinenti che non condividono
parole con la query; non sa niente degli identificativi, perché un numero di
protocollo e uno che differisce di una cifra sono, per il modello, quasi lo stesso
testo. Su archivi di milioni di documenti i vicini si cercano con un indice
approssimato; su poche decine di migliaia basta una scansione esatta.

Il \textbf{recupero ibrido} usa i due insieme, e pone due domande distinte. La
prima è chi entra fra i candidati: solo i documenti trovati dal lessicale, che il
vettore riordina (\emph{re-ranking}), o l'unione dei due insiemi. La seconda è
come si fondono due punteggi su scale diverse. Tre famiglie di fusione:
\begin{itemize}
  \item la \textbf{somma pesata} dei punteggi grezzi, che dipende dalle scale e
        quindi dalla collezione e dalla query;
  \item la \textbf{combinazione convessa} di punteggi normalizzati, per esempio
        min-max su ciascuna lista, $\alpha \cdot s_\text{lex} + (1 - \alpha)
        \cdot s_\text{vec}$, che rende le scale confrontabili ma lascia un peso da
        scegliere;
  \item la \textbf{Reciprocal Rank Fusion}~\cite{cormack2009rrf}, che ignora i
        punteggi e somma $1/(k + r)$ sulle posizioni $r$ del documento nelle
        liste, con $k = 60$: senza parametri da calibrare, ma cieca alla distanza
        fra un primo e un secondo classificato.
\end{itemize}
Bruch, Gai e Ingber~\cite{bruch2023fusion} confrontano le ultime due e trovano
che una combinazione convessa ben calibrata batte RRF, e che RRF è sensibile ai
suoi parametri più di quanto la sua fama di fusione senza parametri suggerisca. Il capitolo~\ref{chapter:ibrido}
aggiunge un'osservazione che quel confronto non fa: il peso giusto sembra dipendere dal
tipo di query, identificativa o in lingua naturale (che nei dati coincide con la collezione).

\section{Valutazione: Cranfield, pooling, metriche, known-item, accordo}\label{sec:valutazione-teoria}

\paragraph{Il paradigma di Cranfield.} Una collezione di valutazione ha tre
parti: i documenti, le query e i giudizi di pertinenza che dicono quali documenti
rispondono a quali query~\cite{cleverdon1967}. Con i giudizi fissati, due sistemi
si confrontano sulle stesse query con la stessa metrica, e il confronto si rifà.

\paragraph{Il pooling.} Giudicare ogni documento per ogni query non è possibile su
una collezione vera. Il \emph{pool} di una query è l'unione dei primi $k$
risultati di più sistemi, e si giudica solo quello; i documenti fuori dal pool si
considerano non pertinenti. Zobel~\cite{zobel1998pooling} ha mostrato che i
confronti fra sistemi restano affidabili, ma che il richiamo misurato è un limite
superiore: i pertinenti che nessun sistema ha trovato non si contano.

\paragraph{Le metriche.} Con $\text{rel}_i$ il grado del documento in posizione
$i$, il guadagno cumulato scontato è
$\text{DCG@}k = \sum_{i=1}^{k} (2^{\text{rel}_i} - 1)/\log_2(i + 1)$, e nDCG@$k$
lo divide per il DCG dell'ordinamento ideale, così che stia fra 0 e
1~\cite{jarvelin2002ndcg}; è la metrica principale di BEIR~\cite{thakur2021beir}.
L'MRR@$k$ è il reciproco della posizione del primo documento pertinente, zero se
oltre $k$: con un solo documento giusto per query è la misura più diretta.
Recall@$k$ è la quota dei documenti pertinenti fra i primi $k$.

\paragraph{Known-item.} Una query \emph{known-item} cerca un documento preciso,
che chi cerca sa esistere: il giudizio arriva gratis, perché il documento di
partenza è l'unica risposta giusta. Le query si possono raccogliere da persone o
generare a partire dal documento; Azzopardi, de Rijke e
Balog~\cite{azzopardi2007known} studiano come generare query simulate e le confrontano con
query reali: le simulate non sempre riproducono i risultati ottenuti con quelle
reali. È la ragione per cui questa
tesi tiene separate le known-item automatiche, che misurano una sola forma di
query, e quelle scritte da persone.

\paragraph{Accordo fra annotatori.} Due annotatori non danno gli stessi giudizi.
Il kappa di Cohen~\cite{cohen1960kappa} misura il loro accordo al netto di quello
atteso per caso, e la versione pesata~\cite{cohen1968weighted} conta meno i
disaccordi fra gradi vicini. Un kappa basso non invalida una collezione, ma dice
quanto le differenze fra i sistemi sono più grandi della variabilità dei giudizi.

\section{Lavori correlati}\label{sec:correlati}

\paragraph{BEIR e le regressioni di Pyserini.} BEIR~\cite{thakur2021beir} raccoglie
diciotto collezioni eterogenee per valutare il recupero senza addestramento
specifico, e mostra che BM25 resta un riferimento difficile da battere fuori
dominio. Le regressioni di Pyserini~\cite{pyserini2cr} ne danno i numeri di BM25
di Lucene, riproducibili con un comando: sono il riferimento usato per SciFact e
NFCorpus.

\paragraph{Lucene, Elasticsearch, Meilisearch.} I motori di produzione con cui
Koskidex si confronta hanno ciascuno le loro scelte documentate: Lucene ed
Elasticsearch per BM25, le espansioni e gli analizzatori; Meilisearch, più
vicino a Koskidex per destinazione, per la tolleranza ai refusi, compresa
l'impostazione \texttt{disableOnNumbers} che la spegne sui numeri. Quando una correzione
di questa tesi ha un precedente in uno di questi motori, lo segue, e lo dice: il
contributo non è inventare una soluzione, ma trovare e misurare il difetto.

\paragraph{Predizione della difficoltà e scelta per query.} Stimare prima di
cercare quanto una query riuscirà è un campo di ricerca, la predizione della
difficoltà delle query~\cite{cronentownsend2002,carmel2010qpp}, e una delle sue
applicazioni è scegliere per ogni query una strategia di recupero diversa. Il
capitolo~\ref{chapter:ibrido} ne fa una versione piccola, con otto
caratteristiche della query e una regressione logistica, e trova che su questi
dati basta una regola di una riga.

\chapter{I sistemi e l'innesto}\label{chapter:sistemi}

Questa tesi misura i difetti di recupero di un motore di ricerca piccolo,
Koskidex, dentro un sistema vero, Documentale. Prima di misurare qualsiasi
difetto serve una condizione: che il motore piccolo, messo al posto di
Elasticsearch, trovi gli stessi documenti. Senza quella condizione ogni
miglioramento misurato su Koskidex si potrebbe liquidare con una frase sola:
\emph{stai migliorando il tuo motore, non il nostro sistema}. Questo capitolo
descrive i due sistemi com'erano all'inizio del lavoro, l'innesto dell'uno
nell'altro e le differenze di matching che è stato necessario trovare e
correggere per arrivare agli stessi risultati.

\section{Koskidex: com'era il 23 settembre 2026}\label{sec:koskidex}

Koskidex è un motore di ricerca full-text scritto in Go, nato come progetto
personale per la barra di ricerca di un piccolo negozio online. Si distribuisce
come un binario solo, senza servizi esterni, e l'unica dipendenza del codice è
\texttt{golang.org/x/text}, per la normalizzazione Unicode. Quando è stato scelto
come oggetto di questa tesi contava circa 2.850 righe di codice e 2.060 di test.

\paragraph{Struttura.} Il cuore è un indice invertito in memoria: per ogni
termine la lista delle occorrenze (\emph{posting}), ciascuna con documento,
campo e posizione. Accanto all'indice ci sono i documenti originali, l'elenco
dei termini di ogni documento (per cancellarlo senza scorrere l'indice intero) e
una mappa da ogni coppia di caratteri consecutivi (bigramma) ai termini che la
contengono, usata per cercare i refusi. La persistenza è un registro delle operazioni (\emph{write-ahead log}) in
JSON, scritto prima di toccare l'indice, più un'istantanea periodica di tutti i
documenti; all'avvio il motore ricarica l'istantanea e riapplica il registro.
Tutto si usa da un'API HTTP: indici, documenti, impostazioni, ricerca con
filtri e faccette.

\paragraph{La ricerca.} Una query si spezza in termini su ogni carattere che non
è una lettera o una cifra, in minuscolo e senza accenti. Ogni termine della
query è \textbf{obbligatorio}: un documento entra fra i risultati solo se li
contiene tutti. Per ogni termine valgono tre forme di corrispondenza:
\begin{itemize}
  \item la forma esatta;
  \item la ricerca per prefisso, per cui \texttt{manut} trova
        \texttt{manutenzione};
  \item i refusi, con la distanza di Damerau~\cite{damerau1964}: uno da quattro
        caratteri in su, due da otto.
\end{itemize}
Il punteggio di un documento è la somma, sui termini, di
\[
  (10 - \text{refusi} + 2 \cdot \text{esatto}) \cdot \text{peso del campo},
\]
dove \emph{esatto} vale 1 se la corrispondenza è esatta. Non c'è frequenza
documentale inversa: un termine che compare in un documento su diecimila pesa
quanto uno che compare in tutti. Un campo \texttt{TF} esisteva nella struttura
delle occorrenze, con il commento \emph{calculated later}, e non era letto da
nessuna parte.

\paragraph{I vettori.} Un documento può avere un vettore, e una ricerca può
portare il vettore della query. In quel caso al punteggio di ogni documento
\emph{già trovato} dalla parte lessicale si aggiunge la similarità del coseno
moltiplicata per 20. È un riordino dei risultati lessicali, non una ricerca
ibrida: con una query che ha almeno un termine, un documento che non ne
contiene le parole non entra mai.

\paragraph{Il punto di partenza, congelato.} Il comportamento di quel giorno è
fissato in un test, \texttt{TestBaselineRankingIsFrozen}, che confronta gli
ordinamenti di un corpus di prova con quelli registrati e che deve passare,
senza modifiche al test, a ogni commit successivo (con un'unica eccezione, tre
minuti dopo, descritta nella sezione~\ref{sec:principi}). Ogni correzione descritta
nei capitoli seguenti sta dietro un'impostazione, con il comportamento di
partenza come valore predefinito. Misurato sulle collezioni pubbliche del
capitolo~\ref{chapter:valutazione}, il punto di partenza fa 0,0246 di nDCG@10 su
SciFact e 0,1659 su NFCorpus.

Letto il codice, tre scelte sembravano i candidati a difetti di recupero:
l'assenza di IDF, il vettore che riordina senza recuperare e la costante 20 che
somma due scale diverse. Con il recupero congiuntivo sono i quattro difetti dei
capitoli~\ref{chapter:lessicale} e~\ref{chapter:ibrido}. Il recupero
congiuntivo, come si vedrà, è stato trovato misurando e non leggendo: la tesi
stava per essere scritta su tre capitoli di ottimizzazione del punteggio di un motore che su SciFact non
restituiva niente in 290 query su 300.

\section{Documentale: come cerca in produzione}\label{sec:documentale}

Documentale è un sistema di gestione documentale sviluppato da Dieffetech: un
backend Laravel, un'interfaccia Next.js, MySQL per i dati ed Elasticsearch per la
ricerca. Chi lo usa archivia atti e documenti in cartelle, con un nome, dei tag,
un riassunto, dei soggetti (gli enti o le persone coinvolte), delle note e dei
metadati liberi.

\paragraph{Due indici.} La ricerca passa da un servizio,
\texttt{ElasticsearchService}, che tiene due indici: uno per i documenti e uno
per le cartelle. L'indice dei documenti ha sei campi di testo, \texttt{name},
\texttt{tags}, \texttt{summary}, \texttt{subjects}, \texttt{notes} e
\texttt{additional\_data}, con l'analizzatore predefinito di Elasticsearch, cioè
il tokenizer standard senza filtri per l'italiano. Il testo integrale dei
documenti non è indicizzato: si cerca sulla scheda.

\paragraph{La query.} La ricerca dei documenti è un \texttt{multi\allowbreak\_match}
di tipo \texttt{best\_fields}~\cite{es-multimatch} con quattro parametri, che
il codice commenta uno per uno come il tipo, e i pesi dei campi:
\begin{itemize}
  \item \texttt{operator: and}, \emph{tutte le parole devono comparire};
  \item \texttt{fuzziness: AUTO}, un refuso da tre caratteri, due da sei;
  \item \texttt{prefix\_length: 1}, la prima lettera non ammette refusi;
  \item \texttt{tie\_breaker: 0.3}: al punteggio del campo migliore si somma
        0,3 volte quello degli altri campi che combaciano, e quindi cambia
        l'ordine dei risultati, non il loro insieme;
  \item i pesi dei campi: \texttt{name} 5, \texttt{tags} 4, \texttt{summary} e
        \texttt{subjects} 3, \texttt{notes} 2, \texttt{additional\_data} 1.
\end{itemize}
Con \texttt{best\_fields} e \texttt{operator: and} tutte le parole della query
devono stare \emph{nello stesso campo}: è un dettaglio che il
capitolo~\ref{chapter:produzione} ritroverà come difetto. La ricerca delle
cartelle è diversa: un \texttt{match} con refusi sul percorso, in disgiunzione,
unito a un \emph{wildcard} \texttt{*termine*}.

\paragraph{Il recupero è congiuntivo anche qui.} Koskidex ed Elasticsearch, in
due basi di codice indipendenti, arrivano alla stessa scelta per la stessa
ragione: in una barra di ricerca, chi scrive due parole si aspetta documenti che
le contengano tutte e due. È la scelta che il capitolo~\ref{chapter:lessicale}
misura come difetto 0 quando le query si allungano.

\paragraph{Il corpus.} Documentale non ha più un archivio di clienti: Dieffetech
lo ha abbandonato nel settembre 2026, e con lui il progetto Elasticsearch in
cloud. Il sistema però gira in locale, con i suoi comandi di importazione e di
indicizzazione, e il contenuto è stato sostituito con un corpus pubblico della
stessa famiglia: 10.018 atti degli albi pretori comunali, 563 del Comune di
Crispiano con il testo integrale e 9.455 di 171 enti del Friuli Venezia Giulia
con i soli metadati. Il corpus è descritto nel
capitolo~\ref{chapter:valutazione}. Qui conta che gli atti entrano in
Documentale dal suo codice di importazione e in Elasticsearch dal suo comando di
indicizzazione, come in produzione: la ricerca misurata è quella dell'app, non
una sua ricostruzione.

\section{L'innesto: un contratto, due motori, gli stessi risultati}\label{sec:innesto}

\paragraph{Il contratto.} L'app non parla con Elasticsearch direttamente ma con
un'interfaccia, \texttt{SearchBackend}, estratta da \texttt{ElasticsearchService}
con una regola sola: le firme sono quelle del servizio in produzione, e se una
firma non torna si corregge l'interfaccia, non il servizio. I metodi sono quelli
che l'app usa davvero: creare i due indici, indicizzare un elemento o un blocco
nel formato \emph{bulk} di Elasticsearch, cancellare, aggiornare, controllare
l'esistenza e cercare (\texttt{fuzzySearch}, che restituisce gli id in ordine di
rilevanza). Accanto a \texttt{ElasticsearchService} c'è \texttt{KoskidexService},
che implementa lo stesso contratto su un client HTTP di Koskidex, e quale dei due
si usa lo decide una variabile d'ambiente, \texttt{SEARCH\_BACKEND}. Nient'altro
dell'app cambia (figura~\ref{fig:innesto}).

\begin{figure}[ht]
  \centering
  \begin{tikzpicture}[node distance=0.9cm and 1.2cm,
      box/.style={draw, rounded corners=2pt, align=center, font=\footnotesize, inner sep=5pt, minimum height=8mm},
      nota/.style={font=\scriptsize\itshape, align=left},
      freccia/.style={-{Stealth[length=5pt]}, line width=.7pt}]
    \node[box] (app) {App Documentale (Laravel)};
    \node[box, below=of app] (sb) {\texttt{SearchBackend}\\ (interfaccia)};
    \node[box, below left=1.1cm and 0.7cm of sb] (es) {\texttt{ElasticsearchService}\\ (produzione, non toccato)};
    \node[box, below right=1.1cm and 0.7cm of sb] (ks) {\texttt{KoskidexService}\\ (aggiunto)};
    \node[box, below=of es] (esm) {Elasticsearch};
    \node[box, below=of ks] (kc) {client HTTP};
    \node[box, below=of kc] (k) {Koskidex\\ indice in memoria};
    \draw[freccia] (app) -- (sb);
    \draw[freccia] (sb) -- (es);
    \draw[freccia] (sb) -- (ks);
    \draw[freccia] (es) -- (esm);
    \draw[freccia] (ks) -- (kc);
    \draw[freccia] (kc) -- (k);
    \node[nota, right=0.4cm of sb] {motore scelto da\\ \texttt{SEARCH\_BACKEND}};
  \end{tikzpicture}
  \caption{L'innesto. L'app parla con un'interfaccia estratta da
  \texttt{ElasticsearchService}, con le stesse firme; \texttt{KoskidexService}
  la implementa su un client HTTP di Koskidex. Cambiando una variabile
  d'ambiente si passa da un motore all'altro e nient'altro dell'app cambia.}
  \label{fig:innesto}
\end{figure}


\paragraph{Prima dell'innesto, un motore completo.} La regola era di non
innestare un Koskidex che non sapesse fare tutto quello che l'app chiede.
Provandolo contro il contratto sono venuti fuori sette buchi, chiusi uno per uno
in Koskidex con test scritti prima e visti fallire: gli id interi, che venivano
scartati in silenzio; i campi lista (i tag, i soggetti), che venivano salvati ma
non indicizzati; un limite di risultati troppo basso e risposte senza punteggio;
la cancellazione di molti documenti in una richiesta; la ricerca per
sottostringa che serve alle cartelle; e le impostazioni per riprodurre il
matching di Elasticsearch, descritte nella sezione seguente.

\paragraph{Le impostazioni di compatibilità.} \texttt{KoskidexService} crea
l'indice dei documenti con le impostazioni che riproducono la query di
produzione: tutti i termini obbligatori, tutti nello stesso campo, le soglie dei
refusi di \texttt{fuzziness: AUTO}, niente ricerca per prefisso, la prima lettera
esatta, il tokenizer standard e gli stessi pesi dei campi. Il punteggio resta
quello di Koskidex: la compatibilità riguarda \emph{quali} documenti si trovano,
non in che ordine.

\paragraph{Il risultato.} Dal vivo, passando dall'app sui 10.018 atti e cambiando
solo \texttt{SEARCH\_BACKEND}, Koskidex restituisce \textbf{gli stessi insiemi di
Elasticsearch su 24 query su 24} del confronto, e su 300 su 300 delle query
\emph{known-item} del capitolo~\ref{chapter:produzione}, una famiglia di query
che il confronto non conteneva. Gli ordinamenti invece coincidono nei primi dieci
risultati solo su 7 query, quelle con al più due risultati: il punteggio
euristico di Koskidex non è BM25. È esattamente la condizione che serve. Da qui
in poi, ogni differenza misurata fra Koskidex e il suo punto di partenza viene
dalla correzione che si sta misurando, non da un motore che trova altro.

La parità ha un rovescio, e va detto subito: \textbf{riproducendo Elasticsearch,
Koskidex ne eredita i difetti}. Sulle 300 query known-item trova gli stessi
insiemi di Elasticsearch, query per query, e quindi manca gli stessi atti. La parità non è il traguardo ma il punto da
cui misurare le correzioni del capitolo~\ref{chapter:produzione}.

\section{Le differenze di matching trovate per arrivarci}\label{sec:differenze}

La parità non era scontata. Su un primo campione di quattordici documenti con la
forma dei dati di Documentale i due motori trovavano gli stessi insiemi; sul
corpus vero di 10.018 atti coincidevano su 8 query su 24. Il modello di recupero
era lo stesso, congiuntivo, ma il matching differiva. Per trovare in cosa, a
Koskidex è stato dato esattamente il contenuto che Elasticsearch aveva
indicizzato, copiato dall'indice, e si è aggiunta una correzione alla volta,
contando a ogni passo le query con lo stesso insieme e i documenti trovati da un
motore solo (tabella~\ref{tab:divergenza}).

\begin{table}[ht]
  \centering
  \small
  \begin{tabular}{lrrr}
    \toprule
    Passo & identiche & solo ES & solo Koskidex \\
    \midrule
    0. Koskidex com'era & 8 / 24 & 124 & 447 \\
    1. parole nello stesso campo & 10 / 24 & 124 & 121 \\
    2. soglie dei refusi di Elasticsearch & 12 / 24 & 0 & 154 \\
    3. niente ricerca per prefisso & 13 / 24 & 0 & 116 \\
    4. prima lettera esatta & 15 / 24 & 0 & 84 \\
    5. tokenizer standard & \textbf{24 / 24} & \textbf{0} & \textbf{0} \\
    \bottomrule
  \end{tabular}
  \caption{Le differenze fra i due motori sulle 24 query del confronto,
  aggiungendo a Koskidex una correzione alla volta: query con gli stessi risultati,
  e risultati trovati da uno solo dei due motori (ES: Elasticsearch). I passi 2-5 sono misurati
  dopo la correzione del difetto dei bigrammi descritta nel testo.}
  \label{tab:divergenza}
\end{table}

\paragraph{Le cinque cause.}
\begin{enumerate}
  \item \textbf{La congiunzione per campo.} Con \texttt{best\_fields} e
        \texttt{operator: and} tutte le parole devono stare nello stesso campo;
        Koskidex le accettava sparse fra i campi. È la causa più grossa: 326
        dei 447 documenti in più.
  \item \textbf{Le soglie dei refusi.} Elasticsearch ammette un refuso da tre
        caratteri e due da sei, Koskidex da quattro e da otto.
  \item \textbf{La ricerca per prefisso}, che Koskidex fa e Elasticsearch, con
        quella query, no.
  \item \textbf{La prima lettera esatta}, per \texttt{prefix\_length: 1}.
  \item \textbf{Il tokenizer.} Il tokenizer standard di Elasticsearch segue le
        regole di segmentazione di Unicode~\cite{uax29}, che tengono dentro la
        parola un apostrofo fra due lettere e uniscono date e decimali:
        \emph{dell'illuminazione} è un termine solo, e \texttt{14.01.2026}
        anche. Koskidex spezzava su ogni carattere che non è lettera o cifra.
\end{enumerate}
Ciascuna è diventata un'impostazione di Koskidex, spenta per default.

\paragraph{Due ipotesi sbagliate, e un difetto di Koskidex.} Dopo il quarto passo
restavano 9 query diverse, con 42 documenti trovati solo da Elasticsearch, tutti
sulla query \texttt{ordinanza 187}. Le due spiegazioni scritte allora
erano plausibili e sbagliate tutte e due. La prima era la tokenizzazione dei
numeri. La seconda era il limite di espansioni dei refusi di Elasticsearch
(\texttt{max\_expansions}), smentita rilanciando la query di produzione con un
limite di diecimila: stessi insiemi. La causa vera era un difetto di Koskidex.
I candidati per i refusi venivano solo dai termini con almeno un bigramma in
comune con la parola cercata, e \texttt{17} non ne ha nessuno con \texttt{187},
pur distando un refuso. Una modifica rompe fino a due bigrammi, una
trasposizione tre, e sotto una certa lunghezza non resta nessuna garanzia: con
le soglie predefinite capitava anche ad \texttt{atre} contro \texttt{arte}. Il
motore prometteva la distanza di Damerau e non la manteneva. La correzione non è
un'impostazione: per le parole troppo corte perché i bigrammi garantiscano un
candidato, ora si scorre il vocabolario intero, con un filtro sulla lunghezza
prima della distanza.

Il quinto passo è arrivato chiedendo a Elasticsearch, con \texttt{\_explain},
perché non trovava un documento preciso: il nome diceva \emph{dell'illuminazione
pubblica}. È la lezione di metodo di questa sezione, e torna nel resto della
tesi: le ipotesi scritte senza verificarle erano ragionevoli, e l'unica che ha
risolto è venuta da una domanda precisa a un documento preciso.

\paragraph{Un difetto di Documentale, visto per la prima volta.} Il tokenizer
standard è la scelta giusta per le date (\texttt{18} non deve combaciare dentro
\texttt{18.01.2026}) e quella sbagliata per le elisioni: chi cerca
\emph{illuminazione} non trova gli atti che dicono \emph{dell'illuminazione}. Il
capitolo~\ref{chapter:produzione} conta quanti sono.

\section{Prestazioni dell'innesto}\label{sec:prestazioni-innesto}

Con l'innesto fatto, lo stesso confronto dà i costi dei due motori passando
dall'app: indicizzazione dei 10.018 atti da indice vuoto con il comando di
Documentale, e le 24 query tre volte per motore (tabella~\ref{tab:innesto}).

\begin{table}[ht]
  \centering
  \begin{tabular}{lrr}
    \toprule
    & Koskidex & Elasticsearch \\
    \midrule
    insiemi identici sulle 24 query & \multicolumn{2}{c}{24 / 24} \\
    primi dieci identici & \multicolumn{2}{c}{7 / 24} \\
    ricerca, mediana per query & 2,7 ms & 22,0 ms \\
    indicizzazione di 10.018 atti & 1,9 s & 1,7 s \\
    \bottomrule
  \end{tabular}
  \caption{L'innesto misurato dall'app, sugli stessi atti e sulle stesse query.
  Koskidex compilato dal commit \texttt{f583d03}, Elasticsearch 9.1.0 locale.}
  \label{tab:innesto}
\end{table}

I tempi vanno letti come ordini di grandezza, per due ragioni: Elasticsearch
gira nella macchina virtuale di Docker e Koskidex nativo, e tutti e due i tempi
comprendono il viaggio HTTP da PHP. Quello che dicono è che, su un archivio di
diecimila atti, il motore piccolo non costa di più. Il confronto alla pari,
con i due motori nella stessa macchina virtuale, sotto carico e su archivi più
grandi, è nella sezione~\ref{sec:latenza}.

\paragraph{La prima misura ha trovato un altro difetto.} La prima
indicizzazione con Koskidex costava 31,8 secondi: 3,2 millisecondi per
documento. Ogni documento sincronizzava il registro delle operazioni
sul disco, e su macOS una sincronizzazione è un \texttt{F\_FULLFSYNC}, che
attende il supporto fisico. Ora una richiesta scrive tutte le sue righe del
registro con una sincronizzazione sola, e se la scrittura fallisce non entra
niente nell'indice: da 31,8 a 1,9 secondi, con lo stesso formato del registro.

È la seconda volta in questo capitolo, dopo i bigrammi, che la misura trova un
difetto che nessuno stava cercando. Non sarà l'ultima, e non è un caso: un
impianto che misura tutto dall'app, con il codice di produzione, passa dagli
stessi percorsi che percorrerebbe un utente, compresi quelli che nessun test
aveva previsto.

\chapter{L'impianto di valutazione}\label{chapter:valutazione}

Ogni numero di questa tesi viene da un file di un archivio, prodotto da codice
committato, e ogni misura che conta ha le sue previsioni scritte prima. Questo
capitolo descrive le regole, le collezioni, le query e gli strumenti; le regole
vengono prima, perché sono loro a dire quanto ci si può fidare del resto.

\section{Principi: ipotesi prima, interruttori, baseline congelato, archivio con provenienza}\label{sec:principi}

\paragraph{Le ipotesi prima della misura.} Ogni modifica che cambia l'ordine dei
risultati ha una voce nel diario di Koskidex, e ogni esperimento un README, con
una sezione \emph{Prima di misurare} committata prima di lanciare la
valutazione, spesso prima di scrivere il codice. Le previsioni sono numeriche,
con una soglia, e dicono cosa vorrebbe dire sbagliarle. A misura fatta si
aggiunge l'esito, previsione per previsione, senza toccare quello che era scritto
prima. Quello che si decide dopo aver visto un numero, un'analisi in più o una
strategia nuova, sta in una sezione che lo dichiara. È la risposta alla domanda
che in una valutazione arriva sempre: le metriche e le soglie sono state scelte
dopo aver visto i risultati? L'appendice~\ref{chapter:appendice-ipotesi} le raccoglie
tutte, e molte sono state smentite.

\paragraph{Gli interruttori.} Ogni correzione sta dietro un'impostazione, con il
comportamento di partenza come valore predefinito. Le impostazioni si salvano con
l'indice, e un indice salvato prima di una correzione non cambia comportamento da
solo. Così ogni misura confronta due configurazioni dello stesso codice, e i
default si decidono alla fine, tutti insieme.

\paragraph{Il baseline congelato.} Il comportamento di Koskidex del 23 settembre
2026 è fissato in un test, \texttt{TestBaselineRankingIsFrozen}, che confronta gli
ordinamenti di un corpus di prova con quelli registrati. Deve passare a ogni
commit, con le impostazioni predefinite e senza che il test venga toccato: è la
prova che nessuna correzione ha cambiato il punto di partenza di nascosto.

\paragraph{L'archivio.} Le esecuzioni vanno in un archivio, nel repository della
tesi, con regole scritte prima di riempirlo:
\begin{itemize}
  \item un file per esecuzione, con l'ora UTC nel nome, mai sovrascritto né
        modificato a mano;
  \item ogni file registra da dove viene: il commit del codice, se l'albero aveva
        modifiche non committate, il corpus con la sua impronta, ogni
        impostazione. Per la tesi valgono solo i file da alberi puliti;
  \item i tempi stanno separati dalle metriche, perché cambiano a ogni esecuzione
        e le metriche no;
  \item solo identificativi, mai titoli: gli atti contengono nomi di persone, e
        il corpus resta fuori dal repository;
  \item solo esecuzioni vere: le prove fatte mentre si scrive codice si lanciano
        con l'archivio spento, e gli strumenti, se l'archivio non è impostato,
        lo dicono invece di tacere.
\end{itemize}
Le eccezioni stanno in un registro: dodici file del primo giorno prodotti da
codice non ancora committato e rifatti, un file rinominato, tre ridotti da 12 MB a
1 MB conservando i primi mille risultati, e così via. Nessuno è stato tolto in
silenzio.

\paragraph{Le guardie.} Una previsione può essere una guardia: una quantità che la
modifica non deve muovere. La prima scritta, «BM25 non cambia il Recall@100», è
scattata perché era posta male (sezione~\ref{sec:difetto1}); da lì il valutatore
registra per ogni query il numero dei candidati prima del taglio, che una modifica
al punteggio non può cambiare e una al recupero sì.

\section{Le collezioni pubbliche: SciFact e NFCorpus}\label{sec:collezioni}

Due collezioni di BEIR~\cite{thakur2021beir}, scaricate con un controllo
dell'impronta degli archivi (tabella~\ref{tab:collezioni}). Servono a una cosa
sola: dire se il recupero lessicale di Koskidex ha un difetto, contro un
riferimento noto, BM25 di Lucene nella configurazione \emph{flat} di BEIR
(0,6789 e 0,3218 di nDCG@10~\cite{pyserini2cr}). Non rispondono alla domanda
della tesi, che vive sul corpus di dominio.

\begin{table}[ht]
  \centering
  \begin{tabular}{lrr}
    \toprule
    & SciFact & NFCorpus \\
    \midrule
    documenti & 5.183 & 3.633 \\
    query di test & 300 & 323 \\
    giudizi di test & 339 (1,1 per query) & 12.334 (38,2 per query) \\
    rilevanza & binaria & graduata (1 e 2) \\
    termini per query, mediana & 12 & 2 \\
    split di addestramento & train, 809 query & train 2.590, dev 324 \\
    \bottomrule
  \end{tabular}
  \caption{Le collezioni pubbliche.}
  \label{tab:collezioni}
\end{table}

\textbf{SciFact} è la prova di fumo: il riferimento è alto, e con un solo
documento pertinente per query una query si verifica a mano in un minuto.
\textbf{NFCorpus} è la prova vera: con 38 giudizi graduati per query il nDCG è
stabile, e la struttura somiglia più a un documentale, dove molti documenti sono
pertinenti in parte. Una trappola: il file
delle query contiene tutti gli split, e le query da valutare si ricavano dai
giudizi, altrimenti si misura su query che non ne hanno.

Sulle collezioni pubbliche Koskidex gira \emph{piatto}: in memoria, titolo e testo
in un campo solo, refusi spenti come nei riferimenti. La calibrazione, quando
c'è, usa gli split di addestramento e sviluppo, mai quello di test.

\section{Il corpus di dominio: gli albi pretori}\label{sec:corpus}

\paragraph{Perché gli albi.} Il corpus doveva venire dall'archivio di Documentale,
ma il progetto è stato abbandonato e non c'è nessun cliente a cui chiedere il
permesso di usarne i documenti. L'albo pretorio ha la forma di un documentale:
l'archivio di una sola organizzazione, con generi ricorrenti (determine, delibere,
ordinanze, avvisi, liquidazioni), un vocabolario suo e persone che cercano un atto
che sanno esistere. Ed è pubblico per obbligo di legge, quindi riscaricabile da
chiunque legga la tesi. Due fonti, verificate scaricandole il 23 settembre 2026:
\begin{itemize}
  \item \textbf{Comune di Crispiano} (TA), feed RSS, licenza CC BY 4.0: 563 atti
        da novembre 2025 a settembre 2026, 390 determine, 82 delibere, 58
        ordinanze. Il testo è nativo nei PDF, estraibile senza riconoscimento
        ottico: da due a sei pagine, fra 5 e 13 mila caratteri;
  \item \textbf{Regione Friuli Venezia Giulia}, API Socrata, licenza IODL 2.0:
        9.455 atti dal 2011, di 171 enti e 46 tipologie, con ente, tipologia,
        ufficio, oggetto e numero, ma senza testo integrale. I 171 enti si
        comportano come i clienti di un documentale condiviso.
\end{itemize}
Scartato AlboPOP, il progetto che converte gli albi in feed comuni: dei 215 feed
schedati ne rispondono 105, e scaricando a campione gli allegati ne arrivano 3 su
14. Gli altri rispondono \emph{200 OK} con una pagina d'errore: andava provato
scaricando, non contando le righe del catalogo.

\paragraph{Dentro Documentale.} Gli atti entrano in Documentale dal suo comando di
importazione, con il suo modello dei documenti, e in Elasticsearch dal suo comando
di indicizzazione, come in produzione. Il corpus per le valutazioni piatte esce
dall'esportazione di Documentale, non da una ricostruzione: un corpus costruito
scavalcando il sistema non direbbe niente sul sistema. L'indice ha 10.018 atti.

\paragraph{Dati personali.} Gli atti contengono nomi di persone: 79 oggetti del
Friuli Venezia Giulia hanno marcatori espliciti come \emph{nato a} o
\emph{codice fiscale}, e le pubblicazioni di
matrimonio di Crispiano riportano nomi, date e luoghi di nascita. Usarli per
misurare un motore è legittimo, ripubblicarli no: il corpus non entra nel
repository, che ne conserva solo gli script per riscaricarlo, e in questa tesi gli
esempi sono anonimi.

\section{Le query: known-item e bisogni aperti}\label{sec:query}

\paragraph{Known-item automatiche.} Chi sa il numero di un atto e il comune che
l'ha emesso lo ritrova? È la ricerca identificativa più comune in un documentale,
e si può fabbricare senza nessuno che annoti: la query nasce da un atto, e
quell'atto è l'unica risposta giusta. La forma è \texttt{<numero> <comune>}, per
esempio \texttt{13605 Fiume Veneto}. Degli atti con un numero che contiene una
cifra si tengono solo quelli la cui coppia (comune, numero) è unica: la
numerazione riparte ogni anno, e una query che descrive due atti non ha una sola
risposta. Su 9.806 candidati ne restano 9.425, e se ne estraggono 300 con un seme
fisso; altre 300, con un altro seme e nessun atto in comune, servono solo ad
addestrare e calibrare. Misurano una forma di query sola, e fabbricata: dicono se
il motore trova un atto di cui si sa il numero, non come lo cercherebbe una
persona.

\paragraph{Known-item umane.} Per le ricerche di un atto di cui si ricorda il
contenuto servono query scritte da persone. Il protocollo
(appendice~\ref{chapter:appendice-annotazione}) è fissato prima della raccolta: 160 atti
estratti con un seme, metà per fonte e dentro ogni fonte per genere, nessuno già
usato dalle known-item automatiche, divisi in quattro lotti da quaranta. A ogni
persona si mostra un atto come lo vedrebbe un impiegato, lo si nasconde, e le si
chiede cosa scriverebbe fra qualche settimana per ritrovarlo. L'atto si nasconde
prima che la persona scriva, altrimenti la query copia il titolo e la trova
qualunque motore. Scrivono persone che non conoscono i motori; l'autore di questa
tesi no, perché le sue query misurerebbero lui. Al momento di scrivere sono
arrivati due lotti su quattro, 76 query su 80 atti: per quattro atti la
risposta è stata «non saprei», un dato anche questo. Per il protocollo basta un
terzo lotto.

\paragraph{Bisogni aperti.} Per una query come \emph{manutenzione strade} non
esiste un atto giusto solo. Le 24 query del confronto, scritte a mano il 23
settembre 2026 in quattro famiglie (corte, medie, in lingua naturale, per numero
d'atto e con refusi), sono nate per vedere se i due motori trovano le stesse
cose, e sono servite a questo nel capitolo~\ref{chapter:sistemi}. Per misurarne la
pertinenza servono i giudizi della sezione seguente.

\section{Giudizi e accordo fra annotatori}\label{sec:giudizi}

\paragraph{Il pool.} Per le 24 query il pool è l'unione dei primi dieci risultati
di dodici configurazioni: le nove di Koskidex piatto misurate nei
capitoli~\ref{chapter:lessicale} e~\ref{chapter:ibrido}, dal baseline al solo
vettore, e tre passate dall'app, Elasticsearch di produzione e Koskidex innestato
con e senza vettori. Ne escono 905 righe da giudicare.

\paragraph{Il protocollo.} Prima di vedere un solo risultato, per ogni query si
scrive in una o due frasi cosa cerca chi la digita nel documentale di un comune,
e si committa: è l'unica difesa contro un bisogno che diventa «quello che i motori
hanno trovato». Poi si giudica ogni atto su tre gradi: 2 se risponde al bisogno
per intero, 1 se è pertinente ma parziale, 0 se non risponde, anche quando
contiene le parole della query. Le righe sono ordinate per identificativo, non per
posizione, e non dicono quale motore ha trovato l'atto; nel dubbio fra due gradi
vale il più basso, con una nota. Un giudizio non si cambia dopo aver visto le
metriche.

\paragraph{Il secondo annotatore.} Un solo annotatore misura anche il suo modo di
giudicare. Un secondo annotatore giudica, senza vedere i giudizi del primo,
query intere estratte con un seme dal foglio già annotato, almeno 150 righe. Si
riportano il kappa di Cohen sui tre gradi, il kappa pesato linearmente, perché
fra 1 e 2 si sbaglia meno che fra 0 e 2, e l'accordo semplice, comunque vengano.

\paragraph{Com'è andata.} Il protocollo voleva un annotatore umano su tutto il
pool. Non è andata così, e lo scrivo come è andata. I 24 bisogni li ha scritti
un modello linguistico, Claude Opus 5.5, in una sessione separata e prima che
il foglio del pool fosse aperto, e io li ho rivisti prima di committarli. I 905
giudizi li ha dati lo stesso modello, con sette istanze indipendenti, ciascuna
su un gruppo di query, con la scala e le regole dell'appendice~\ref{chapter:appendice-annotazione},
gli atti in ordine di identificativo e senza sapere quale motore li avesse
trovati; il prompt, parola per parola, è nell'archivio accanto ai giudizi.
L'annotatore umano sono stato io, nel ruolo del secondo: il campione estratto
con il seme fissato prima di qualunque giudizio è di quattro query intere, 178
righe, giudicate con la pagina di annotazione senza vedere i giudizi del
modello.

Prima di vedere un mio giudizio avevo scritto le previsioni e una regola:
con un kappa pesato di almeno 0,60 il pool giudicato dal modello avrebbe retto
il capitolo~\ref{chapter:confronto}, sotto 0,40 no. I lavori sui giudici
automatici trovano accordi con gli annotatori umani molto variabili, spesso
moderati~\cite{thomas2023llm,faggioli2023perspectives}. L'accordo semplice è
del 54,5\%, il kappa di Cohen 0,262, quello pesato 0,370: sotto la soglia.
Quattro previsioni su cinque sono cadute, e la più istruttiva è quella sul
verso: avevo previsto un modello più generoso di me, come riporta la
letteratura, ed è più severo, con 107 righe a 0 contro le mie 81. Il
disaccordo sta soprattutto sul grado intermedio: io ho dato 1 a 47 righe, il
modello a 26, e siamo d'accordo su 6. E dipende dalla query: su una il
modello sta sistematicamente un grado sotto di me, su un'altra sopra.

Bisogni e giudice vengono dallo stesso modello, quindi il kappa non misura
solo il giudice: misura anche quanto io leggo i bisogni come li ha scritti
lui. In ogni caso, per la regola fissata prima, i giudizi del modello non
reggono da soli il confronto finale, e il resto del pool non l'ho giudicato:
il capitolo~\ref{chapter:confronto} si regge sulle known-item umane, dove la
risposta giusta è l'atto da cui la query è nata e non serve un giudizio. Il
pool resta al più un confronto secondario, dichiarato come tale e con il kappa
accanto.

\paragraph{Le metriche.} Tre, calcolate dallo stesso codice per ogni motore:
\begin{itemize}
  \item \textbf{nDCG@10}, la metrica principale di BEIR, che pesa i gradi e la
        posizione nei primi dieci;
  \item \textbf{MRR@10}, il reciproco della posizione del primo documento
        pertinente, zero oltre il decimo: con un solo atto giusto, come nelle
        known-item, è la misura naturale;
  \item \textbf{Recall@100}, la quota dei pertinenti nei primi cento.
\end{itemize}
Accanto, sempre, le query a vuoto e il numero di candidati: una media non
distingue un motore che ordina male da uno che non trova, e i due chiedono
correzioni opposte.

\section{Gli strumenti}\label{sec:strumenti}

\paragraph{Koskidex piatto.} \texttt{scripts/evaluate} carica una collezione in
formato BEIR in memoria, la indicizza con le impostazioni date da riga di comando,
esegue le query e scrive metriche, candidati e tempi per query. Con i vettori
calcola gli embedding una volta e li tiene in una cache.

\paragraph{Dall'app.} \texttt{app:eval-run-queries}, un comando aggiunto a
Documentale, esegue un file di query con la ricerca di produzione, sul motore
scelto dalla configurazione, e scrive per ogni query gli identificativi trovati e
il tempo della chiamata. Lo stesso \texttt{scripts/evaluate}, con un'opzione, ne
calcola le metriche: i motori passati dall'app e Koskidex piatto si contano con lo
stesso codice.

\paragraph{Confronto e pool.} \texttt{scripts/compare} esegue le stesse query su
più configurazioni di Koskidex in memoria; \texttt{scripts/pool} costruisce il
foglio dei giudizi dalle configurazioni di Koskidex e dai rapporti dell'app.

\paragraph{Raccolta e giudizi.} Le query umane e i giudizi non passano da un
foglio di calcolo. Due pagine HTML, da aprire nel browser senza rete, fanno da
sole i passi dei protocolli: mostrano un atto alla volta, lo nascondono quando
serve, applicano l'ordine e nascondono il motore, registrano i tempi e le
sessioni, e alla fine scaricano un file da rimandare. La pagina dei giudizi si
rifiuta di esistere finché i bisogni non sono tutti scritti e committati. Due
script importano i file ricevuti, controllano che corrispondano alla pagina da cui
vengono, e scrivono le query e i giudizi in formato BEIR. Le pagine contengono il
testo degli atti e non entrano nel repository.

\paragraph{La macchina.} Tutti i tempi sono misurati su un Apple M2 con 8 core e
16 GB. Koskidex gira nativo su tutti i core; Elasticsearch 9.1.0 e MySQL in una
macchina virtuale Docker con 4 CPU, Elasticsearch con 1 GB di heap, e il suo
tempo include la chiamata HTTP da PHP. I tempi dei due motori non si confrontano
alla pari, e dove stanno fianco a fianco lo si scrive.

\chapter{Il recupero lessicale}\label{chapter:lessicale}

Questo capitolo corregge la parte lessicale di Koskidex, misurandola sulle due
collezioni pubbliche e sulle known-item del capitolo~\ref{chapter:valutazione}.
Il riferimento per le collezioni pubbliche è BM25 di Lucene nella configurazione
\emph{flat} di BEIR~\cite{thakur2021beir}, come la riportano le regressioni di
Pyserini~\cite{pyserini2cr}: 0,6789 di nDCG@10 su SciFact e 0,3218 su NFCorpus.
Koskidex non è tenuto a riprodurli: ha un tokenizer suo, e i riferimenti hanno lo
stemmer e le stopword inglesi di Lucene. Servono a dire quanto manca, e a
chiedersi perché.

Ogni correzione segue la stessa regola: sta dietro un'impostazione con il
comportamento di partenza come default, e le previsioni sono scritte e
committate nel diario di Koskidex prima della misura. Le previsioni sbagliate
restano nel testo, perché qui sono state le più istruttive.

\section{Difetto 0: il recupero congiuntivo}\label{sec:difetto0}

\paragraph{La prima misura.} Il punto di partenza, misurato prima di toccare il
motore, doveva solo collaudare l'impianto. Avevo previsto che SciFact stesse
nettamente sotto il riferimento, fra 0,30 e 0,55, e soprattutto che il richiamo
fosse molto più alto del nDCG: il motore senza IDF avrebbe trovato i documenti e
li avrebbe ordinati male. Avevo anche scritto che, se il richiamo fosse stato
basso, il problema non sarebbe stato l'ordinamento ma il recupero, e la tesi
sarebbe cambiata. È andata così (tabella~\ref{tab:difetto0}): su SciFact
Recall@100 e nDCG@10 sono lo stesso numero, 0,024. Il motore non ordinava male:
non trovava.

\paragraph{La causa.} Ogni termine della query è obbligatorio
(sezione~\ref{sec:koskidex}). Una query di SciFact ha in media 12,5 termini, e
nessun documento li contiene tutti: 290 query su 300 tornano vuote. La prova che
è la lunghezza della query, e non altro, sta nella stessa collezione: su NFCorpus
le 163 query con risultati hanno in media 2,0 parole, le 160 vuote 4,7; su
SciFact le 10 con risultati 8,7, le 290 vuote 12,6. Non è un errore di
programmazione: per la barra di ricerca di un negozio, due o tre parole,
restituire chi le contiene tutte è giusto. È il tipo di query a rivelarlo, ed è
il regime di un documentale, dove una persona scrive in lingua naturale.

Questo difetto non era fra i tre letti nel codice, e viene prima di tutti: con il
97\% di query a vuoto nessun punteggio migliore può fare niente, perché BM25 su un
recupero congiuntivo riordinerebbe il nulla.

\paragraph{La correzione.} Un'impostazione, \texttt{RetrievalMode}, con
\texttt{all} come default e \texttt{any} per il recupero disgiuntivo: un documento
entra se contiene almeno un termine, e il punteggio fa il resto.

\begin{table}[ht]
  \centering
  \begin{tabular}{lrrrr}
    \toprule
    & \multicolumn{2}{c}{SciFact} & \multicolumn{2}{c}{NFCorpus} \\
    \cmidrule(lr){2-3} \cmidrule(lr){4-5}
    & \texttt{all} & \texttt{any} & \texttt{all} & \texttt{any} \\
    \midrule
    nDCG@10 & 0,0246 & \textbf{0,4936} & 0,1659 & \textbf{0,2246} \\
    Recall@100 & 0,0242 & \textbf{0,7571} & 0,0967 & \textbf{0,1898} \\
    MRR@10 & 0,0267 & \textbf{0,4641} & 0,2644 & \textbf{0,3774} \\
    query a vuoto & 290 & \textbf{0} & 160 & \textbf{24} \\
    \midrule
    riferimento BM25 & \multicolumn{2}{c}{0,6789} & \multicolumn{2}{c}{0,3218} \\
    \bottomrule
  \end{tabular}
  \caption{Il difetto 0: recupero congiuntivo (\texttt{all}, il comportamento di
  partenza) e disgiuntivo (\texttt{any}), con il punteggio euristico.}
  \label{tab:difetto0}
\end{table}

SciFact fa venti volte il nDCG di prima e trentuno volte il richiamo; tutte e due
le collezioni arrivano al 70-73\% del riferimento, partendo dal 4\% e dal 52\%.
Due previsioni su cinque sono cadute. SciFact è andato sopra la fascia che avevo
dichiarato. E su NFCorpus avevo previsto che il nDCG potesse peggiorare: le query
sono corte, e pensavo che l'AND facesse da filtro di precisione, con il 50\% di
query a vuoto come prezzo di risposte buone. È migliorato su tutte e tre le
metriche, e il motivo è che il filtro era ridondante con il punteggio che c'era
già: chi trova più termini prende più punti, e i documenti con tutti i termini
restano in testa da soli. L'AND non aggiungeva precisione, toglieva richiamo.

\paragraph{Il residuo.} Le 24 query di NFCorpus ancora vuote sono tutte di un
termine solo e raro: \emph{eggnog}, \emph{halibut}, \emph{okra},
\emph{Fosamax}. Da un campione di dieci avevo concluso che nove non compaiono nel
corpus in nessuna forma e che solo \emph{leeks} sarebbe stata salvata da uno
stemmer. Il campione era troppo piccolo: la sezione~\ref{sec:stopword} ne
recupera nove.

\section{Difetto 1: niente IDF, e BM25}\label{sec:difetto1}

\paragraph{Cosa mancava.} Il punteggio euristico non ha frequenza documentale
inversa, frequenza del termine o normalizzazione sulla lunghezza. BM25
\cite{robertson2009bm25} le chiede tutte e tre, e l'indice non ne aveva nessuna:
la frequenza documentale di un termine non era mantenuta, l'elenco dei termini di
un documento era deduplicato e quindi non dava la lunghezza, e il campo della
frequenza del termine, come detto, non era mai riempito. La correzione le
mantiene durante l'indicizzazione e aggiunge un'impostazione,
\texttt{ScoringMode}, con l'euristico come default. I parametri sono quelli
predefiniti di Lucene, $k_1 = 1{,}2$ e $b = 0{,}75$, non calibrati: qui si misura
la formula standard.

\begin{table}[ht]
  \centering
  \begin{tabular}{lrrrr}
    \toprule
    & \multicolumn{2}{c}{SciFact} & \multicolumn{2}{c}{NFCorpus} \\
    \cmidrule(lr){2-3} \cmidrule(lr){4-5}
    & euristico & BM25 & euristico & BM25 \\
    \midrule
    nDCG@10 & 0,4936 & \textbf{0,6197} & 0,2246 & \textbf{0,2810} \\
    Recall@100 & 0,7571 & \textbf{0,8746} & 0,1898 & \textbf{0,2280} \\
    MRR@10 & 0,4641 & \textbf{0,5868} & 0,3774 & \textbf{0,4710} \\
    \bottomrule
  \end{tabular}
  \caption{Il difetto 1, con il recupero disgiuntivo: il punteggio euristico e
  BM25 ($k_1 = 1{,}2$, $b = 0{,}75$).}
  \label{tab:difetto1}
\end{table}

BM25 porta SciFact al 91\% del riferimento e NFCorpus all'87\%, dentro le fasce
previste (tabella~\ref{tab:difetto1}). Su un piccolo caso di prova separa due
documenti che l'euristico lasciava a pari merito sulla query \texttt{2026/0173}, e
mette davanti il più corto, il cui titolo è quasi solo l'identificativo: la
normalizzazione sulla lunghezza fa il suo mestiere.

\paragraph{La guardia che è scattata.} Avevo scritto una guardia: BM25 cambia come
si ordina, non cosa si recupera, quindi Recall@100 doveva restare entro $\pm0{,}02$.
Si è mosso di +0,117 su SciFact e +0,038 su NFCorpus. Una guardia che scatta vuol
dire che il risultato è falso o che la guardia era scritta male, e le due ipotesi
producono lo stesso sintomo. Tre riscontri dicono la seconda:
\begin{itemize}
  \item il taglio morde: su SciFact tutte le 300 query arrivano a cento risultati,
        su NFCorpus 202 su 323;
  \item sotto il taglio non si muove niente: le 121 query di NFCorpus con meno di
        cento risultati hanno lo stesso richiamo fino al dodicesimo decimale;
  \item il codice non lo consente: il filtro che decide chi resta legge solo
        quanti termini un documento contiene, e il punteggio non elimina nessuno.
\end{itemize}
Recall@100 non misura quanto si è recuperato ma quanti documenti pertinenti sono
entrati nei primi cento, e quando i candidati sono molti più di cento è una
metrica di ordinamento come il nDCG. Su SciFact in recupero disgiuntivo la query
mediana trova 5.182 documenti su 5.183. Il +0,117, letto bene, è un secondo
risultato: BM25 non trova più documenti pertinenti, li porta dentro i primi cento.

Da qui la guardia giusta, che vale per tutto il resto della tesi: il valutatore
registra per ogni query il numero dei candidati prima del taglio. Una modifica al
punteggio non deve cambiarlo per nessuna query; una modifica al recupero lo
cambia, ed è lì che va guardato. Fra il recupero disgiuntivo con l'euristico e
con BM25 non cambia su nessuna delle 623 query.

\section{Le espansioni pesate con il termine trovato}\label{sec:espansioni}

\paragraph{Il sintomo.} Sulle 300 known-item \texttt{<numero> <comune>}, con
Koskidex piatto e il recupero disgiuntivo, BM25 mette l'atto giusto primo nel
61,7\% delle query e l'euristico nel 95,7\%: MRR@10 0,709 contro 0,975. Lo stesso
segno era uscito su \emph{ordinanza 187}.

\paragraph{La causa.} Nelle 115 query in cui BM25 non mette l'atto primo, in 61
il primo classificato non contiene nessun termine della query tale e quale:
combacia solo per prefisso. \texttt{190 Crispiano} è vinta da un atto che contiene
\texttt{1900129}, \texttt{2 Pradamano} da uno che contiene \texttt{2026}.
Koskidex fa la ricerca per prefisso anche con i refusi spenti, e attribuisce il
match al termine trovato; BM25 usa poi la frequenza documentale di quel termine.
Un codice lungo e raro che comincia con le cifre cercate ha un IDF altissimo, e
batte l'atto che contiene il numero esatto, che è più comune. L'euristico non ci
cade perché dà $+2$ ai match esatti.

Tre controfattuali, misurati sulle stesse query e quindi non adottabili, dicono
quanto pesano le alternative ovvie. Togliere la normalizzazione della lunghezza
($b = 0$) porta l'MRR@10 da 0,709 a 0,731, e smentisce la mia previsione che
fosse la causa principale. Spegnere la ricerca per prefisso lo porta a 0,854. Una
saturazione forte della frequenza ($k_1 = 0{,}3$) a 0,900, per ragioni che non ho
misurato.

\paragraph{La correzione.} Lucene conosce lo stesso problema per le espansioni
delle ricerche con refusi, e lo risolve mescolando le frequenze documentali dei
termini espansi (\texttt{TopTermsBlendedFreqScoringRewrite}). Un'impostazione,
\texttt{BM25Expansion}, fa lo stesso: con \texttt{blended}, tutte le espansioni
di un termine della query, il termine esatto compreso, usano la frequenza
documentale più alta fra quelle dei termini che ha trovato. Un codice raro che
comincia con \texttt{190} pesa quanto \texttt{190}, non di più. Cambiano solo i
punteggi: gli insiemi dei candidati restano identici, query per query, e la
guardia lo conferma.

\begin{table}[ht]
  \centering
  \begin{tabular}{lrrr}
    \toprule
    & di serie & frequenza mescolata & differenza \\
    \midrule
    known-item, MRR@10 & 0,7094 & \textbf{0,8331} & +0,124 \\
    SciFact, nDCG@10 & 0,6197 & \textbf{0,6694} & +0,050 \\
    NFCorpus, nDCG@10 & 0,2810 & \textbf{0,3049} & +0,024 \\
    \bottomrule
  \end{tabular}
  \caption{BM25 con le espansioni pesate col termine trovato (di serie) e con la
  frequenza mescolata, senza analisi.}
  \label{tab:espansioni}
\end{table}

Sulle known-item la correzione vale +0,124 ma resta sotto lo spegnimento del
prefisso (0,854) e lontana dall'euristico (tabella~\ref{tab:espansioni}). La
sorpresa è sulle collezioni pubbliche. Avevo previsto che cambiassero meno di
0,01, perché sulle query in inglese le espansioni per prefisso mi sembravano rare;
SciFact guadagna 0,050, cinque volte tanto. Le espansioni non sono rare: sono
dappertutto, e costavano cinque punti di nDCG@10. Quanto, lo dice la sezione
seguente.

Resta, qui, una regola di parità arbitraria: se un documento contiene due
espansioni alla stessa distanza dalla parola cercata, gli viene accreditata la
prima che si incontra, e quindi la sua frequenza nel documento. Con l'ordine in
cui i termini sono entrati nell'indice, quello di queste misure, SciFact fa
0,6694; con l'ordine lessicografico, 0,6670, per due query su 300 in cui un
documento pertinente scende di un posto (sezione~\ref{sec:latenza}). I due
valori sono ugualmente arbitrari: 0,0024 è quanto quella regola pesa, ed è il
margine con cui leggere le differenze di questo ordine su SciFact.

Una regola che non sceglie c'è, ed è quella della \texttt{SynonymQuery} di
Lucene: sommare nel documento le occorrenze di tutte le espansioni e usare la
frequenza documentale più alta. Con un'altra impostazione,
\texttt{BM25Expansion} a \texttt{synonym}, il punteggio non dipende più
dall'ordine, e gli insiemi restano gli stessi. SciFact fa 0,6667, quasi quanto
con l'ordine lessicografico, e con le stopword 0,6734 invece di 0,6757; NFCorpus
e le known-item si spostano di meno di 0,003. Avevo previsto che sommare le
varianti di una parola aiutasse come uno stemmer leggero, e non aiuta. Le due
regole senza arbitrio sono d'accordo fra loro, e lo 0,6694 misurato è il valore,
un po' più alto, dell'ordine di ingresso: i numeri di questo capitolo vanno
letti con quei tre millesimi di margine
(\texttt{2026-09-28\_espansioni-sinonimo}).

\section{Le stopword: un guadagno che veniva da altrove}\label{sec:stopword}

\paragraph{La prima misura.} Per capire quanto del divario dai riferimenti fosse
dell'analizzatore, Koskidex ha preso le stopword inglesi e uno stemmer, misurati
separati e insieme. Avevo previsto che lo stemmer portasse molto più delle
stopword, e le stopword quasi niente, sotto +0,01: con BM25 l'IDF schiaccia già i
termini frequentissimi. È andata al contrario (tabella~\ref{tab:stopword}).

\begin{table}[ht]
  \centering
  \begin{tabular}{lrrrr}
    \toprule
    & \multicolumn{2}{c}{SciFact} & \multicolumn{2}{c}{NFCorpus} \\
    \cmidrule(lr){2-3} \cmidrule(lr){4-5}
    analisi & nDCG@10 & candidati & nDCG@10 & a vuoto \\
    \midrule
    nessuna & 0,6197 & 4.833 & 0,2810 & 24 \\
    solo stopword & \textbf{0,6641} & 2.522 & 0,2900 & 24 \\
    solo stemmer & 0,6243 & 4.908 & 0,2861 & 15 \\
    entrambi & 0,6585 & 3.086 & \textbf{0,2941} & 15 \\
    \bottomrule
  \end{tabular}
  \caption{Stopword e stemmer inglesi con BM25 di serie. \emph{Candidati}: media
  dei documenti trovati per query.}
  \label{tab:stopword}
\end{table}

Le stopword fanno dieci volte lo stemmer su SciFact. La spiegazione che ne avevo
dato era di recupero, non di punteggio: i candidati medi passano da 4.833 a 2.522,
perché in recupero disgiuntivo ogni documento che contiene \emph{the} è un
candidato, e le 25 query senza stopword non guadagnano quasi niente. Lo stemmer
invece salva nove delle 24 query vuote di NFCorpus, quasi tutte plurali il cui
singolare nel corpus c'è: \emph{leeks}, \emph{pineapples}, \emph{turnips},
\emph{bagels}.

\paragraph{La misura che la corregge.} Due giorni dopo la frequenza mescolata ha
mostrato che una stopword di due o tre lettere è anche un prefisso larghissimo:
\emph{the} combacia con \emph{theory}, \emph{in} con \emph{insulin}. Con le
espansioni corrette il guadagno delle stopword su SciFact scende da +0,044 a
+0,006. Era un indizio, non una prova, perché la frequenza mescolata cambia il peso
di tutte le espansioni. La prova è il controfattuale diretto: con la ricerca per
prefisso spenta le stopword portano SciFact da 0,6634 a 0,6664, cioè +0,003.

Il meccanismo, contato: 271 query di test su 300 contengono una stopword di
almeno due lettere, che Koskidex espande. \emph{of} compare in 173 query e in
5.173 documenti su 5.183, ma porta 26 espansioni con frequenza mediana di due
documenti; \emph{in} ne porta 933, di cui 332 in un documento solo. BM25 pesava
ogni espansione con il suo IDF: chiedere \emph{of} voleva dire premiare chi
contiene \emph{offspring}. Le query senza stopword espandibili portano lo 0,7\%
del guadagno. \textbf{Le stopword su SciFact non correggevano un eccesso di peso
delle parole funzionali, ma il difetto delle espansioni}, lo stesso che la
frequenza mescolata corregge per tutti i termini. In questa tesi il loro guadagno
non va letto come un effetto dell'analisi lessicale.

\paragraph{Dove arriva.} Il confronto giusto con i riferimenti, che non fanno
ricerca per prefisso, è con la frequenza mescolata: con le stopword SciFact fa
0,6757, il 99,5\% di 0,6789, e NFCorpus 0,3062, il 95\% di 0,3218. Il difetto 1 è
chiuso. Resta, e non la forzo, una cosa non spiegata: su SciFact aggiungere lo
stemmer alle stopword abbassa il nDCG@10 di 0,006 e alza il Recall@100 da 0,879
a 0,910, mentre su NFCorpus migliora tutti e due.

\section{Identificativi e lingua naturale: nessuna impostazione unica}\label{sec:nessuna-unica}

\paragraph{Il divario che resta.} Con la frequenza mescolata le known-item fanno
0,833 contro lo 0,975 dell'euristico. Prima di scrivere una correzione ho guardato
chi passa davanti all'atto giusto nei 69 fallimenti: in 36 vince un atto senza il
numero, e 35 di questi ripetono il nome del comune più dell'atto giusto; solo in
11 un numero trovato per prefisso. È il recupero disgiuntivo: BM25 somma i
contributi, e un atto con il solo nome del comune, ripetuto, supera uno che ha
numero e comune una volta sola. L'euristico non ci cade perché mette prima chi ha
tutti i termini. Con il recupero congiuntivo BM25 fa 0,960, e col prefisso spento
0,976, cioè quanto l'euristico: \textbf{sulle known-item BM25 non è peggio
dell'euristico, lo è il recupero disgiuntivo}. Il confronto 0,709 contro 0,975
era in realtà fra un recupero disgiuntivo e uno congiuntivo.

\paragraph{Due vie di mezzo.} Il recupero congiuntivo però è il difetto 0. Fra i due
estremi ci sono due impostazioni con un precedente, provate con un valore fissato
prima di misurare (tabella~\ref{tab:intermedio}):
\begin{itemize}
  \item \texttt{MinimumShouldMatch}, con la sintassi di Elasticsearch e il valore
        d'esempio della sua documentazione, \texttt{2<-25\% 9<-3}: fino a due
        termini tutti, da tre a nove tutti meno il 25\%, oltre nove tutti meno tre;
  \item \texttt{Coordination}, il fattore \emph{coord} della similarità classica
        di Lucene, tolto con il passaggio a BM25: il punteggio moltiplicato per la
        quota di termini della query che il documento contiene.
\end{itemize}

\begin{table}[ht]
  \centering
  \begin{tabular}{lrrr}
    \toprule
    BM25, frequenza mescolata & known-item & SciFact & NFCorpus \\
    \midrule
    \texttt{any} & 0,833 & \textbf{0,6694} & \textbf{0,3049} \\
    \texttt{all} & \textbf{0,960} & 0,0246 & 0,1939 \\
    \texttt{any}, \texttt{2<-25\% 9<-3} & 0,889 & 0,2556 & 0,2268 \\
    \texttt{any}, coordinazione & 0,944 & 0,6422 & 0,2935 \\
    \bottomrule
  \end{tabular}
  \caption{Fra il recupero disgiuntivo e il congiuntivo: MRR@10 per le known-item,
  nDCG@10 per le collezioni pubbliche.}
  \label{tab:intermedio}
\end{table}

Due previsioni su cinque hanno tenuto. \texttt{minimum\_should\_match} è
inservibile su tutte e tre: lascia a vuoto 169 query di SciFact e 100 di NFCorpus,
e sulle known-item da tre termini in su ne lascia cadere uno, spesso il numero. Un
vincolo sul numero di termini che va bene per due parole è sbagliato per dodici.
La coordinazione è il compromesso migliore, perché chiude quasi tutto il divario
sulle known-item senza togliere documenti, ma costa 0,027 su SciFact e 0,011 su
NFCorpus: premia chi ha più parole della query, e sulle frasi le parole meno
importanti contano quanto le altre. È la ragione per cui Lucene l'ha tolta.

\textbf{Nessuna impostazione unica va bene per le query identificative e per
quelle in lingua naturale.} Restano tutte spente, e la domanda diventa come
riconoscere il tipo di query e scegliere di conseguenza. Con il solo lessicale la
risposta è una regola di una riga; con i vettori la scelta si allarga al peso del
vettore, e se ne occupa la sezione~\ref{sec:scelta}.

\section{Analisi lessicale italiana}\label{sec:italiana}

Koskidex ha ora i tre stadi dell'analizzatore \texttt{italian} di Elasticsearch,
uguali a quelli di Lucene e tutti spenti per default: il filtro delle elisioni con
gli articoli di Lucene, le stopword italiane di Snowball, e lo stemmer leggero di
Savoy~\cite{savoy2002light}, controllato sulle 35.494 parole del vocabolario di
prova di Lucene.

Delle tre, solo l'elisione è misurata. Con il tokenizer compatibile con
Elasticsearch \emph{dell'infanzia} resta un termine solo, e chi cerca
\emph{infanzia} non trova l'atto: il difetto di produzione della
sezione~\ref{sec:elisioni}. Innestato in Documentale, Koskidex senza elisione
perde le stesse 5.775 coppie di Elasticsearch, e con l'elisione ne perde le stesse
36 di Elasticsearch col filtro. La parità vale per il matching, non per gli
accenti: in 20 parole su 507 Koskidex trova solo di più, perché toglie gli accenti
ed Elasticsearch no, e \emph{identita} trova \emph{identità}.

Stopword e stemmer italiani non si possono misurare con le known-item
automatiche, fatte di un numero e del nome di un comune, che non li mettono alla
prova. Aspettano le known-item scritte da persone.

\section{Scheda contro testo intero, e la lunghezza dei documenti}\label{sec:scheda-testo}

Il corpus degli albi mescola due fonti: 9.455 schede del Friuli Venezia Giulia a
soli metadati, e 563 atti di Crispiano di cui c'è anche il testo intero, fra i 5 e
i 13 mila caratteri. Due domande restano aperte, e si misurano sui soli atti di
Crispiano, perché mescolare le due fonti in un indice solo darebbe una risposta
falsa. La prima: quanto guadagna il recupero quando al posto della scheda c'è il
documento. La seconda: se $b = 0{,}75$, tarato su collezioni di articoli, vada bene
quando le schede corte schiacciano la lunghezza media. Tutte e due chiedono query
scritte da persone su quegli atti, e aspettano le known-item umane.

\chapter{La ricerca in produzione: quattro difetti di Documentale}\label{chapter:produzione}

I difetti di Koskidex dei capitoli~\ref{chapter:lessicale} e~\ref{chapter:ibrido},
tranne il recupero congiuntivo, erano ipotesi scritte leggendo il codice e poi
misurate. Quelli di questo
capitolo no: sono venuti fuori per strada, lavorando all'innesto e fabbricando
le prime query sul corpus di dominio, e sono difetti della ricerca di Documentale
\emph{in produzione}, cioè di Elasticsearch interrogato come lo interroga l'app.
Sono quattro, e tutti colpiscono le ricerche più ordinarie di un documentale:
un atto di cui si sa il numero, una parola italiana preceduta da un articolo,
un numero insieme al nome del comune, una data. Il quarto è venuto fuori per
ultimo, ed è misurato in un altro modo (sezione~\ref{sec:date}).

Il metodo è lo stesso per tutti. Il difetto si misura su Elasticsearch, passando
dall'app o con la sua query esatta; la causa si isola con un controllo che cambia
una cosa sola; la correzione si misura in Koskidex innestato, dietro una variabile
d'ambiente di Documentale spenta per default. Misurarla in Koskidex ha un senso
preciso grazie alla parità del capitolo~\ref{chapter:sistemi}: partendo dagli
stessi insiemi di Elasticsearch, ogni differenza viene dalla correzione.

\section{Il numero d'atto}\label{sec:numero-atto}

In un documentale si cerca spesso un atto di cui si conosce il numero. Due atti
del Comune di Crispiano fanno da esempio: l'\emph{Ordinanza N.~187} e la
\emph{Determina N.~1223}. La tabella~\ref{tab:numero-atto} dà la posizione
dell'atto giusto con la query di produzione e con due varianti che isolano le
cause.

\begin{table}[ht]
  \centering
  \begin{tabular}{lcc}
    \toprule
    Configurazione & \texttt{ordinanza 187} & \texttt{determina 1223} \\
    \midrule
    Elasticsearch, produzione & \textbf{12° su 85} & \textbf{5° su 8} \\
    Elasticsearch senza refusi & 1° su 1 & 1° su 1 \\
    Elasticsearch senza il campo \texttt{name} & 1° su 84 & 5° su 8 \\
    \bottomrule
  \end{tabular}
  \caption{Posizione dell'atto giusto per due query con il numero d'atto.}
  \label{tab:numero-atto}
\end{table}

\paragraph{La causa sono i refusi sui numeri.} Per \texttt{fuzziness: AUTO}
\texttt{187} è una parola di tre caratteri e ammette un refuso: combacia con
\texttt{18}, \texttt{17}, \texttt{137}, \texttt{186}. Spenti i refusi, l'atto
giusto è primo e unico in tutti e due i casi. Un refuso è una tolleranza pensata
per le parole, dove \emph{manutenzone} è quasi certamente \emph{manutenzione}; per
un numero di protocollo una cifra diversa è un altro atto.

\paragraph{Il nome troncato la aggrava.} Prima di indicizzarlo, l'app passa il nome
del documento da \texttt{pathinfo()} per togliere l'estensione del file. L'atto
che vince \texttt{ordinanza 187} si chiama \emph{ORDINANZA N.~18.2026}, e
\texttt{pathinfo()} prende \texttt{.2026} per un'estensione: nell'indice diventa
\emph{ORDINANZA N.~18}, e nel campo che pesa di più, \texttt{name}, \texttt{18}
combacia con \texttt{187} per un refuso. Per la stessa ragione il nome
dell'atto giusto, dentro l'indice, il numero l'ha perso. Tolto quel campo dalla query, per
l'ordinanza l'atto giusto torna primo; per la determina no, perché lì i
concorrenti hanno il numero vicino nel riassunto.

Due atti sono un esempio, non una misura. La sezione~\ref{sec:numero-comune}
ripete la domanda su trecento atti, e lì il difetto dei refusi sui numeri
ricompare, dietro un altro.

\section{Le elisioni}\label{sec:elisioni}

\paragraph{Il difetto.} Il tokenizer standard di Elasticsearch tiene dentro la
parola un apostrofo fra due lettere (sezione~\ref{sec:differenze}): nell'indice
\emph{dell'illuminazione} è un termine solo, e chi cerca \emph{illuminazione} non
trova un atto che la nomina solo così. L'analizzatore \texttt{italian} di
Elasticsearch ha un filtro apposta, \texttt{elision}, con l'elenco degli articoli
italiani~\cite{es-italian}, ma l'indice di Documentale usa l'analizzatore
predefinito.

\paragraph{Quanti atti tocca.} La misura è fatta sull'indice che l'app ha costruito,
senza ricostruire il tokenizer. Dai termini che Elasticsearch ha scritto davvero
si prendono quelli elisi (lettere, apostrofo, lettere) e per ognuno la parola
dopo l'apostrofo. Una coppia (atto, parola) è \emph{a rischio} se nell'atto la
parola compare solo elisa. È \emph{persa} se la query di produzione con quella sola
parola non trova l'atto: si interroga il motore, invece di contare i termini,
perché i refusi o un plurale nello stesso atto potrebbero salvarla. Il
controllo della causa è un indice temporaneo con lo stesso mapping più il filtro
\texttt{elision}: una coppia che lì torna era persa per l'elisione e non per
altro. Le coppie in cui la forma elisa sta solo nel nome dell'ente
(\emph{Gradisca d'Isonzo}) si contano anche a parte, perché chi cerca un comune ne
scrive probabilmente il nome con l'apostrofo.

\begin{table}[ht]
  \centering
  \begin{tabular}{lrr}
    \toprule
    & atti & \% \\
    \midrule
    nell'indice & 10.018 & \\
    con almeno un termine eliso & 4.280 & 42,7 \\
    \textbf{con almeno una coppia persa} & \textbf{4.033} & \textbf{40,3} \\
    \quad escluso il solo nome dell'ente & 3.826 & 38,2 \\
    \midrule
    & coppie & \\
    \midrule
    a rischio & 5.913 & \\
    salvate dalla ricerca (refusi, forme vicine) & 138 & 2,3 \\
    \textbf{perse} & \textbf{5.775} & \\
    perse che il filtro \texttt{elision} recupera & 5.739 & 99,4 \\
    \bottomrule
  \end{tabular}
  \caption{Le elisioni nell'Elasticsearch di produzione, su 507 parole.}
  \label{tab:elisioni}
\end{table}

\textbf{Quattro atti su dieci hanno almeno una parola che la ricerca di quella
parola non trova} (tabella~\ref{tab:elisioni}). Non dipende dalla fonte: a
Crispiano sono il 34,6\% degli atti, nel Friuli Venezia Giulia il 40,6\%. È la
lingua. La ricerca non salva quasi niente, il 2,3\% delle coppie, e il filtro ne
recupera il 99,4\%: la causa è quella. Delle 36 coppie che restano, 17 sono
elisioni fuori dall'elenco degli articoli (\emph{sant'}, \emph{cinquant'},
\emph{quest'}), 13 sono parole attaccate per uno spazio mancante
(\emph{scuoladell'infanzia}), le altre accenti scritti come apostrofo
(\emph{conformita'}) e prefissi fuori elenco.

Le parole più colpite dicono chi se ne accorge. \emph{Infanzia} è persa in tre atti
su quattro fra quelli che la nominano (168 su 220), perché si scrive quasi sempre
\emph{scuola dell'infanzia} o \emph{nido d'infanzia}; \emph{associazione} in 170 su
230; \emph{art}, cioè i rinvii normativi come \emph{dell'art.~36}, in 595 su 1.406.
\emph{Illuminazione}, la parola da cui tutto è partito, è fra le meno colpite: 12
atti su 99.

\paragraph{La correzione, nei due motori.} In Elasticsearch è il filtro
\texttt{elision} nel mapping e una reindicizzazione. In Koskidex innestato è
un'impostazione, \texttt{ElisionArticles}, con gli stessi 21 articoli
dell'analizzatore \texttt{italian}, che Documentale accende con
\texttt{KOSKIDEX\_ELISION=true}. Misurata con le stesse 507 parole:
\begin{itemize}
  \item senza, Koskidex perde \textbf{le stesse 5.775 coppie} di Elasticsearch in
        produzione, non solo lo stesso numero;
  \item con, ne perde \textbf{le stesse 36} di Elasticsearch col filtro;
  \item sulle 24 query del confronto nessuna query perde un documento, e 8 ne
        guadagnano 65 in tutto: \emph{illuminazione pubblica} passa da 62 a 73
        atti, \emph{contributo associazioni} da 17 a 48;
  \item sulle 300 query known-item della sezione seguente non cambia niente, nemmeno
        l'ordine dei primi dieci.
\end{itemize}
Una differenza fra i due motori resta, ed è a favore di chi cerca: Koskidex toglie
gli accenti, Elasticsearch con l'analizzatore predefinito no. Su 20 delle 507
parole Koskidex trova qualcosa di più, 195 atti in tutto, 67 dei quali per
\emph{unità}: chi scrive \emph{identita} senza accento trova \emph{identità}.

\section{Numero e comune: il campo unico e i refusi sui numeri}\label{sec:numero-comune}

\paragraph{Le query.} Il caso della sezione~\ref{sec:numero-atto} si può
fabbricare in grande senza che nessuno annoti. Da 300 atti estratti a caso con un
seme fisso si scrive la query \texttt{<numero> <comune>}, per esempio
\texttt{13605 Fiume Veneto}; l'atto di partenza è l'unica risposta giusta. Si
tengono solo gli atti la cui coppia (comune, numero) è unica nel corpus, perché
la numerazione riparte ogni anno. Le query e il loro uso sono descritti nel
capitolo~\ref{chapter:valutazione}; qui la metrica principale è l'MRR@10, con la
quota di query con l'atto al primo posto, entro i primi dieci e a vuoto.

\paragraph{Il difetto.} In produzione \textbf{l'atto giusto è entro i primi dieci
nel 2\% delle query}, e sette query su dieci non trovano niente
(tabella~\ref{tab:numero-comune}). Nelle 89 che trovano qualcosa l'atto giusto c'è
in 6. La causa non è il numero ma come la query combina i campi: con
\texttt{best\_fields} e \texttt{operator: and} tutte le parole devono stare nello
stesso campo, ma il numero sta in \texttt{additional\_data} e il comune in
\texttt{subjects}, e nessun campo li ha tutti e due. Il controllo lo conferma:
riscritta la query con un \texttt{multi\_match} per parola, con gli stessi refusi,
la stessa prima lettera esatta, gli stessi pesi e tutte le parole ancora
obbligatorie, cade solo il vincolo del campo unico, e l'atto giusto torna in
tutte e 300 le query. Koskidex innestato, con le impostazioni di compatibilità,
trova gli stessi insiemi di Elasticsearch su tutte e 300: la parità eredita
anche questo difetto.

\paragraph{Due correzioni, perché i difetti sono due.} Spento in Koskidex
innestato il solo vincolo del campo unico
(\texttt{KOSKIDEX\_ALL\_TERMS\_IN\_ONE\_FIELD=false}), l'atto torna entro i primi
dieci nel 93\% delle query e nessuna resta a vuoto, ma è primo solo nel 38\%.
Nelle 185 query in cui non è primo, in 154 lo scavalca un atto con il numero a un
refuso: per \texttt{1209 Crispiano} vince un atto con \texttt{1109}. È il difetto
della sezione~\ref{sec:numero-atto}, che il vincolo del campo unico nascondeva.
La seconda correzione è un'impostazione di Koskidex con lo stesso significato
di \texttt{disableOnNumbers} di Meilisearch~\cite{meili-typo},
\texttt{disable\_on\_numbers}: un termine della query fatto di sole cifre non
ammette refusi, mentre i codici misti (\texttt{a651}) li tengono. In Documentale
la accende \texttt{KOSKIDEX\_TYPOS\_ON\_NUMBERS=false}.

\begin{table}[ht]
  \centering
  \small
  \begin{tabular}{lrrrr}
    \toprule
    Dall'app & MRR@10 & primo & entro 10 & vuote \\
    \midrule
    Elasticsearch, produzione & 0,017 & 1,3\% & 2,0\% & 70\% \\
    Koskidex, campo unico & 0,018 & 1,7\% & 2,0\% & 70\% \\
    Koskidex, campo unico, n.\ esatti & 0,020 & 2,0\% & 2,0\% & 95\% \\
    Koskidex, parole libere & 0,548 & 38,3\% & 93,3\% & 0 \\
    \textbf{Koskidex, parole libere, n.\ esatti} & \textbf{0,950} & \textbf{91,7\%} & \textbf{99,7\%} & \textbf{0} \\
    \midrule
    Elasticsearch, parole libere (controllo) & 0,534 & 36,3\% & 96,0\% & 0 \\
    \bottomrule
  \end{tabular}
  \caption{Le 300 query \texttt{<numero> <comune>}. \emph{Campo unico}: tutte le
  parole nello stesso campo, come in produzione; \emph{parole libere}: ciascuna in
  un campo qualsiasi; senza \emph{numeri esatti} (n.\ esatti) i numeri ammettono refusi. Le
  prime cinque righe passano dall'app; il controllo interroga Elasticsearch con la
  query riscritta parola per parola.}
  \label{tab:numero-comune}
\end{table}

\textbf{Le due correzioni insieme portano l'atto giusto al primo posto in 275
query su 300}, da 4. Nessuna delle due basta da sola: il vincolo del campo unico
decide se l'atto si trova, i refusi sui numeri se è primo. La riga con il solo
campo unico e i numeri esatti dice anche un'altra cosa: togliere i refusi ai
numeri, lasciando il vincolo, porta le query a vuoto dal 70\% al 95\%. Quel 30\%
che in produzione rispondeva, rispondeva soprattutto con numeri sbagliati.

Nelle 25 query in cui l'atto giusto resta dietro, in 21 il primo classificato ha
lo stesso numero, esatto, in un campo che pesa di più (il numero dell'atto giusto
sta in \texttt{additional\_data}, che pesa 1); in 3 lo ha esatto in un campo che
non pesa di più; in una sola non lo ha esatto, ed è l'unico codice misto fra i
perdenti. Sulle 24 query del confronto il campo libero allarga 14 insiemi su 24,
di 375 atti in tutto, senza toglierne nessuno; i numeri esatti cambiano poi
soltanto le due query con un numero d'atto: \texttt{ordinanza 187} e
\texttt{determina 1223} restituiscono l'atto giusto, primo e unico.

\section{Le date, in tre formati}\label{sec:date}

Nelle schede degli atti la stessa data compare in tre formati, tutti frequenti:
\texttt{14/01/2026} (3.097 volte, in 1.465 atti), \texttt{14 gennaio 2026}
(1.834 volte, in 949) e \texttt{14.01.2026} (2.103 volte, in 982). Gli importi
hanno quasi sempre i punti delle migliaia (312 volte, in 201 atti), e quasi mai
no (2 volte, in 2). Il tokenizer standard di Elasticsearch tiene
\texttt{14.01.2026} come un termine solo e spezza \texttt{14/01/2026} in tre, e
nessun tokenizer sa che \emph{gennaio} è \texttt{01}: chi scrive una data in un
formato non trova gli atti che la scrivono in un altro. Il tokenizer di
Koskidex spezza tutto, e collega barra e punto, ma il giorno di ogni data
diventa un numero che combacia con qualunque ricerca per numero.

\paragraph{Come si misura, e perché diversamente.} Qui il metodo del capitolo
non si può seguire: una query \texttt{<data> <comune>} su Elasticsearch di
produzione andrebbe quasi sempre a vuoto per il difetto della
sezione~\ref{sec:numero-comune}, perché il comune non sta nel campo della data,
e misurerebbe quello. Si misura quindi Koskidex piatto con il tokenizer
standard, che è quello di Elasticsearch, e con il suo; per la parità del
capitolo~\ref{chapter:sistemi} il primo si comporta come Elasticsearch sui
termini. Le query sono sintetiche, come le known-item automatiche: da un atto
si prende una data che in quel comune compare in quell'atto solo, e la si
scrive nei tre formati, cinquanta atti per formato di partenza; allo stesso modo
gli importi, con e senza i punti delle migliaia, su 52 atti: 50 con i punti e i
soli 2 che non li hanno.

\paragraph{Il difetto.} Con il tokenizer standard la data nel formato dell'atto
trova l'atto entro i primi dieci nel 98-100\% delle query; in un altro formato
nello 0-8\%, quasi sempre a vuoto. Con il tokenizer di Koskidex barra e punto si
trovano fra loro (94\%), il mese in lettere no (0-8\%). Gli importi con e senza
punti non si trovano mai.

\paragraph{La correzione.} Due impostazioni di Koskidex, spente per default,
riscrivono il testo degli atti e delle query prima del tokenizer: ogni data
valida, in uno dei quattro formati e con i mesi in italiano, diventa un numero
solo, \texttt{aaaammgg}, e agli importi si tolgono i punti delle migliaia. Con
tutti e due i tokenizer ogni coppia di formati trova l'atto entro i primi dieci
nel 98-100\% dei casi, gli importi nel 100\%, e le 300 ricerche
\texttt{<numero> <comune>} non peggiorano (MRR@10 da 0,975 a 0,977). Due
previsioni su otto sono cadute. I numeri piccoli, da 1 a 31, perdono meno
candidati del previsto, perché il comune nella query li restringe già. E
l'indicizzazione costa il 29-34\% in più: avevo previsto al più il 20\%, e
saltare i testi che non possono contenere un formato l'ha solo ridotta dal
41-46\%, perché quasi ogni scheda contiene davvero cifre, barre e punti.

Le query sono fabbricate: dicono che il motore trova una data in qualunque
formato, non quanto spesso una persona cerca per data, e nelle known-item umane
raccolte finora non succede mai: nessuna delle 76 contiene una data. E la correzione ha un prezzo che nessuna
previsione misurava: la data diventa un termine solo, quindi il suo anno non
combacia più con una ricerca che contiene l'anno da solo, come \emph{bilancio
2026}, a meno che l'anno compaia anche fuori da una data.

\paragraph{Dall'app.} In Documentale le due impostazioni sono due variabili,
spente e fuori dal profilo della sezione seguente, misurate dall'app con il
profilo acceso (\texttt{2026-09-28\_date-app}). Le date si comportano come in
Koskidex piatto: da 0-8\% a 98-100\% di atti entro i primi dieci in un altro
formato, MRR@10 da 0,315 a 0,979, con il 6\% in più di indicizzazione; le
known-item automatiche migliorano appena, le 76 umane raccolte finora e le 21
query del confronto senza cifre non cambiano di un posto. Gli importi no, e
per una ragione che Koskidex piatto non poteva mostrare, perché
\texttt{scripts/evaluate} cerca senza refusi e l'app con la tolleranza
automatica: nell'app gli importi si trovano fra formati anche spenti,
perché \texttt{5056,03} dista un carattere da \texttt{5.056,03}, e un termine
con la virgola non è fatto di sole cifre, quindi i numeri esatti del
profilo non lo riguardano. Lo stesso refuso fa combaciare importi diversi, e
togliere i punti peggiora le cose: gli importi diventano più corti e più
vicini, e l'MRR@10 degli importi scende da 0,917 a 0,891. Il difetto vero
degli importi è che ammettono refusi.

\paragraph{Gli importi senza refusi.} È il difetto della
sezione~\ref{sec:numero-atto} su un'altra forma di numero, e la correzione è la
stessa: un'impostazione di Koskidex, \texttt{disable\_on\_amounts}, spenta per
default, nega i refusi ai termini fatti di cifre, punti e virgole, e in
Documentale è una variabile fuori dal profilo
(\texttt{2026-09-28\_importi-esatti}). Da sola toglie il ponte che il refuso
dava per caso, e gli importi fra formati non si trovano più (0 su 52); insieme
alla normalizzazione degli importi si trovano tutti (52 su 52), l'MRR@10 degli
importi sale da 0,917 a 0,990, e i risultati sulle 104 query scendono da 175 a
106: quelli tolti sono atti con un importo diverso. Sette previsioni su otto;
quella caduta è caduta in meglio. Anche \texttt{14.01.2026} è un termine con i
punti, che prima ammetteva refusi e trovava \texttt{14.01.2025}: con
l'impostazione 14 ricerche di date col punto migliorano, nessuna peggiora, e i
risultati sulle query delle date scendono da 594 a 241. Le known-item
automatiche, le umane raccolte finora e le 21 query del confronto senza cifre
non cambiano.

Nell'app quindi le date vanno normalizzate, e gli importi normalizzati ed
esatti. Del prezzo dell'anno da solo le query raccolte finora non dicono niente:
nessuna delle umane contiene un anno, e l'unica del confronto che lo contiene
non trova niente in nessun caso. Per questo le tre impostazioni restano fuori
dal profilo finché non girano le known-item umane complete.

\section{Cosa cambia per chi cerca}\label{sec:cosa-cambia}

Le tre correzioni misurate dall'app sono misurate una alla volta, e una configurazione fatta di
pezzi misurati separatamente potrebbe non fare la somma dei pezzi. In Documentale
esiste per questo un profilo, \texttt{KOSKIDEX\_PROFILO=consigliata}, che fa
partire insieme le tre variabili dal valore corretto, senza cambiare nessun
default: senza profilo l'app si comporta come in produzione, e ogni variabile
impostata da sola vince sul profilo. Il profilo è misurato come un tutto, da un
indice vuoto, con le stesse query:
\begin{itemize}
  \item sulle 300 query known-item fa \textbf{esattamente} quello che facevano
        parole libere e numeri esatti, fino all'ordine dei primi dieci risultati di
        ogni query: l'elisione non sposta niente;
  \item sulle 24 query del confronto, rispetto a parole libere e numeri esatti,
        nessuna query perde un documento e 7 ne guadagnano, quelle toccate
        dalle elisioni (l'ottava della sezione~\ref{sec:elisioni},
        \texttt{ordinanza 187}, con i numeri esatti ha già il solo atto giusto);
  \item \texttt{ordinanza 187} e \texttt{determina 1223} danno il solo atto giusto.
\end{itemize}

\begin{table}[ht]
  \centering
  \begin{tabular}{lrr}
    \toprule
    & produzione & consigliata \\
    \midrule
    \texttt{ordinanza 187} & 12° su 85 & 1° su 1 \\
    \texttt{determina 1223} & 5° su 8 & 1° su 1 \\
    \texttt{<numero> <comune>}, atto primo & 1,3\% & 91,7\% \\
    \texttt{<numero> <comune>}, a vuoto & 70\% & 0 \\
    coppie (atto, parola) perse per l'elisione & 5.775 & 36 \\
    atti con una parola introvabile & 40,3\% & 0,4\% \\
    \bottomrule
  \end{tabular}
  \caption{La ricerca di Documentale in produzione e con il profilo consigliato di
  Koskidex innestato. Le righe delle elisioni vengono dalla misura delle 507
  parole della sezione~\ref{sec:elisioni}.}
  \label{tab:cosa-cambia}
\end{table}

La tabella~\ref{tab:cosa-cambia} riassume cosa cambia. Il numero degli atti con una
parola introvabile, con il profilo, è quello delle 36 coppie residue, che stanno in
36 atti diversi.

\paragraph{Cosa non dice.} Tutte queste misure riguardano \emph{se} l'atto si
trova e, per le query identificative, \emph{dove}. Non dicono niente dell'ordine
dei risultati sulle ricerche in lingua naturale, come \emph{manutenzione strade}:
lì non c'è un atto giusto solo, e per misurare servono giudizi di pertinenza, che
sono l'oggetto del capitolo~\ref{chapter:confronto}. Non dicono neanche che
Elasticsearch non si possa correggere: il filtro delle elisioni e la query
riscritta parola per parola sono misurati anche lì, con lo stesso effetto. Per i
numeri servirebbe separare nella query i termini fatti di sole cifre e cercarli
senza refusi, e questo in Elasticsearch non è stato misurato.

\paragraph{Cosa vuol dire.} I quattro difetti hanno una cosa in comune: nessuno è
un errore di programmazione. Il recupero congiuntivo per campo, la tolleranza ai
refusi uguale per parole e numeri, il tokenizer che rispetta l'apostrofo e tiene
insieme le date con il punto sono tutte impostazioni ragionevoli di un motore
generico, e ciascuna si rompe su una proprietà precisa dei documenti
amministrativi italiani: i metadati sparsi in campi diversi, i numeri di
protocollo, gli articoli elisi, le date scritte in tre modi. Per trovarle è bastato
fare le ricerche che fa chi usa un documentale e contare, con il codice dell'app,
quante andavano a vuoto.

\chapter{Recupero ibrido e fusione}\label{chapter:ibrido}

Il capitolo~\ref{chapter:lessicale} ha portato il recupero lessicale di Koskidex
vicino ai riferimenti pubblicati. Restano i due difetti che riguardano i vettori.
Il difetto 2: il vettore della query riordina i documenti che il lessicale ha già
trovato, ma non ne porta di nuovi. Il difetto 3: il punteggio lessicale e la
similarità si sommano con una costante, 20, messa a occhio, su due scale che non
hanno niente in comune. Questo capitolo li misura tutti e due, ne misura il costo,
e arriva a una conclusione che nessuno dei due, da solo, faceva prevedere: il modo
giusto di usare il vettore dipende dal tipo di query.

\section{Gli embedding e la suddivisione dei documenti lunghi}\label{sec:embedding}

\paragraph{Il modello.} Gli embedding vengono da un modello locale, \texttt{bge-m3}
\cite{chen2024bgem3}, servito da Ollama sulla stessa macchina. La scelta di un
modello locale invece di un servizio esterno ha tre ragioni: i pesi e la versione
sono fissi, e ogni esito registra il nome e l'impronta del modello
(\texttt{790764642607}); i documenti non escono dalla macchina; e un motore che si
vanta di non avere dipendenze non guadagnerebbe niente a dipendere da un servizio
a pagamento. \texttt{bge-m3} è multilingue, dà vettori di 1.024 dimensioni già
normalizzati, ha un contesto di 8.192 token e non chiede prefissi diversi per
query e documenti.

\paragraph{Nel motore.} Koskidex ha un'impostazione, \texttt{Embedder}, spenta per
default: accesa, il server calcola il vettore di ogni documento quando lo aggiunge
e quello della query quando la ricerca lo chiede, e li tiene in una cache indicizzata
dal modello e dal testo, così che uno stesso testo si calcoli una volta sola. Il
motore vede solo i vettori.
Documentale la accende con una variabile d'ambiente,
\texttt{KOSKIDEX\_EMBEDDER\_MODEL}. Il testo di cui si calcola il vettore è quello
indicizzato: nelle valutazioni piatte titolo e testo in un campo solo, in
Documentale i sei campi della scheda. Calcolare la prima volta i vettori richiede
minuti per collezione; dalla cache, frazioni di secondo.

\paragraph{Quanto testo arriva al vettore.} Il contesto di 8.192 token è quello
del modello, non quello che Ollama usa: senza opzioni nella richiesta, Ollama
legge i primi 2.048 token di un testo e scarta il resto. Chiedere un contesto più
grande con \texttt{num\_ctx} non basta, perché Ollama taglia al minimo fra
\texttt{num\_ctx} e la dimensione del batch, \texttt{num\_batch}, che vale anch'essa
2.048. L'ho scoperto tardi, e per caso: una prova su dati inventati dava numeri
identici con e senza il contesto più grande. L'esperimento che lo misura
(\texttt{2026-09-25\_contesto-ollama}) conta i token letti documento per
documento. Nessuna delle 9.455 schede del Friuli Venezia Giulia arriva alla
soglia (la più lunga ha 486 token), né quelle di Crispiano, che non superano i
1.058 caratteri; e in
SciFact e NFCorpus la superano 5 documenti su 8.816, di poche centinaia di token:
i vettori delle misure di questo capitolo hanno letto tutto il testo, o quasi. Dei 563 testi integrali
degli atti di Crispiano, invece, il 59\% supera i 2.048 token (mediana 2.168), e
solo 3 arrivano a 8.192. Koskidex ha ora un'impostazione, \texttt{embedder.context},
che manda a Ollama \texttt{num\_ctx} e \texttt{num\_batch} con lo stesso valore;
a zero la richiesta resta quella di sempre.

\paragraph{La suddivisione dei documenti lunghi.} Un modello di embedding ha un
contesto limitato, e i documenti più lunghi di quel contesto vanno spezzati in
parti, con un vettore per parte. Per le schede, che Documentale indicizza oggi,
la questione non si pone. Per il testo integrale, con il contesto portato a 8.192
token, resterebbero tagliati 3 atti su 563. Se leggere tutto il testo, o
spezzarlo, migliori la pertinenza l'ha misurato la scheda contro il testo
intero (sezione~\ref{sec:scheda-testo}), sulle 50 known-item umane di
Crispiano: con i vettori dei primi 8.192 token A fa 0,515 di MRR@10, con quelli
dei primi 2.048 fa 0,558, e sulla scheda 0,613. Quello che sta dopo i primi
2.048 token, su queste query, pesa contro; per la regola scritta prima della
misura, spezzare i documenti lunghi in parti non serve.

\section{Difetto 2: dal re-ranking al recupero ibrido}\label{sec:difetto2}

\paragraph{Il difetto.} Con il vettore della query, Koskidex aggiunge
$20 \cdot \text{sim}$ al punteggio di ogni documento già trovato dal lessicale. Un
documento pertinente che non condivide parole con la query non entra mai: dal vivo,
con l'embedder acceso, \emph{traffico} non trova l'ordinanza di chiusura di una
strada, né da sola né in ibrido.

\paragraph{La correzione.} Un'impostazione, \texttt{HybridMode}, con il
comportamento di partenza come default:
\begin{itemize}
  \item vuota (\textbf{R}, re-ranking): il vettore riordina i candidati lessicali;
  \item \texttt{union} (\textbf{U}): i \texttt{VectorTopK} documenti più vicini per
        similarità del coseno entrano fra i candidati, accanto a quelli lessicali;
  \item \texttt{vector} (\textbf{V}): solo i più vicini, il lessicale ignorato. È il
        riferimento: quanto vale il modello da solo.
\end{itemize}
\texttt{VectorTopK} vale 100, fissato prima di misurare perché è la profondità di
Recall@100. La fusione non cambia: un documento portato solo dal vettore vale
$20 \cdot \text{sim}$, come nella ricerca solo vettoriale; qui cambia solo chi
entra fra i candidati. La scansione dei vettori è esaustiva, senza indice
approssimato: il suo costo si misura nella sezione~\ref{sec:costo}.

\paragraph{Con il recupero disgiuntivo.} La base lessicale (\textbf{L}) è quella
corretta del capitolo~\ref{chapter:lessicale}: BM25, recupero disgiuntivo,
frequenza mescolata. La tabella~\ref{tab:difetto2-any} dà le quattro configurazioni.

\begin{table}[ht]
  \centering
  \begin{tabular}{lrrrr}
    \toprule
    & L & R & U & V \\
    \midrule
    SciFact nDCG@10 & 0,6694 & \textbf{0,6913} & 0,6913 & 0,6436 \\
    SciFact Recall@100 & 0,8859 & 0,9116 & 0,9116 & 0,9037 \\
    NFCorpus nDCG@10 & 0,3049 & 0,3284 & \textbf{0,3340} & 0,3149 \\
    NFCorpus Recall@100 & 0,2391 & 0,2499 & \textbf{0,2809} & 0,2830 \\
    NFCorpus, query a vuoto & 24 & 24 & \textbf{0} & 0 \\
    known-item MRR@10 & 0,8331 & \textbf{0,8464} & 0,8464 & 0,1837 \\
    \bottomrule
  \end{tabular}
  \caption{Il difetto 2 con il recupero disgiuntivo: lessicale (L), re-ranking
  (R, il comportamento di partenza), unione (U) e solo vettore (V).}
  \label{tab:difetto2-any}
\end{table}

Il re-ranking vale più di due punti su tutte e due le collezioni pubbliche,
\emph{con} la fusione a occhio del difetto 3. L'unione aggiunge qualcosa solo dove
il lessicale lascia buchi. Su SciFact non cambia niente, e la ragione è istruttiva:
con il recupero disgiuntivo il lessicale trova in mediana 5.182 documenti su 5.183,
i cento più vicini sono già fra i candidati, e il vettore in mediana non ne porta nessuno. Il
difetto 2 lì non si vede perché il difetto 0 è stato corretto dalla parte
opposta: quando il lessicale prende tutto, non c'è niente da recuperare. Su
NFCorpus, dove le query sono corte, il lessicale trova in mediana 560 documenti e
lascia 24 query senza risultati; il vettore porta in mediana 44 candidati nuovi, il
richiamo sale di 0,031 e le query a vuoto spariscono. Il modello da solo fa meglio
del lessicale su NFCorpus e peggio su SciFact, e sulle known-item, dove tutto sta in
un numero, non sa quasi niente (0,18).

\paragraph{Con il recupero congiuntivo.} Documentale cerca in modo congiuntivo, e lì
il lessicale lascia a vuoto la maggior parte delle query lunghe. La stessa misura,
cambiando solo il recupero, dà la tabella~\ref{tab:difetto2-all}.

\begin{table}[ht]
  \centering
  \begin{tabular}{lrrrr}
    \toprule
    Recupero congiuntivo & L & R & U & U disgiuntivo \\
    \midrule
    SciFact nDCG@10 & 0,0246 & 0,0246 & \textbf{0,6483} & 0,6913 \\
    SciFact, query a vuoto & 290 & 290 & \textbf{0} & 0 \\
    NFCorpus nDCG@10 & 0,1939 & 0,1973 & \textbf{0,3333} & 0,3340 \\
    NFCorpus, query a vuoto & 160 & 160 & \textbf{0} & 0 \\
    known-item MRR@10 & 0,9596 & 0,9579 & 0,9579 & 0,8464 \\
    \bottomrule
  \end{tabular}
  \caption{Il difetto 2 con il recupero congiuntivo di Documentale. L'ultima colonna
  è l'unione della tabella~\ref{tab:difetto2-any}, per confronto.}
  \label{tab:difetto2-all}
\end{table}

\textbf{Con il recupero congiuntivo il difetto 2 è il più grande dei difetti di
Koskidex}: l'unione vale +0,62 di nDCG@10 su SciFact e +0,14 su NFCorpus, e porta a
zero le query a vuoto. Il re-ranking non può niente: riordina i pochi documenti che
il lessicale ha trovato. E c'è un risultato che non avevo previsto, nell'ultima
riga. Il recupero congiuntivo con l'unione tiene le known-item a 0,958, dove quello
disgiuntivo con l'unione faceva 0,846, e sulle collezioni pubbliche resta vicino al
disgiuntivo ($-$0,043 su SciFact, $-$0,001 su NFCorpus). Il motivo per cui il recupero
disgiuntivo serviva alle query in lingua naturale era che il congiuntivo le
lasciava vuote; l'unione toglie quel motivo, perché i documenti li porta il
vettore. È anche la configurazione più vicina a Documentale, che cerca già in modo
congiuntivo: alla ricerca di produzione, per correggere il difetto 2, basta
aggiungere i vettori.

\section{Difetto 3: fondere scale incomparabili}\label{sec:difetto3}

\paragraph{Il difetto.} La somma $\text{BM25} + 20 \cdot \text{sim}$ mette insieme un
punteggio che cresce con la lunghezza della query e una similarità fra $-1$ e $1$,
con una costante scelta a occhio. Su una query di una parola il vettore pesa molto,
su una di dodici quasi niente.

\paragraph{Tre fusioni.} Dietro un'impostazione, \texttt{FusionMode}, con la somma
di partenza come default:
\begin{itemize}
  \item la \textbf{somma}, con la costante in un'impostazione, \texttt{VectorWeight},
        che vuota vale 20;
  \item la \textbf{Reciprocal Rank Fusion}~\cite{cormack2009rrf}, che somma
        $1/(60 + r)$ per ciascuna delle due liste, dove $r$ è la posizione del
        documento in quella lista: non guarda i punteggi, solo le posizioni. La
        costante 60 è quella dell'articolo, predefinita anche in Elasticsearch;
  \item la \textbf{combinazione convessa}: ciascuna lista normalizzata min-max sui
        propri candidati, poi $\alpha \cdot \text{lessicale} + (1-\alpha) \cdot
        \text{vettore}$, con $\alpha$ in un'impostazione che vuota vale 0,5.
\end{itemize}
Le due fusioni nuove lavorano sull'unione dei candidati, e un documento assente da
una lista prende zero da quella lista.

\paragraph{Calibrare senza guardare il test.} La costante della somma e $\alpha$
si scelgono su split separati da quelli di test: SciFact train (809 query), NFCorpus
dev (324) e 300 known-item generate da atti diversi da quelli del test. Le griglie
sono fissate prima: la costante in $\{0, 5, 10, 20, 40, 80, 160, 320\}$, $\alpha$ da
0 a 1 a passi di 0,1. A parità di metrica vince il valore più vicino al predefinito,
per non spostare il comportamento di partenza senza un guadagno
(tabella~\ref{tab:calibrazione}).

\begin{table}[ht]
  \centering
  \small
  \setlength{\tabcolsep}{4pt}
  \begin{tabular}{lrrrrrrrr}
    \toprule
    costante & 0 & 5 & 10 & 20 & 40 & 80 & 160 & 320 \\
    \midrule
    SciFact train & 0,678 & 0,686 & 0,696 & 0,708 & 0,723 & 0,742 & \textbf{0,753} & 0,750 \\
    NFCorpus dev & 0,276 & 0,297 & 0,303 & 0,313 & 0,321 & 0,326 & \textbf{0,326} & 0,320 \\
    known-item train & 0,861 & 0,867 & \textbf{0,872} & 0,868 & 0,814 & 0,690 & 0,505 & 0,353 \\
    \midrule
    $\alpha$ & 0 & 0,2 & 0,4 & 0,5 & 0,6 & 0,8 & 0,9 & 1 \\
    \midrule
    SciFact train & 0,712 & 0,731 & 0,743 & \textbf{0,748} & 0,745 & 0,720 & 0,703 & 0,678 \\
    NFCorpus dev & 0,303 & 0,316 & 0,320 & \textbf{0,323} & 0,320 & 0,307 & 0,299 & 0,276 \\
    known-item train & 0,195 & 0,283 & 0,444 & 0,578 & 0,719 & 0,850 & \textbf{0,863} & 0,861 \\
    \bottomrule
  \end{tabular}
  \caption{La calibrazione della somma e della combinazione convessa: nDCG@10, MRR@10
  per le known-item. La griglia di $\alpha$ è completa nei dati, qui ridotta.}
  \label{tab:calibrazione}
\end{table}

\paragraph{Sul test.} I valori calibrati, misurati una volta sola sul test, danno la
tabella~\ref{tab:difetto3}.

\begin{table}[ht]
  \centering
  \small
  \begin{tabular}{lrrrrr}
    \toprule
    & \multicolumn{2}{c}{somma} & RRF & \multicolumn{2}{c}{convessa} \\
    \cmidrule(lr){2-3} \cmidrule(lr){5-6}
    & 20 & calibrata & & 0,5 & calibrata \\
    \midrule
    SciFact & 0,6913 & \textbf{0,7067} (160) & 0,6866 & 0,7025 & 0,7025 (0,5) \\
    NFCorpus & 0,3340 & 0,3441 (160) & \textbf{0,3447} & 0,3416 & 0,3416 (0,5) \\
    known-item & 0,8464 & \textbf{0,8499} (10) & 0,4623 & 0,5723 & 0,8466 (0,9) \\
    \bottomrule
  \end{tabular}
  \caption{Il difetto 3 sul test, con il recupero disgiuntivo e l'unione: nDCG@10,
  MRR@10 per le known-item. Fra parentesi il valore calibrato.}
  \label{tab:difetto3}
\end{table}

\textbf{Il 20 era sbagliato di un fattore otto per le query in lingua naturale e
giusto, per caso, per quelle identificative.} Calibrare la costante vale +0,015 su
SciFact e +0,010 su NFCorpus; la stessa costante, 160, sul train delle known-item porta
l'MRR da 0,87 a 0,50. RRF, che di parametri non ne ha, fa come la somma calibrata su
NFCorpus, sotto la somma a occhio su SciFact, e dimezza le known-item: pesa la
posizione vettoriale dell'atto giusto, pessima, quanto quella lessicale. La
normalizzazione della combinazione convessa non risolve: $\alpha = 0{,}5$ va bene per
le collezioni pubbliche e sulle known-item fa 0,57, ne serve 0,9 (figura~\ref{fig:fusione}).

\begin{figure}[ht]
  \centering
  \begin{tikzpicture}
    \begin{groupplot}[group style={group size=3 by 1, horizontal sep=1.3cm},
        tesi, width=.32\textwidth, height=5.4cm, xmin=-.3, xmax=7.3,
        xtick={0,...,7}, xticklabels={0,5,10,20,40,80,160,320},
        x tick label style={font=\tiny, rotate=90, anchor=east}, xlabel={costante}]
      \nextgroupplot[title={\footnotesize SciFact, nDCG@10}]
        \addplot[color=black, mark=*, mark size=1.5pt] table[x=posizione, y=scifact] {dati/fusione.dat};
      \nextgroupplot[title={\footnotesize NFCorpus, nDCG@10}]
        \addplot[color=black, mark=*, mark size=1.5pt] table[x=posizione, y=nfcorpus] {dati/fusione.dat};
      \nextgroupplot[title={\footnotesize known-item, MRR@10}]
        \addplot[color=black, mark=*, mark size=1.5pt] table[x=posizione, y=known] {dati/fusione.dat};
    \end{groupplot}
  \end{tikzpicture}
  \caption[La costante della somma pesata sulle collezioni di calibrazione]{La costante della somma pesata al variare del suo valore, sulle
  tre collezioni di calibrazione (SciFact e known-item, sezione di addestramento; NFCorpus, sviluppo). Il valore di
  partenza di Koskidex è 20: vicino all'ottimo sulle known-item, e otto volte più
  basso dell'ottimo sulle query in lingua naturale. Dati: \texttt{latex/dati/fusione.dat}.}
  \label{fig:fusione}
\end{figure}


Avevo previsto il contrario sulla parte che conta. La mia ipotesi era che la
costante dipendesse dalla lunghezza delle query, più alta su SciFact (dodici termini
in mediana) che su NFCorpus (due), e che solo la normalizzazione rendesse il peso
trasferibile fra le due. La prima è caduta: su SciFact e NFCorpus la costante
migliore è la stessa, 160. La seconda ha tenuto ($\alpha = 0{,}5$ su tutte e due),
ma senza il merito che le davo, perché anche le costanti coincidono. La scala che conta non è quella
fra query lunghe e corte ma quella fra \emph{tipi} di query. Su una query
identificativa il vettore non sa niente del numero e va tenuto basso; su una in
lingua naturale va pesato otto volte più di quanto facesse il motore. Nessuna
fusione con un parametro unico regge i due tipi, e la correzione del difetto 3 non
è una formula ma una scelta: il peso del vettore per tipo di query, che è
l'argomento della sezione~\ref{sec:scelta}.

\section{Il costo: scansione, latenza, memoria}\label{sec:costo}

\paragraph{La scansione.} Koskidex non ha un indice approssimato dei vettori: i più
vicini si trovano calcolando la similarità con tutti i documenti, in $O(n \cdot d)$.
È una scelta deliberata, per un archivio da decine di migliaia di documenti. Nelle
misure del difetto 2 l'unione costa, rispetto al re-ranking, 6,5~ms di mediana in
più su SciFact, 5,1 su NFCorpus e 15,3 sugli atti: una scansione, lineare nel
numero di documenti. Le fusioni nuove costano meno della somma, da 1 a 7~ms, perché
scandiscono i vettori una volta sola invece di calcolare prima la similarità dei
candidati lessicali e poi quella di tutti.

\paragraph{Un costo che le valutazioni non vedevano.} I vettori che arrivano al
motore via HTTP, o dall'istantanea su disco, sono liste generiche
(in Go \texttt{[]interface\{\}}), e il motore li ricopiava in un vettore di numeri
\emph{per ogni documento a ogni ricerca}, due volte con l'unione. Nelle valutazioni
piatte i vettori arrivano già come numeri, e il costo non si vedeva. Un benchmark
con 10.000 documenti da 1.024 dimensioni, dati come li dà l'API, lo isola
(tabella~\ref{tab:costo}).

\begin{table}[ht]
  \centering
  \begin{tabular}{lrrr}
    \toprule
    & prima & dopo & \\
    \midrule
    re-ranking, ms per ricerca & 57,6 & 24,7 & 2,33$\times$ \\
    unione, ms per ricerca & 102,6 & 42,7 & 2,40$\times$ \\
    unione, MB allocati per ricerca & 177,3 & 13,5 & $-92\%$ \\
    memoria viva dopo il caricamento, MB & 270,5 & 94,7 & $-65\%$ \\
    \bottomrule
  \end{tabular}
  \caption{Convertire i vettori una volta sola, all'aggiunta del documento. Mediane
  di cinque esecuzioni, Apple M2.}
  \label{tab:costo}
\end{table}

Convertirli una volta sola, all'aggiunta, dimezza la ricerca e toglie due terzi
della memoria, con punteggi identici al bit. Del tempo che resta, il 69\% è il
calcolo della similarità. Calcolare le norme dei documenti una volta sola, invece
che a ogni ricerca, sembrava il passo ovvio e non ha dato quasi niente: $-7\%$ con
l'unione, niente col re-ranking. Una spiegazione che regge i numeri, ma che non ho
misurato direttamente, è che il ciclo non è limitato dal numero di operazioni ma
dalla latenza di una catena di 1.024 somme dipendenti: le tre somme di prima correvano in parallelo, e toglierne due
lascia la stessa catena. Accorciarla vorrebbe più accumulatori, cioè sommare in un
altro ordine, e punteggi diversi nelle ultime cifre. A 10.000 documenti da 1.024
dimensioni una scansione esatta costa quindi 12-15~ms su un portatile:
l'unione, che ne fa due, costa 39,6~ms contro i 24,6 del re-ranking, che ne fa una (dopo aver precalcolato le norme).

\paragraph{In Documentale.} Dal vivo, sulle 24 query del confronto, il re-ranking
vettoriale sembrava portare la mediana a 33~ms, contro gli 8-11 senza vettori. Rimisurato in A/B, con il
motore di prima e quello di dopo sugli stessi dati, i due binari non differiscono
in modo misurabile (le passate ballano di 2~ms nei due versi), e con i vettori
delle query già calcolati la mediana è di 9,4-9,9~ms. I 23~ms in più erano il vettore della query, che Ollama
calcola la prima volta che una query si vede. Con il re-ranking e il recupero
congiuntivo i candidati lessicali sono pochi, e la conversione dei vettori non
pesava. Il costo del vettore, in un documentale, è la query, non la ricerca.

\section{Scegliere per tipo di query}\label{sec:scelta}

\paragraph{Senza vettori.} Già il capitolo~\ref{chapter:lessicale} aveva trovato che
nessuna impostazione lessicale unica va bene per le query identificative e per
quelle in lingua naturale: il recupero congiuntivo è giusto per
\texttt{1209 Crispiano} e disastroso per le frasi di SciFact. Scegliendo per ogni
query fra il recupero disgiuntivo e il congiuntivo, una regressione logistica su
otto caratteristiche della query (lunghezza, cifre, IDF dei termini, se il
congiuntivo torna vuoto) arriva a 0,0006 dall'oracolo che sceglie a posteriori.
Ma l'oracolo stesso non va oltre la scelta per collezione, e una regola scritta a
mano, \emph{recupero congiuntivo se la query ha una cifra e il congiuntivo trova
qualcosa}, fa esattamente quanto il classificatore.

\paragraph{Con l'unione e il peso del vettore.} I due difetti di questo capitolo
cambiano la domanda. Con l'unione il recupero congiuntivo non torna più vuoto, e il
peso del vettore giusto dipende dal tipo di query. Le due configurazioni fra cui
scegliere diventano: \textbf{A}, recupero disgiuntivo, unione, peso 160, per le
frasi; \textbf{B}, recupero congiuntivo, unione, peso 10, per gli identificativi.
I pesi vengono dalla calibrazione della sezione~\ref{sec:difetto3}, senza
ricalibrare. Le strategie, il classificatore e i suoi parametri sono gli stessi di
prima; la caratteristica \emph{il congiuntivo torna vuoto} si calcola sul solo lessicale,
perché B con l'unione non è mai vuota (tabella~\ref{tab:scelta}).

\begin{table}[ht]
  \centering
  \small
  \begin{tabular}{lrrrr}
    \toprule
    strategia & known-item & SciFact & NFCorpus & media \\
    \midrule
    sempre A (disgiuntivo, 160) & 0,5589 & \textbf{0,7067} & \textbf{0,3441} & 0,5366 \\
    sempre B (congiuntivo, 10) & \textbf{0,9680} & 0,6483 & 0,3325 & 0,6496 \\
    sempre congiuntivo, peso 20 & 0,9675 & 0,6483 & 0,3333 & 0,6497 \\
    oracolo & 0,9731 & 0,7280 & 0,3655 & 0,6889 \\
    \textbf{la regola} & \textbf{0,9680} & 0,7067 & 0,3438 & \textbf{0,6728} \\
    classificatore & 0,9680 & 0,7036 & 0,3432 & 0,6716 \\
    \midrule
    classificatore senza vettori & 0,9687 & 0,6694 & 0,3038 & 0,6473 \\
    \bottomrule
  \end{tabular}
  \caption{Scegliere per ogni query fra due configurazioni ibride: nDCG@10 sul
  test. La regola sceglie B se la query ha una cifra e il recupero lessicale
  congiuntivo trova qualcosa, altrimenti A.}
  \label{tab:scelta}
\end{table}

\textbf{Scegliere il peso del vettore per tipo di query vale +0,023 di \mbox{nDCG@10} medio
sulla migliore configurazione fissa}, e +0,026 sulla scelta senza vettori. Per
prenderlo basta ancora la regola di una riga; il classificatore impara la stessa
regola e fa appena peggio.

E adesso l'oracolo ha spazio oltre il tipo di query: 0,016 sopra la regola, quasi
tutto sulle frasi (+0,021 su SciFact e +0,022 su NFCorpus), dove senza vettori non ne
aveva. A e B sono ormai due ordinamenti diversi della stessa frase, il lessicale con
un vettore forte contro quasi solo il vettore, e B vince su 29 query di SciFact su
300 e su 82 di NFCorpus su 323. Nessuna caratteristica della query lo sa
riconoscere: per saperlo bisognerebbe guardare i risultati, non la query. È un
limite, e va scritto come tale.

\paragraph{Il rischio che resta aperto.} Qui ogni tipo di query sta in una collezione
diversa: le known-item hanno tutte una cifra e un recupero congiuntivo che trova
l'atto, mentre le frasi di SciFact, anche quando hanno una cifra (114 query su
300), col congiuntivo tornano quasi sempre vuote. Una regola che riconosce la collezione e una che riconosce il tipo
di query danno gli stessi numeri, e su questi dati non si possono distinguere. La
prova giusta sono query dei due tipi nella stessa collezione, cioè le known-item
scritte da persone sul corpus degli albi, dove \emph{ordinanza chiusura via Roma
luglio} non ha numeri e una frase può averne.

\paragraph{La prova sulle query umane.} Le 104 known-item umane
(\texttt{2026-09-25\_scelta-umane}), con il classificatore congelato, cioè con i
pesi imparati sopra e nessuna query umana usata per imparare, danno una prova
più debole di quella che serviva: solo 4 query su 104 contengono una cifra, e i
due tipi nella stessa collezione ci sono appena. Su queste query A e B danno lo
stesso nDCG@10 in 73 casi su 104, e la scelta vale poco in tutte e due le
direzioni. La regola sceglie B una volta sola, e fa 0,6727 contro 0,6719 di
sempre A; il classificatore sceglie B su 59 query e fa 0,6769, appena 0,004
sopra la regola; l'oracolo arriva a 0,7045, 0,032 sopra la regola. Tre
previsioni su quattro: è caduta quella che metteva il margine dell'oracolo fra
le query senza cifra, dove B batte A su 12 query su 100, non su almeno il
15\%. Il rischio della collezione resta quindi aperto, e la scelta per tipo di
query resta fuori dal motore e fuori dal profilo consigliato del
capitolo~\ref{chapter:produzione}.

\chapter{Il confronto finale con Elasticsearch}\label{chapter:confronto}

\section{Le configurazioni a confronto}\label{sec:configurazioni}

\todo[inline]{Bozza del 28/09, prima del terzo lotto: struttura e ragionamento
fissati, i numeri (\dacompletare) arrivano da
\texttt{2026-09-25\_known-item-umane} a raccolta chiusa.}

Il confronto finale si regge sulle known-item umane
(sezione~\ref{sec:query}): \dacompletare{} query scritte da persone che non
conoscono i motori, una per atto, su \dacompletare{} atti estratti con un seme,
metà per fonte. Per ogni query c'è un atto giusto solo, quindi la misura
principale è l'MRR@10, accanto a tre quantità che un utente sente di più:
l'atto al primo posto, l'atto fra i primi dieci, la ricerca a vuoto. I bisogni
aperti delle 24 query restano un confronto secondario: i giudizi del modello
concordano con quelli umani con un kappa pesato di 0,370
(sezione~\ref{sec:giudizi}), troppo poco per reggere un capitolo, e il resto
del pool non ha un giudizio umano.

Le configurazioni sono di due tipi, e si misurano in due modi.

\paragraph{Dall'app.} Due, sulle stesse query passate da
\texttt{app:eval-run-queries}, cioè dalla ricerca vera di Documentale:
\begin{itemize}
  \item \textbf{ES}: Elasticsearch come lo usa Documentale in produzione, con il
        \texttt{multi\_match} \texttt{best\_fields} e l'operatore
        \texttt{and} (capitolo~\ref{chapter:sistemi});
  \item \textbf{KC}: Koskidex innestato con il profilo consigliato
        (sezione~\ref{sec:cosa-cambia}): parole libere di stare in campi
        diversi, numeri esatti, elisione; senza vettori.
\end{itemize}

\paragraph{Koskidex piatto.} Sei, con \texttt{scripts/evaluate} sul corpus a
schede che esce dall'esportazione di Documentale, per vedere dove porterebbero
le correzioni dei capitoli~\ref{chapter:lessicale} e~\ref{chapter:ibrido} se
entrassero nell'app: \textbf{K0}, il punto di partenza (congiuntivo,
punteggio euristico); \textbf{LA} e \textbf{LT}, BM25 con la frequenza
mescolata in recupero disgiuntivo e congiuntivo; \textbf{A} e \textbf{B},
le due configurazioni ibride del capitolo~\ref{chapter:ibrido}, quella per le
frasi e quella per gli identificativi; \textbf{V}, i soli vettori. Le
configurazioni piatte non passano dall'app e non hanno i refusi, quindi fra
loro si confrontano alla pari, con ES e KC solo nell'ordine di grandezza.

\paragraph{Il tetto del riordino.} Un'ultima coppia non è una proposta:
\textbf{RR-LA} e \textbf{RR-A} sono i primi cento di LA e di A riordinati da un
cross-encoder multilingue (\texttt{bge-reranker-v2-m3}), che legge query e
documento insieme e costa secondi a query su questo hardware
(\texttt{2026-09-28\_reranker}). Dicono quanto vale, al massimo, riordinare
quello che il lessicale trova.

Ogni metrica si riporta su tutte le query, per fonte (Crispiano, dove chi
scriveva aveva letto il testo intero e l'indice ha la scheda, e Friuli Venezia
Giulia, dove l'atto è la sola scheda) e per query con e senza un numero. Le
sette previsioni sono state scritte il 25 settembre, prima che arrivasse una
sola risposta.

\section{Pertinenza, per famiglia di query e per fonte}\label{sec:pertinenza}

\begin{table}[ht]
  \centering
  \small
  \begin{tabular}{lrrrrr}
    \toprule
    & tutte & Crispiano & Friuli V.G. & con numero & senza numero \\
    \midrule
    ES & \dacompletare & \dacompletare & \dacompletare & \dacompletare & \dacompletare \\
    KC & \dacompletare & \dacompletare & \dacompletare & \dacompletare & \dacompletare \\
    \midrule
    K0 & \dacompletare & \dacompletare & \dacompletare & \dacompletare & \dacompletare \\
    LA & \dacompletare & \dacompletare & \dacompletare & \dacompletare & \dacompletare \\
    LT & \dacompletare & \dacompletare & \dacompletare & \dacompletare & \dacompletare \\
    A & \dacompletare & \dacompletare & \dacompletare & \dacompletare & \dacompletare \\
    B & \dacompletare & \dacompletare & \dacompletare & \dacompletare & \dacompletare \\
    V & \dacompletare & \dacompletare & \dacompletare & \dacompletare & \dacompletare \\
    \midrule
    RR-LA & \dacompletare & \dacompletare & \dacompletare & \dacompletare & \dacompletare \\
    RR-A & \dacompletare & \dacompletare & \dacompletare & \dacompletare & \dacompletare \\
    \bottomrule
  \end{tabular}
  \caption{Known-item umane: MRR@10 e, fra parentesi, la quota di ricerche a
  vuoto. ES e KC dall'app, le altre con Koskidex piatto sul corpus a schede.}
  \label{tab:known-item-umane}
\end{table}

\paragraph{Che cosa scrive chi cerca.} \dacompletare{} query su
\dacompletare{} contengono un numero; avevo previsto meno di una su quattro,
perché chi ricorda un atto dopo settimane ne ricorda l'argomento, non il
protocollo. \dacompletare{} persone su \dacompletare{} hanno risposto «non
saprei» per almeno un atto.
\todo[inline]{Previsione 1. Se le query con un numero sono poche, le colonne
«con numero» vanno lette come aneddoti: dirlo con il loro conteggio.}

\paragraph{Il campo unico anche sulle ricerche per contenuto.} Il difetto della
sezione~\ref{sec:numero-comune} era stato trovato sulle ricerche per numero e
comune. Qui la domanda è se colpisca anche chi cerca per argomento: con
l'operatore \texttt{and} di produzione tutte le parole devono stare nello
stesso campo della scheda, e chi scrive aggiunge facilmente il comune, un anno o
una parola del testo. ES va a vuoto nel \dacompletare{} delle query e trova
l'atto fra i primi dieci nel \dacompletare; KC, con le parole libere di stare
in campi diversi e l'elisione, va a vuoto nel \dacompletare{} e fa
\dacompletare{} di MRR@10 sopra ES.
\todo[inline]{Previsioni 2 e 3: più del 30\% a vuoto per ES, KC almeno 0,10
sopra con almeno 10 punti di vuoti in meno. Se tengono tutte e due, il
difetto del campo unico riguarda anche le ricerche per contenuto: è la frase
da scrivere.}

\paragraph{Congiuntivo o disgiuntivo.} Una parola in più, o scritta
diversamente da come è nell'atto, svuota il recupero congiuntivo. Sulle query
senza numero LA fa \dacompletare{} e LT \dacompletare; su quelle con un numero
\dacompletare.
\todo[inline]{Previsione 4: LA almeno 0,05 sopra LT senza numero, LT sopra LA
con il numero. Collegare alla sezione~\ref{sec:scelta}: la regola per tipo di
query e il classificatore sulle stesse query (\texttt{scelta-umane}).}

\paragraph{I vettori.} A, che aggiunge ai candidati di LA i cento più vicini,
fa \dacompletare{} sulle query senza numero, \dacompletare{} sopra LA; i soli
vettori (V) \dacompletare, e sulle query con un numero \dacompletare.
\todo[inline]{Previsioni 5 e 7: A almeno 0,03 sopra LA senza numero; V entro
0,10 da LA senza numero e sotto 0,30 con il numero. È la verifica, su query
vere, di quello che il capitolo~\ref{chapter:ibrido} aveva visto sulle
collezioni pubbliche.}

\paragraph{Le due fonti.} Su Crispiano chi scriveva aveva letto il testo
intero, e l'indice ha la scheda: le parole che ricorda possono non esserci. LA
fa \dacompletare{} su Crispiano e \dacompletare{} sul Friuli Venezia Giulia, KC
\dacompletare{} e \dacompletare.
\todo[inline]{Previsione 6: almeno 0,10 sotto su Crispiano per LA e KC.
Rimandare alla sezione~\ref{sec:scheda-testo} per scheda contro testo intero
sui soli atti di Crispiano, e alla sezione~\ref{sec:italiana} per stopword e
stemmer italiani sulle stesse query.}

\paragraph{Quanto vale riordinare.} Il cross-encoder porta LA da
\dacompletare{} a \dacompletare{} di MRR@10 e A da \dacompletare{} a
\dacompletare; degli atti che il primo stadio aveva fra l'undicesimo e il
centesimo posto ne porta fra i primi dieci \dacompletare{} su \dacompletare.
Sulle collezioni pubbliche il guadagno è più piccolo e ha un tetto: SciFact
passa da 0,676 a 0,723 di nDCG@10, NFCorpus da 0,306 a 0,331, e su NFCorpus
il primo stadio porta fra i cento soltanto il 24\% dei documenti giusti, gli
unici che il riordino può vedere. Sulle ricerche \texttt{<numero> <comune>}
porta LA da 0,833 a 0,953, contro 0,960 del recupero congiuntivo: degli atti
fra l'undicesimo e il centesimo posto ne porta fra i primi dieci 12 su 13.
Avevo previsto che restasse sotto LT perché un cross-encoder non pesa un
numero esatto più di uno simile; resta sotto, ma di 0,007, e il ragionamento
era sbagliato. Il prezzo è di 5 secondi di mediana a query sulle schede e di
12 sugli abstract di SciFact, contro il millisecondo di Koskidex: la misura di
un tetto, non una configurazione da mettere in produzione.
\todo[inline]{Previsioni umane di \texttt{reranker}: almeno +0,05 su LA e
+0,03 su A, almeno metà degli atti dall'11° al 100° posto portati fra i primi
dieci.}

\section{Latenza, memoria, dimensione dell'indice}\label{sec:latenza}

La sezione~\ref{sec:prestazioni-innesto} ha dato i tempi dei due motori
passando dall'app, con Koskidex nativo ed Elasticsearch nella macchina
virtuale: un ordine di grandezza, non un confronto alla pari. Qui i due motori
girano nella stessa macchina virtuale, con le stesse 4 CPU e la stessa rete:
Koskidex è lo stesso binario compilato per Linux, in un container accanto a
quello di Elasticsearch, con l'indice riempito dall'app nella configurazione
consigliata. Le richieste sono quelle che manda l'app, copiate dal codice: per
Elasticsearch il \texttt{multi\_match} di produzione con i sei campi pesati,
che restituisce anche il testo di ogni risultato; per Koskidex la ricerca con
i soli identificativi. Le query sono 400: le 300 known-item automatiche, le 24
del confronto e le 76 known-item umane raccolte fino al 28 settembre. Tutto
sta in \texttt{risultati/esperimenti/2026-09-28\_carico/}, previsioni
comprese.

Koskidex tiene in memoria le ultime 1.024 risposte, e non c'è un'impostazione
per spegnere la cache: ripetere le stesse query misurerebbe la cache. Ogni
richiesta porta quindi un numero diverso di spazi in coda alla query, che il
tokenizer ignora e la chiave della cache no; lo strumento verifica, prima di
misurare, che con gli spazi nessuna delle 400 risposte cambi, su nessuno dei
due motori.

\paragraph{Una ricerca alla volta.} Ogni query cinque volte, in un ordine
mescolato con un seme fisso. Il p50 di Koskidex è un terzo di quello di
Elasticsearch, e l'avevo previsto; avevo previsto anche che la differenza
stesse soprattutto nel peso della risposta, perché Elasticsearch restituisce il
testo degli atti che l'app poi scarta. Non è così: senza il testo Elasticsearch
non va più veloce (6,77~ms), perché con l'operatore \texttt{and} le risposte
sono piccole, 160 byte di mediana. La differenza sta nel motore. La coda invece
era quasi la stessa (tabella~\ref{tab:carico-sequenziale}, colonna
\emph{prima}): 68 query su 400 superavano gli 8~ms, ed erano le query lunghe,
piene di parole comuni.

\begin{table}[ht]
  \centering
  \begin{tabular}{lrrr}
    \toprule
    & p50 & p95 & p99 \\
    \midrule
    Elasticsearch, come l'app & 6,32 & 17,66 & 24,26 \\
    Elasticsearch, senza il testo & 6,77 & 19,60 & 26,09 \\
    Koskidex, prima & 2,12 & 12,36 & 18,31 \\
    Koskidex, dopo & \textbf{1,12} & \textbf{2,74} & \textbf{4,79} \\
    \bottomrule
  \end{tabular}
  \caption{Una ricerca alla volta, in millisecondi, 400 query per cinque
  passate, i due motori nella stessa macchina virtuale. \emph{Prima} è
  Koskidex \texttt{ffa38ab}, \emph{dopo} \texttt{510b9d2}, con le due correzioni
  di questa sezione.}
  \label{tab:carico-sequenziale}
\end{table}

\paragraph{Sotto carico.} Da 1 a 32 client che mandano ricerche senza pause,
20 secondi per livello. Avevo previsto che Koskidex reggesse almeno una volta e
mezza le ricerche al secondo di Elasticsearch. Ne reggeva 1,14 volte: 653
contro 571 (tabella~\ref{tab:carico}). Con un client solo usava già il 116\% di
CPU, più di un core per una ricerca alla volta, e a 32 client il suo p99 era il
doppio di quello di Elasticsearch.

\begin{table}[ht]
  \centering
  \begin{tabular}{rrrrrrr}
    \toprule
    & \multicolumn{2}{c}{Elasticsearch} & \multicolumn{2}{c}{Koskidex, prima} &
    \multicolumn{2}{c}{Koskidex, dopo} \\
    \cmidrule(lr){2-3} \cmidrule(lr){4-5} \cmidrule(lr){6-7}
    client & ricerche/s & p99 & ricerche/s & p99 & ricerche/s & p99 \\
    \midrule
    1 & 134,6 & 22,8 & 233,3 & 19,2 & 705,9 & 4,9 \\
    4 & 471,7 & 26,6 & 588,9 & 29,1 & 1.707,9 & 7,9 \\
    8 & 571,1 & 46,5 & 653,0 & 48,5 & 2.313,8 & 13,9 \\
    16 & 550,4 & 69,6 & 622,6 & 105,0 & 2.588,7 & 25,3 \\
    32 & 559,4 & 103,1 & 647,2 & 214,6 & \textbf{2.614,0} & \textbf{42,8} \\
    \bottomrule
  \end{tabular}
  \caption{Sotto carico, stesse 4 CPU: ricerche al secondo e p99 in
  millisecondi. Elasticsearch e Koskidex prima saturano le CPU dagli 8 client
  in su (380-390\%); Koskidex dopo arriva al 319\% a 8 client e al 360\% a
  32, e il resto va alla rete.}
  \label{tab:carico}
\end{table}

\paragraph{Dove spendeva Koskidex.} Il profilo a 8 client ha dato la risposta, e
non era quella che avevo previsto: avevo scommesso sull'espansione dei refusi e
sulla scrittura del JSON della risposta. Koskidex allocava 2,57~MB a ricerca, e
quella memoria si pagava in CPU: il 20\% dei campioni era nel lock
dell'allocatore di Go e un altro 7\% nel garbage collector. Il 43\% delle
allocazioni veniva dalla funzione che, per ogni parola della query, costruisce
una voce per ogni documento che la contiene, anche per \emph{di} o \emph{per},
e solo dopo scarta i documenti a cui manca una parola; un altro 40\% dalla
distanza fra parole, che allocava una matrice intera a ogni confronto, e dalla
raccolta dei candidati con refuso.

\paragraph{Due correzioni che non cambiano un risultato.} A differenza delle
correzioni dei capitoli precedenti non stanno dietro un'impostazione, e il
motivo è la condizione stessa del lavoro: non devono cambiare né l'insieme né l'ordine né un
punteggio, e se ne cambiassero uno non sarebbero queste correzioni. La prima
calcola la distanza con tre righe riusate al posto della matrice, la stessa
ricorrenza, e prende i bigrammi come sottostringhe. La seconda, con il recupero
congiuntivo, trova prima i documenti che contengono tutte le parole, partendo
da quella con meno occorrenze come fa Elasticsearch, e calcola il punteggio
solo per loro, parola per parola nell'ordine della query: le somme in virgola
mobile restano le stesse bit per bit. La prova è stata fatta prima di misurare
i tempi: l'impronta di ogni risposta del server, identificativi e punteggi
byte per byte, per 429 query (le 400 più 29 inventate con gli operatori, parole
di una lettera, numeri, refusi e parole inesistenti), coincide prima e dopo
per tutte, a 10.018 documenti e a 100.000; le metriche di
\texttt{scripts/evaluate}, query per query, coincidono su SciFact, NFCorpus e le
known-item automatiche nelle tre configurazioni di confronto; i test di
Koskidex passano, compreso quello che congela l'ordine di partenza.

Le otto previsioni scritte prima (\texttt{2026-09-28\_prestazioni}) sono tutte
confermate, quasi sempre con molto margine: le soglie erano prudenti. Le
allocazioni scendono da 2,57 a 0,44~MB a ricerca, l'83\% in meno. Una ricerca
alla volta costa 1,12~ms di p50 e 2,74 di p95, e il p95 scende di più del p50:
le correzioni hanno tolto proprio il costo delle query lunghe, che faceva la
coda. Sotto carico Koskidex regge 2.614 ricerche al secondo, 4,6 volte
Elasticsearch sulle stesse CPU, e a 32 client il suo p99 è 43~ms contro 103.
Con un client la CPU scende dal 116\% all'83\%. Nel profilo nuovo il lock
dell'allocatore e il garbage collector escono dalle prime righe, e il 43\% dei
campioni è nelle chiamate di sistema della rete: il motore non è più la parte
lenta della richiesta. Nativo, sugli 8 core del Mac, Koskidex passa da 975 a
4.100 ricerche al secondo.

\paragraph{Una terza correzione, e due previsioni sbagliate.} Delle allocazioni
che restavano, l'86\% veniva da una riga sola: l'insieme che toglie i doppioni
fra le parole del vocabolario che condividono un bigramma con quella cercata,
migliaia per i bigrammi comuni, quasi tutte scartate subito dopo perché troppo
lunghe o troppo corte. La correzione scarta prima e toglie i doppioni solo fra
quelle che restano; il controllo dipende solo dalla parola, quindi le parole
trovate sono le stesse, nello stesso ordine
(\texttt{2026-09-28\_allocazioni}). La prova è la stessa: 429 impronte su 429
a tutte e due le dimensioni, 9 valutazioni su 9, più un test che confronta
42.420 ricerche di parole con il percorso di prima. Misurato contro un
\emph{prima} rifatto nella stessa mezz'ora, le allocazioni scendono da 0,435 a
0,081~MB a ricerca, il p50 di una ricerca alla volta nel container da 1,14 a
0,92~ms, e la capacità nel container da 2.622 a 3.779 ricerche al secondo,
+44\%, 6,6 volte Elasticsearch; nativo da 4.090 a 6.376. Cinque
previsioni su sette. Una soglia era una quota, meno del 10\% dei byte allocati
per quella funzione: è scesa al 30\%, ma perché è sceso tutto il resto con lei,
e in valore assoluto è passata da 385 a 25~KB a ricerca. L'altra si aspettava
un guadagno forte a 100.000 documenti, dove pensavo che le liste dei bigrammi
fossero dieci volte più lunghe; ma la copia ripete gli stessi testi, il
vocabolario non cresce, e la coda lì viene dalle liste dei documenti: p95 e p99
scendono solo del 4 e del 7\%.

\paragraph{Memoria e archivi più grandi.} Per vedere come cresce il costo ho
copiato gli stessi atti fino a 50.000 e 100.000 documenti, con identificativi
diversi e lo stesso testo: dati inventati, che per tempi e memoria vanno bene e
per la pertinenza no (tabella~\ref{tab:crescita}).

\begin{table}[ht]
  \centering
  \small
  \begin{tabular}{lrrrrrr}
    \toprule
    & \multicolumn{2}{c}{memoria} & \multicolumn{2}{c}{indice su disco} &
    \multicolumn{2}{c}{indicizzazione} \\
    \cmidrule(lr){2-3} \cmidrule(lr){4-5} \cmidrule(lr){6-7}
    documenti & ES & Koskidex & ES & Koskidex & ES & Koskidex \\
    \midrule
    10.018 & 1.431 MB & 155 MB & 3,5 MB & 6,3 MB & 1,1-1,7 s & 1,9-2,1 s \\
    50.000 & 1.439 MB & 651 MB & 17,9 MB & 32,8 MB & 2,8 s & 6,2 s \\
    100.000 & 1.433 MB & 1.213 MB & 35,5 MB & 65,6 MB & 5,3 s & 12,9 s \\
    \bottomrule
  \end{tabular}
  \caption{Memoria a riposo dopo un riavvio, indice su disco e tempo di
  indicizzazione, con gli atti copiati. A 10.018 documenti l'indicizzazione è
  quella dall'app.}
  \label{tab:crescita}
\end{table}

Alla dimensione dell'albo Koskidex occupa un decimo della memoria di
Elasticsearch, e sotto carico 272~MB al massimo contro 1.566, 195 dopo la terza
correzione. Ma tiene tutto in
memoria, con gli identificativi dei documenti e i nomi dei campi come stringhe
in ogni occorrenza: circa 12~KB a documento, e a 100.000 documenti la memoria è
quasi quella di Elasticsearch, che ha un heap fisso. Le correzioni non la
toccano (1.213 contro 1.184~MB), perché non toccano le strutture dell'indice.
Toccano invece la coda: a 100.000 documenti il p95 di Koskidex era di 133~ms e
il p99 di 201, dieci volte quelli a 10.018, contro i 21 e 36 di Elasticsearch;
dopo sono 10,4 e 15,5, sotto quelli di Elasticsearch, e 9,9 e 13,6 dopo la
terza.

\paragraph{Un difetto trovato per strada.} Il test che confronta i due percorsi
su query casuali ha trovato una cosa che non cercava, e che ha avuto il suo
esperimento (\texttt{2026-09-28\_ordine-fisso}). Per una parola corta con
un refuso ammesso, quattro lettere o meno con la tolleranza automatica,
Koskidex non usa i bigrammi ma scorre tutto il vocabolario, e il vocabolario è
una mappa di Go, che si percorre in un ordine diverso a ogni chiamata. L'ordine
delle parole evidenziate cambia da una ricerca all'altra; e se un documento
contiene due parole alla stessa distanza da quella cercata, quale delle due gli
viene accreditata, e quindi il suo punteggio BM25, dipende dal caso: in un test
con due documenti che contengono le stesse due parole con frequenze diverse,
i due documenti si scambiano di posto da una chiamata all'altra. Lo stesso
succede con la ricerca per sottostringa, che scorre il vocabolario allo stesso
modo, e dopo un riavvio anche altrove, perché l'istantanea su disco scrive i
documenti in ordine di mappa. La correzione è un'impostazione,
\texttt{StableTermOrder}, spenta per default, che visita i termini trovati in
ordine lessicografico. Sei previsioni su sette hanno tenuto: nessun numero di
\texttt{scripts/evaluate} dipendeva dall'ordine della mappa, perché lo
strumento cerca senza refusi, e cinque ripetizioni non variano mai, spente o
accese; Documentale, con il punteggio euristico, dà le stesse 429 risposte;
il costo sta nel rumore. È caduta quella che dava per uguali le medie con
l'impostazione accesa e spenta: su SciFact con BM25 disgiuntivo passano da
0,6694 a 0,6670. Anche l'ordine di ingresso nell'indice, che le misure della
tesi hanno sempre usato, è arbitrario, e decide a parità di distanza quale
espansione per prefisso viene accreditata; la differenza, due query su 300, è
quanto pesa quella regola (sezione~\ref{sec:espansioni}).

\paragraph{Cosa ne segue per Documentale.} All'albo di un comune, diecimila
atti, Koskidex costa un decimo della memoria di Elasticsearch e risponde in
meno di un quinto del tempo, e sotto carico regge più di sei volte le ricerche
sulle stesse CPU. A centomila documenti il tempo regge ancora, la memoria no: lì
serve un indice più compatto, con identificativi e campi interi, che
ridurrebbe anche l'indice su disco, oggi il doppio di quello di
Elasticsearch, e l'indicizzazione, 2,4 volte più lenta. I limiti delle
misure: il client gira sullo stesso Mac, fuori dalla macchina virtuale, e tiene
le connessioni aperte, mentre l'app in PHP ne apre in genere una per ricerca;
Elasticsearch ha l'heap fissato a 1~GB, e con meno memoria si comporterebbe
diversamente.

\section{Dove un motore piccolo è competitivo, e dove no}\label{sec:competitivo}

\todo[inline]{Bozza del 28/09: le affermazioni sui costi hanno già i loro
numeri (sezione~\ref{sec:latenza}); quelle sulla pertinenza aspettano la
sezione~\ref{sec:pertinenza}.}

La domanda della tesi non era se Koskidex batte Elasticsearch, ma se un motore
piccolo, innestato in un documentale al posto di uno grande, regge, e dove
smette di reggere. Le risposte sono di tre tipi.

\paragraph{Dove regge.} All'albo di un comune, diecimila atti, Koskidex
risponde in 0,92~ms di p50 contro 6,32, regge 6,6 volte le ricerche al secondo
sulle stesse quattro CPU e occupa un decimo della memoria. Sulle ricerche per
numero e comune il profilo consigliato porta l'atto giusto al primo posto nel
91,7\% dei casi, contro l'1,3\% di produzione, e le date si trovano in
qualunque formato. Sulle ricerche per contenuto scritte da persone KC fa
\dacompletare{} di MRR@10 contro \dacompletare{} di ES.
\todo[inline]{Se la previsione 3 cade, questa frase cambia verso: dire allora
che la differenza di pertinenza sta nelle impostazioni di produzione, non nel
motore, e che Elasticsearch con le stesse correzioni farebbe lo stesso
(la parità del capitolo~\ref{chapter:sistemi}).}

\paragraph{Dove il vantaggio non è del motore.} Quasi tutto il guadagno di
pertinenza viene da impostazioni che Elasticsearch ha già, o potrebbe avere:
parole in campi diversi con \texttt{cross\_fields}, l'elisione con il suo
filtro, le date con un \emph{char filter}; i numeri esatti non hanno
un'opzione, ma si ottengono separando nella query i termini numerici e
cercandoli senza fuzziness. Il capitolo~\ref{chapter:sistemi} ha
mostrato che con le stesse impostazioni i due motori trovano gli stessi
insiemi. Il merito di Koskidex, in questa tesi, è stato di rendere visibili e
misurabili i difetti di una configurazione, non di esserne immune.

\paragraph{Dove non regge.} La memoria: Koskidex tiene tutto in memoria, con
identificativi e nomi dei campi come stringhe in ogni occorrenza, e a centomila
documenti ne usa 1.213~MB, quasi quanto l'heap fisso di Elasticsearch. L'indice
su disco è il doppio e l'indicizzazione 2,4 volte più lenta. Mancano le cose che
un documentale grande dà per scontate e che questa tesi non ha misurato:
repliche, shard, un ecosistema di strumenti. E la pertinenza sulle ricerche
aperte resta non misurata con giudizi umani.
\todo[inline]{Chiudere con una frase sola: per chi, con quali numeri di
archivio e di utenti, un motore piccolo è la scelta giusta.}


%%%%%%%%%%%%%%%%%%%%%%%%%%%%%%%% Conclusioni
\pagestyle{plain}
\chapter{Conclusioni}\label{chapter:conclusione}

\section{Le risposte alla domanda}\label{sec:risposte}

La domanda (sezione~\ref{sec:domanda}) si articolava in quattro parti, e
tutte e quattro hanno una risposta misurata; per la pertinenza sulle ricerche
scritte da persone la risposta viene dalle known-item umane del
capitolo~\ref{chapter:confronto}.

\paragraph{I difetti di Koskidex.} Dei quattro difetti di Koskidex, quello che
viene prima di tutti non era fra i tre sospettati leggendo il codice. Il recupero congiuntivo
lasciava a vuoto 290 query di SciFact su 300: correggerlo porta il nDCG@10 da
0,0246 a 0,4936, e solo dopo ha senso parlare di punteggio. BM25 lo porta a
0,6197, e la correzione di un difetto dentro BM25, le espansioni pesate con la
frequenza del termine trovato, a 0,6694; con le stopword Koskidex arriva a
0,6757, il 99,5\% del riferimento di Lucene, e su NFCorpus al 95\%. Sui vettori, il
re-ranking vale due punti con il recupero disgiuntivo; con quello congiuntivo di
Documentale non vale niente, e l'unione dei candidati vale +0,62 su SciFact. La
costante della fusione era otto volte troppo bassa per le query in lingua naturale
e giusta, per caso, per quelle identificative.

\paragraph{I difetti di Documentale.} La ricerca di produzione ha quattro difetti
che colpiscono le ricerche più ordinarie: un numero d'atto trova gli atti con un
numero a un refuso, un articolo eliso rende introvabile una parola in quattro atti
su dieci, numero e comune insieme non trovano quasi mai l'atto, perché devono
stare nello stesso campo, e una data scritta in un formato diverso da quello
dell'atto non lo trova quasi mai. Con Koskidex innestato e tre impostazioni, raccolte in un
profilo (per le ricerche per numero e comune ne bastano due), l'atto giusto passa dall'1,3\% al 91,7\% di prime posizioni sulle
ricerche per numero e comune, e le parole introvabili dal 40,3\% allo 0,4\% degli
atti; con un'impostazione per le date, misurata anche dall'app, una data si
trova in qualunque formato. Gli importi, nell'app, si trovano già fra formati,
ma per un refuso ammesso che fa combaciare anche importi diversi: negare i
refusi agli importi, come ai numeri d'atto, e togliere i punti delle migliaia
dà l'importo giusto in qualunque formato, e in 102 query su 104 soltanto quello.
Nessuno dei quattro è un errore di programmazione: sono impostazioni
ragionevoli di un motore generico, che si rompono su proprietà precise dei
documenti amministrativi italiani.

\paragraph{Una configurazione per tutto.} Non c'è. Il recupero congiuntivo è giusto
per le ricerche identificative e disastroso per le frasi; le vie di mezzo di
Lucene ed Elasticsearch non reggono tutte e due; il peso giusto del vettore
dipende dal tipo di query. Sceglierlo con una regola di una riga, \emph{recupero
congiuntivo e vettore leggero se la query ha una cifra e il congiuntivo trova
qualcosa}, vale +0,023 di nDCG@10 medio sulla migliore configurazione fissa. Se la
regola riconosca il tipo di query o la collezione, su questi dati non si può dire,
e le query umane non lo decidono: solo 4 su 104 hanno una cifra: su quelle quattro la regola fa 0,021 di
nDCG@10 più di sempre A, e su tutte le 104 solo 0,0008.

\paragraph{Il confronto con Elasticsearch.} Nella stessa macchina virtuale, con
le stesse CPU e le richieste dell'app, Koskidex risponde in 0,9~ms di p50
contro 6,3, regge circa sette volte le ricerche al secondo e occupa un decimo della
memoria. Non era così prima di misurarlo sotto carico: allocava 2,6~MB a
ricerca, e il vantaggio si fermava a 1,14 volte. Quattro correzioni che non cambiano
un solo risultato, verificato risposta per risposta, hanno tolto quel costo:
ora alloca 0,05~MB a ricerca.
Il limite è la memoria: a centomila documenti è quasi quella di
Elasticsearch. Sulle 104 ricerche scritte da persone Koskidex innestato fa
0,432 di MRR@10 contro 0,326 di Elasticsearch di produzione, e lascia a vuoto
tre ricerche su dieci invece di quasi una su due: il difetto del campo unico
colpisce anche chi cerca per argomento. Con le stesse correzioni anche
Elasticsearch arriva a 0,449: sulle ricerche per contenuto il vantaggio sta
nelle impostazioni, non nel motore. Il margine più grande sta nel recupero e nei
vettori. Dall'app, il recupero disgiuntivo porta le ricerche scritte da persone
a 0,502 ma fa crollare quelle per numero e comune a 0,147, e solo il recupero
congiuntivo con i vettori a peso 10 tiene insieme le due cose, con 0,607 e 0,956,
al prezzo di circa 42~ms a ricerca e di otto minuti per indicizzare i 10.018
atti, se i vettori non sono già in cache. Fuori dall'app, lo stesso motore in recupero disgiuntivo fa 0,559 e con i
vettori 0,613, e l'atto è fra i primi dieci nell'86\% delle ricerche. Il limite
che resta a Documentale è il recupero congiuntivo senza vettori che ha scelto,
non il motore.
Due risultati sono andati contro le attese: sugli atti di Crispiano il testo
intero ritrova l'atto peggio della scheda, e leggere più testo nei vettori
peggiora invece di migliorare.

\section{Limiti}\label{sec:limiti}

\begin{itemize}
  \item \textbf{Query fabbricate.} Le known-item automatiche hanno una sola forma,
        \texttt{<numero> <comune>}, e le 24 query del confronto le ha scritte
        l'autore. Misurano bene quello che misurano, e niente di più. Le query
        scritte da persone servono proprio a questo, ma sono 104, di tre
        persone: la raccolta si è chiusa con tre lotti su quattro, il minimo
        del protocollo. Solo 4 contengono un numero, troppo poche per dire
        qualcosa di quelle, e le differenze sotto qualche centesimo di MRR@10
        stanno nel rumore di una collezione così.
  \item \textbf{Un solo corpus di dominio.} Gli albi pretori hanno la forma di un
        documentale, ma sono gli atti di comuni: altri archivi, con contratti o
        corrispondenza, hanno altri vocabolari e altre ricerche. Il testo intero
        c'è solo per 563 atti su 10.018.
  \item \textbf{Tipi di query separati per collezione.} Le query identificative
        stanno tutte sugli albi, quelle in lingua naturale sulle collezioni
        pubbliche: una regola che riconosce la collezione e una che riconosce il
        tipo di query danno gli stessi numeri, e le query umane, quasi tutte
        senza cifre, non bastano a separarle. È il limite più serio della
        scelta per tipo di query, e la ragione per cui resta fuori dal motore.
  \item \textbf{I giudizi del pool.} Bisogni e giudizi del pool li ha dati un
        modello linguistico, e con il mio campione l'accordo è modesto (kappa
        pesato 0,370): il pool non regge da solo un confronto, e il capitolo
        finale si regge sulle known-item umane. Un kappa basso con un solo
        annotatore umano non dice chi dei due ha ragione, e i bisogni scritti
        dallo stesso modello che giudica lo rendono difficile da leggere.
  \item \textbf{Il richiamo del pool è un limite superiore.} I documenti pertinenti
        che nessuna delle dodici configurazioni ha messo fra i primi dieci non
        sono giudicati, e contano come non pertinenti.
  \item \textbf{Elasticsearch corretto a metà.} Il confronto con le stesse
        correzioni vale per le parole libere in campi diversi, le elisioni e i
        numeri senza refusi. Le date e gli importi non sono stati normalizzati
        in Elasticsearch, che chiede un \emph{char filter} per formato, e le
        clausole per parola sono una delle traduzioni possibili delle
        impostazioni di Koskidex.
  \item \textbf{I tempi.} Nei capitoli di recupero Koskidex gira nativo ed
        Elasticsearch in una macchina virtuale, e i tempi dicono l'ordine di
        grandezza. Il confronto alla pari della
        sezione~\ref{sec:latenza} ha i suoi limiti: il client sullo stesso Mac,
        con le connessioni aperte, l'heap di Elasticsearch fissato a 1~GB, gli
        archivi più grandi fatti di atti copiati.
  \item \textbf{Un solo modello di embedding.} \texttt{bge-m3} in locale è una
        scelta, non un confronto: con un altro modello i pesi calibrati del
        capitolo~\ref{chapter:ibrido} potrebbero cambiare.
\end{itemize}

\section{Sviluppi futuri}\label{sec:sviluppi}

\begin{itemize}
  \item \textbf{Il recupero congiuntivo di Documentale.} Sulle query scritte da
        persone è il limite che resta dopo il profilo consigliato, e il
        capitolo~\ref{chapter:confronto} lo ha misurato dentro l'app come
        profilo intero: il recupero disgiuntivo non si può adottare, perché
        fa crollare le ricerche per numero d'atto, e l'unione con i vettori a
        peso 10 dà 0,607 senza toccarle. Resta da decidere se il costo, un
        servizio Ollama e otto minuti di indicizzazione senza cache, sia accettabile per un
        documentale. Stopword e stemmer italiani, che sul congiuntivo tolgono
        7,7 punti di ricerche a vuoto, non sono ancora stati misurati dentro
        l'app.
  \item \textbf{Più query scritte da persone}, e più persone: con cento query
        un effetto di qualche centesimo non si distingue dal rumore, e le
        ricerche con un numero mancano quasi del tutto.
  \item \textbf{Un secondo stadio di riordino.} Il capitolo~\ref{chapter:confronto}
        misura fuori linea quanto vale riordinare i primi cento candidati con un
        \emph{cross-encoder}: il re-ranking funziona se è progettato come
        stadio, con pochi candidati e un modello preciso. Portarlo nell'app
        vuol dire un servizio a parte, non dentro Koskidex, che resterebbe
        senza dipendenze, e un costo da abbattere: 5-12 secondi a query sulla
        GPU di un portatile con cento candidati. Con venti sono circa 1,2
        secondi sulle schede e il guadagno è quasi lo stesso, ma resta troppo
        per una ricerca dal vivo: restano testi più corti, un modello più
        piccolo o una GPU di un server.
  \item \textbf{La scansione dei vettori.} A 10.000 documenti una scansione esatta
        costa 12-15~ms per passata; più accumulatori nel prodotto scalare la
        accorcerebbero, a prezzo di punteggi diversi nelle ultime cifre, e oltre
        qualche centinaio di migliaia di documenti servirebbe un indice
        approssimato.
  \item \textbf{Un indice più compatto.} Identificativi dei documenti e nomi dei
        campi come interi invece che come stringhe in ogni occorrenza: è la
        condizione per usare Koskidex oltre i centomila documenti, dove oggi la
        memoria è quasi quella di Elasticsearch, e ridurrebbe l'indice su disco
        e il tempo di indicizzazione. Cambia il formato su disco, e per questo è
        rimasto fuori dalle correzioni della sezione~\ref{sec:latenza}.
  \item \textbf{Koskidex in produzione.} La tesi misura il recupero; un motore in
        produzione chiede anche un'istantanea che non riscriva l'indice intero,
        l'aggiornamento parziale di un documento, e un indice che non stia tutto
        in memoria.
\end{itemize}


%%%%%%%%%%%%%%%%%%%%%%%%%%%%%%%% Appendici
\appendix
\chapter{Come rifare ogni numero}\label{chapter:appendice-riproducibilita}
\sloppy

Ogni numero di questa tesi viene da un file dell'archivio \texttt{risultati/}, nel
repository della tesi, e ogni file registra il commit del codice che l'ha prodotto
e se l'albero era pulito (sezione~\ref{sec:principi}). Questa appendice dice come
si rifà un file partendo da zero.

\section*{I tre repository}

\begin{itemize}
  \item \textbf{Koskidex}: il motore, con \texttt{scripts/evaluate},
        \texttt{scripts/compare} e \texttt{scripts/pool}, e il diario delle
        modifiche al ranking in \texttt{eval/DIARIO.md}. Richiede Go e nessuna
        dipendenza oltre \texttt{golang.org/x/text}.
  \item \textbf{Documentale}, ramo \texttt{martin/tesi-magistrale}: i comandi di
        importazione, esportazione e valutazione aggiunti per la tesi, e
        l'innesto di Koskidex scelto da \texttt{SEARCH\_BACKEND}.
  \item \textbf{Il repository della tesi}: l'archivio \texttt{risultati/}, un
        README e uno script di analisi per ogni esperimento, gli strumenti di
        raccolta e di annotazione, questo testo.
\end{itemize}

\section*{L'ambiente}

La macchina delle misure è descritta in \texttt{risultati/macchina.md}: Apple M2
(salvo la prova di carico e il riordino su scheda grafica, misurati su un secondo
computer, capitolo~\ref{chapter:confronto}), Go 1.27.1, PHP 8.5.5. MySQL 8.4 ed Elasticsearch 9.1.0 girano in Docker,
Elasticsearch a nodo singolo con 1 GB di heap e la sicurezza spenta; gli embedding
vengono da Ollama con il modello \texttt{bge-m3}, impronta \texttt{790764642607}.
Gli strumenti scrivono nell'archivio solo se la variabile \texttt{TESI\_RISULTATI}
indica la cartella \texttt{risultati/}; se manca, lo dicono.

\section*{I dati}

Nessun dato sta nei repository: si riscaricano.
\begin{enumerate}
  \item \textbf{SciFact e NFCorpus}: \texttt{eval/corpora/fetch.sh} di Koskidex
        scarica gli archivi di BEIR e ne controlla l'impronta MD5.
  \item \textbf{Gli albi pretori}: \texttt{eval/corpora/fetch-albo.py} scarica il
        feed di Crispiano con i PDF e l'API del Friuli Venezia Giulia, ed estrae
        il testo dei PDF con \texttt{pdftotext}. Le fonti si aggiornano, e un
        nuovo scaricamento può dare atti in più: il corpus usato nelle misure ha
        la sua impronta SHA-256 registrata nei file dell'archivio.
  \item \textbf{Dentro Documentale}: \texttt{app:import-albo-corpus} importa gli
        atti con il modello dei documenti dell'app;
        \texttt{app:export-to-elastic-search} li indicizza in Elasticsearch, e lo
        script \texttt{strumenti/indicizza-koskidex.sh} fa lo stesso su Koskidex
        passando dall'app; \texttt{app:export-eval-corpus} esporta il corpus in
        formato BEIR per le valutazioni piatte.
\end{enumerate}

\section*{Le misure}

\begin{itemize}
  \item \textbf{Koskidex piatto}, collezioni pubbliche e albi:
        \texttt{go run ./scripts/evaluate} con le impostazioni da riga di comando
        (recupero, punteggio, analisi, vettori, fusione). Ogni file dell'archivio
        riporta nel blocco \texttt{config} le opzioni esatte.
  \item \textbf{Dall'app}: \texttt{app:eval-run-queries} esegue un file di query
        con la ricerca di produzione, su Elasticsearch o su Koskidex secondo
        \texttt{SEARCH\_BACKEND}; \texttt{scripts/evaluate -rankings} ne calcola
        le metriche con lo stesso codice di Koskidex piatto.
  \item \textbf{Gli esperimenti}: ogni cartella di
        \texttt{risultati/esperimenti/} ha un README con domanda, metodo,
        previsioni ed esito, e di solito uno script \texttt{analizza.py} (in qualche
        cartella con un altro nome) che legge i file
        dell'archivio e scrive un file di esito con il proprio commit. Il README
        cita i file di partenza per nome, e rifare l'esito vuol dire rilanciare lo
        script sullo stesso commit.
\end{itemize}

Il README di \texttt{risultati/} riporta i comandi esatti, i formati dei file e il
registro di ogni file tolto, rifatto o rinominato.

\chapter{Istruzioni di annotazione}\label{chapter:appendice-annotazione}
\sloppy

Le istruzioni sono state scritte il 25 settembre 2026, prima di darle a chiunque e
prima di annotare una riga, e stanno nel repository della tesi
(\texttt{istruzioni-annotazione.md}). Qui sono riportate nella sostanza. Una
regola che durante il lavoro si rivela sbagliata non si corregge a voce: si cambia
il file, con un commit che dice cosa cambia e da quale riga dell'annotazione vale.

Le parti 2 e 3 non sono state applicate come sono scritte: i bisogni e i
giudizi del pool li ha dati un modello linguistico, e io ho giudicato, nel ruolo
del secondo annotatore, un campione di 178 righe. La
sezione~\ref{sec:giudizi} racconta com'è andata e cosa ne segue.

\section*{Parte 1: raccogliere le query known-item}

\paragraph{Chi.} Almeno tre persone che non conoscano i motori e non abbiano visto
i risultati di una ricerca sul corpus. L'autore della tesi non scrive query: sa
quali parole trovano cosa, e le sue query misurerebbero lui.

\paragraph{Quali atti.} Un campione estratto a caso con un seme fissato: metà degli
atti da ciascuna fonte, e dentro ogni fonte per genere in proporzione; nessun atto
già usato dalle known-item automatiche; nessuno scelto a mano. La prima versione
diceva «in proporzione al corpus», e avrebbe dato nove atti di Crispiano su 160:
cambiata prima di qualunque raccolta.

\paragraph{Come, per ogni atto.}
\begin{enumerate}
  \item Si mostra l'atto come lo vedrebbe un impiegato: titolo e testo per
        Crispiano, la scheda per il Friuli Venezia Giulia. Senza identificativi,
        punteggi, o il nome del motore.
  \item La persona lo legge per il tempo che le serve, poi l'atto si nasconde.
  \item Le si chiede: \emph{«Fra qualche settimana ti serve di nuovo questo atto.
        Cosa scriveresti nella barra di ricerca per ritrovarlo?»}, e si registra la
        query esattamente com'è, refusi compresi.
  \item Si registra se la query contiene un numero.
\end{enumerate}
Non si suggerisce niente, non si chiede di riformulare, non si prova la query
davanti alla persona, non si scarta una query perché sembra brutta. «Non saprei
cosa scrivere» è un dato: si registra la query vuota con il motivo. Si registrano
anche lo pseudonimo della persona, mai il nome, e i secondi di lettura.

\section*{Parte 2: giudicare il pool}

\paragraph{Prima, il bisogno.} Per ogni query, prima di vedere un solo risultato,
si scrive in una o due frasi cosa cerca chi la digita nel documentale di un comune,
e cosa non conta. Le descrizioni si committano prima di aprire il foglio: senza, il
bisogno diventa «quello che i motori hanno trovato».

\paragraph{Il foglio.} Un atto alla volta, con il bisogno sopra, nell'ordine
dell'identificativo e non della posizione, senza dire quale motore l'ha trovato;
il testo intero a richiesta, e nel dubbio si apre sempre. Sessioni di un'ora al
massimo.

\paragraph{La scala.}

\begin{center}
\small
\begin{tabular}{p{1.2cm}p{4.4cm}p{5.8cm}}
  \toprule
  grado & significato & esempio, per \emph{manutenzione strade} \\
  \midrule
  2 & risponde al bisogno per intero & la determina che affida la manutenzione
  ordinaria delle strade comunali \\
  1 & pertinente ma parziale & la liquidazione di un singolo intervento;
  l'ordinanza di chiusura di una via per quei lavori \\
  0 & non risponde, anche se contiene le parole della query & un avviso di
  concorso che nomina la manutenzione fra le mansioni \\
  \bottomrule
\end{tabular}
\end{center}

\paragraph{Casi dubbi.} Si giudica l'atto, non le parole. Un atto e la sua
rettifica si giudicano ciascuno per conto proprio. Un comune diverso da quello che
si ha in mente non abbassa il grado, a meno che la query nomini un comune. Un
grado vuoto vuol dire «non l'ho guardato», mai «non pertinente». Nel dubbio fra
due gradi, il più basso, con una nota.

\paragraph{Dopo.} Un giudizio non si cambia dopo aver visto le metriche: un errore
trovato dopo si corregge in un file nuovo, con la riga e il motivo nel registro
dell'archivio.

\section*{Parte 3: il secondo annotatore e il kappa}

Si estraggono con un seme alcune query intere dal foglio già annotato, almeno 150
righe: query intere, perché ogni giudizio si dà rispetto a un bisogno. Il secondo
annotatore riceve i bisogni, queste istruzioni e un foglio vuoto, e non vede i
giudizi del primo né ne parla con lui finché non ha finito. Si calcolano il kappa
di Cohen sui tre gradi, il kappa pesato linearmente e l'accordo semplice, e si
riportano comunque. Solo dopo i due discutono i disaccordi; le metriche della tesi
usano i giudizi del primo annotatore, e si dice di quanto cambiano con quelli del
secondo e con quelli concordati.

\section*{Cosa non si fa mai}

Annotare dopo aver visto le metriche di quella collezione. Giudicare in ordine di
ranking, o sapendo quale motore ha trovato il documento. Correggere un giudizio in
silenzio. Scegliere a mano gli atti da mostrare o le query da tenere. Mettere in un
file dell'archivio il nome di chi ha scritto una query.

\chapter{Registro delle ipotesi}\label{chapter:appendice-ipotesi}

Ogni esperimento di questa tesi ha scritto le sue previsioni prima della misura,
con una soglia numerica, e le ha committate nel README dell'esperimento o nella
voce del diario di Koskidex; a misura fatta ha aggiunto l'esito, previsione per
previsione, senza toccare quello che era scritto prima. La
tabella~\ref{tab:ipotesi} le raccoglie tutte, quelle già misurate: la prima colonna dice da dove
viene la previsione (la cartella in \texttt{risultati/esperimenti/} senza la
data, o la voce del diario), la seconda la riassume in una riga con la sua
soglia, la terza riporta il numero misurato e il verdetto, confermata, smentita
o a metà quando una previsione in più parti ha tenuto solo in alcune. Il testo
completo di ogni previsione, con il ragionamento che la sosteneva, sta nella
fonte citata, committato prima della misura; le poche aggiunte dopo una prima
analisi, ma prima della misura che giudicano, sono segnate come tali. Dove la
fonte registra il numero ma non il verdetto, o un verdetto che non tiene conto
di tutte le parti della previsione, il verdetto è mio, ottenuto confrontando il
numero con la soglia, ed è segnato con un asterisco.

{\small
\begin{longtable}{@{}>{\RaggedRight\arraybackslash}p{3.1cm}>{\RaggedRight\arraybackslash}p{6.1cm}>{\RaggedRight\arraybackslash}p{3.6cm}@{}}
  \caption{Le previsioni di ogni esperimento, scritte prima della misura, e il
  loro esito.}\label{tab:ipotesi}\\
  \toprule
  esperimento o voce & previsione & esito \\
  \midrule
  \endfirsthead
  \toprule
  esperimento o voce & previsione & esito \\
  \midrule
  \endhead
  Diario, pari merito
    & 1. \texttt{d1} e \texttt{d3} si fermano su \texttt{[d1 d3]}
    & 40 esecuzioni su 40: confermata \\
    & 2. Nessun altro caso congelato cambia
    & test congelati verdi 20 volte: confermata \\
    & 3. Venti esecuzioni consecutive danno sempre lo stesso ordine
    & confermata (le 40 esecuzioni della 1) \\
  \addlinespace
  Diario, baseline
    & 1. SciFact nettamente sotto 0,6789, fra 0,30 e 0,55
    & 0,0246, sotto la fascia: smentita$^{*}$ \\
    & 2. NFCorpus sotto 0,3218, con un distacco assoluto minore
    & 0,1659: confermata$^{*}$ \\
    & 3. Recall@100 molto più alto di nDCG@10
    & SciFact 0,0242 contro 0,0246: smentita \\
    & 4. Nessuna query saltata su SciFact (339 giudizi a rilevanza 1)
    & 0 saltate, 339 giudizi a rilevanza 1: confermata$^{*}$ \\
  \addlinespace
  Diario, difetto 0
    & 1. Le query a vuoto crollano quasi a zero
    & 0 su SciFact, 24 su 323 su NFCorpus: confermata \\
    & 2. SciFact nDCG@10 fra 0,10 e 0,45
    & 0,4936: smentita \\
    & 3. SciFact Recall@100 sopra 0,50
    & 0,7571: confermata \\
    & 4. NFCorpus: richiamo su, nDCG@10 può peggiorare
    & migliora su tutte e tre le metriche: smentita \\
    & 5. \texttt{TestBaselineRankingIsFrozen} verde a default
    & confermata \\
  \addlinespace
  Diario, difetto 1
    & 1. SciFact nDCG@10 fra 0,55 e 0,70
    & 0,6197: confermata \\
    & 2. NFCorpus nDCG@10 fra 0,25 e 0,33
    & 0,2810: confermata \\
    & 3. Guardia: Recall@100 fermo entro $\pm$0,02
    & +0,1175 e +0,038: smentita, guardia posta male \\
    & 4. \texttt{TestBaselineIdentifier\allowbreak TieBreaksByDocID} fallisce
    & mal posta: gira a default \texttt{legacy} e non poteva fallire; il
      contenuto era giusto: a metà$^{*}$ \\
    & 5. \texttt{TestBaselineRankingIsFrozen} verde a default
    & tre esecuzioni: confermata \\
  \addlinespace
  Diario, stopword e stemmer (E1)
    & 1. SciFact fra 0,63 e 0,70
    & 0,6585 con tutte e due (0,6641 con le sole stopword): confermata \\
    & 2. NFCorpus fra 0,29 e 0,34
    & 0,2941 con tutte e due: confermata \\
    & 3. Lo stemmer porta molto più delle stopword, le stopword meno di +0,01
    & SciFact: stopword +0,0444, stemmer +0,0046: smentita, ribaltata \\
    & 4. I candidati aumentano
    & su con lo stemmer, da 4833 a 2522 con le stopword: smentita \\
    & 5. Guardia: le query a vuoto non aumentano
    & 0 e 24 invariate, 15 con lo stemmer: confermata \\
    & 6. Delle 24 query vuote di NFCorpus se ne recupera una
    & nove: smentita \\
    & 7. \texttt{TestBaselineRankingIsFrozen} verde a default
    & esito non registrato \\
  \addlinespace
  \texttt{bm25-numeri}
    & 1. In più di metà dei fallimenti il primo classificato non contiene il
      numero
    & 107 su 115: confermata nel numero, non nella spiegazione \\
    & 2. I numeri delle query fallite hanno una df mediana più alta
    & 17 contro 6: confermata \\
    & 3. $b = 0$ porta MRR@10 da 0,709 a oltre 0,84
    & 0,731: smentita \\
    & 4. $k_1 = 0{,}3$ aiuta meno di $b = 0$
    & 0,900: smentita \\
  \addlinespace
  Diario, frequenza mescolata
    & 1. Known-item MRR@10 sopra 0,85
    & 0,833: smentita di poco \\
    & 2. Known-item sotto l'euristico
    & 0,833 contro 0,975: confermata \\
    & 3. SciFact e NFCorpus cambiano meno di 0,01 di nDCG@10
    & SciFact +0,050: smentita, nel verso buono \\
    & 4. Gli insiemi non cambiano
    & candidati identici query per query: confermata \\
  \addlinespace
  \texttt{known-item-\allowbreak auto}
    & 1. Elasticsearch: atto entro 10 in meno del 20\%, più di metà a vuoto
    & 2,0\% e 70\%: confermata \\
    & 2. Koskidex innestato: stessi insiemi di Elasticsearch su almeno il
      95\%
    & 300 su 300: confermata \\
    & 3. Koskidex piatto: atto entro 10 in più del 70\%, tutte e tre le
      configurazioni
    & 99,7\%, 99,7\% e 90,3\%: confermata \\
    & 4. BM25 ed euristico a meno di 0,05 di MRR@10
    & 0,709 contro 0,975: smentita \\
  \addlinespace
  \texttt{known-item-\allowbreak divario}
    & 1. Frequenza mescolata: in più di metà dei fallimenti vince un numero
      trovato per prefisso
    & 11 su 69: smentita \\
    & 2. Prefisso spento: in più di due terzi dei fallimenti il primo non ha
      il numero
    & 38 su 60 (63\%): smentita di poco \\
    & 3. In quei casi il primo ripete il comune più dell'atto, in più di due
      terzi
    & 37 su 38: confermata \\
    & 4. (aggiunta dopo la prima analisi) \texttt{all} e frequenza mescolata:
      atto primo in almeno 255, MRR@10 almeno 0,90
    & 281 e 0,960: confermata \\
    & 5. (aggiunta dopo la prima analisi) \texttt{all} e prefisso spento:
      almeno 255 e 0,90
    & 288 e 0,976: confermata \\
  \addlinespace
  \texttt{recupero-\allowbreak intermedio}
    & 1. \texttt{minimum\_should\_match}: known-item MRR@10 almeno 0,93
    & 0,889: smentita \\
    & 2. \texttt{minimum\_should\_match}: SciFact giù di più di 0,02
    & da 0,6694 a 0,2556: confermata \\
    & 3. \texttt{minimum\_should\_match}: NFCorpus giù di meno di 0,02
    & giù di 0,078: smentita \\
    & 4. Coordinazione: known-item MRR@10 almeno 0,93
    & 0,944: confermata \\
    & 5. Coordinazione: SciFact e NFCorpus entro 0,01 dalla base
    & $-$0,027 e $-$0,011: smentita \\
  \addlinespace
  \texttt{stopword-\allowbreak espansioni}
    & 1. Senza prefisso, le stopword portano SciFact meno di 0,01
    & +0,003: confermata \\
    & 2. Più di metà delle query di SciFact ha una stopword espandibile
    & 271 su 300: confermata \\
    & 3. Le query senza portano meno del 10\% del guadagno
    & 0,7\%: confermata \\
  \addlinespace
  \texttt{scelta-per-\allowbreak query}
    & 1. Oracolo: known-item entro 0,02 da sempre B; SciFact e NFCorpus
      almeno 0,01 sopra sempre A
    & uguale a B; 0 e 0,0007: a metà \\
    & 2. Regola della cifra: uguale a B sulle known-item, oltre 0,05 sotto A su
      SciFact, entro 0,005 su NFCorpus
    & $-$0,259 e $-$0,004: confermata \\
    & 3. Ripiego: uguale a B sulle known-item, oltre 0,01 sotto A su SciFact e
      NFCorpus
    & NFCorpus $-$0,017, SciFact $-$0,003: a metà \\
    & 4. Classificatore: entro 0,01 da B sulle known-item ed entro 0,005 da A
      sulle altre, meglio della cifra su SciFact
    & 0,9687, 0,6694, 0,3038: confermata \\
  \addlinespace
  \texttt{elisioni}
    & 1. Atti con un termine eliso fra il 30\% e il 60\%
    & 42,7\%: confermata \\
    & 2. La ricerca salva meno del 30\% delle coppie a rischio
    & 2,3\%: confermata \\
    & 3. Il filtro \texttt{elision} recupera più del 90\% delle coppie perse
    & 99,4\%: confermata \\
    & 4. Almeno cinque delle dieci parole previste fra le prime venti
    & sei: confermata \\
    & 5. Più di un atto su dieci con una coppia persa
    & 40,3\%: confermata \\
  \addlinespace
  \texttt{elisioni-\allowbreak koskidex}
    & 1. \texttt{elis0} e produzione con lo stesso insieme per almeno 500
      parole su 507
    & 486: smentita \\
    & 2. Coppie perse da \texttt{elis0} quelle di Elasticsearch, 5.775
      $\pm$ 58
    & le stesse 5.775: confermata \\
    & 3. \texttt{elis1} recupera almeno il 99\% delle coppie perse
    & 99,4\%, le stesse 36 residue: confermata \\
    & 4. \texttt{elis1} e controllo con lo stesso insieme per almeno 480
      parole
    & 485: confermata \\
    & 5. Documenti in più fra 4.600 e 5.700
    & 5.200: confermata \\
    & 6. Nessuna delle 24 query perde un documento, almeno una ne guadagna
    & 0 persi, 65 guadagnati in 8 query: confermata \\
    & 7. Known-item, parole libere: MRR@10 cambia di meno di 0,01
    & 0,9498 in entrambe: confermata \\
  \addlinespace
  \texttt{koskidex-campi-\allowbreak liberi}
    & 1. Atto entro 10 oltre il 90\%, query a vuoto sotto il 5\%
    & 93,3\% e zero: confermata \\
    & 2. Atto primo oltre il 60\%
    & 38,3\%: smentita \\
    & 3. Sulle 24 query gli insiemi solo crescono, almeno un terzo cambia
    & nessun documento perso, 14 su 24: confermata \\
  \addlinespace
  \texttt{refusi-numeri}
    & 1. Parole libere: atto primo oltre l'80\% e non oltre il 92\%, MRR@10
      oltre 0,85
    & 91,7\% e 0,950: confermata \\
    & 2. Atto entro 10 non sotto il 93,3\%, nessuna query a vuoto,
      Recall@100 almeno 0,99
    & 99,7\%, zero, 1,0: confermata \\
    & 3. Delle 24 cambiano solo le query con un numero, e si restringono;
      \texttt{doc-0001} e \texttt{doc-0002} primi
    & cambiano due, la terza è vuota comunque: confermata \\
    & 4. Campo unico: atto entro 10 non oltre il 3\%, a vuoto almeno il 70\%
    & 2\% e 95\%: confermata nei numeri \\
  \addlinespace
  \texttt{ibrido-unione}
    & 1. Re-ranking almeno +0,02 su SciFact e +0,01 su NFCorpus
    & +0,0219 e +0,0235: confermata \\
    & 2. Re-ranking almeno 0,05 sotto sulle known-item
    & +0,0133: smentita \\
    & 3. Unione almeno +0,01 di Recall@100 su SciFact e NFCorpus, nDCG@10
      entro 0,01
    & SciFact invariato, NFCorpus +0,031: a metà \\
    & 4. Known-item: unione entro 0,01 dal re-ranking
    & identica: confermata \\
    & 5. Solo vettori: SciFact fra 0,55 e 0,72, NFCorpus fra 0,28 e
      0,37, known-item sotto 0,2
    & 0,6436, 0,3149, 0,1837: confermata \\
    & 6. Costo dell'unione: meno di 10 ms su SciFact e NFCorpus, 20 sulle
      known-item
    & +6,5, +5,1 e +15,3 ms: confermata \\
  \addlinespace
  \texttt{fusione}
    & 1. Costante migliore diversa di almeno quattro volte fra SciFact e
      NFCorpus
    & 160 su entrambe: smentita \\
    & 2. $\alpha$ migliori entro 0,2
    & 0,5 su entrambe: confermata \\
    & 3. Known-item: $\alpha$ almeno 0,8, costante al più 10
    & 0,9 e 10: confermata \\
    & 4. Trasferire la costante costa più che trasferire $\alpha$
    & nessuna delle due costa: smentita \\
    & 5. RRF almeno +0,005 su SciFact e NFCorpus, almeno 0,05 sotto sulle
      known-item
    & $-$0,0047, +0,0107, $-$0,384: a metà \\
    & 6. Convessa calibrata almeno +0,005 sulle pubbliche, non oltre 0,005
      sotto sulle known-item, almeno pari a RRF
    & +0,0112, +0,0076, +0,0001; su NFCorpus 0,0031 sotto RRF: a metà$^{*}$ \\
    & 7. Le fusioni nuove costano meno di 2 ms in più dell'unione
    & costano meno, da 1 a 7 ms: confermata \\
  \addlinespace
  \texttt{ibrido-\allowbreak congiuntivo}
    & 1. SciFact: unione almeno 0,60, re-ranking entro 0,01 dal lessicale
    & 0,6483 e +0,000: confermata \\
    & 2. NFCorpus: unione almeno 0,30 senza query a vuoto, re-ranking meno di
      +0,02
    & 0,3333, zero, +0,0035: confermata \\
    & 3. Known-item: re-ranking entro 0,01 dal lessicale, unione entro 0,01
      dal re-ranking
    & $-$0,0017 e +0,000: confermata \\
    & 4. Unione almeno 0,5 sopra il re-ranking su SciFact
    & +0,6236: confermata \\
    & 5. Unione con \texttt{all} più veloce che con \texttt{any}
    & 8, 1,2 e 1,2 ms in meno: confermata \\
  \addlinespace
  \texttt{scelta-ibrida}
    & 1. Regola a due condizioni almeno 0,665 di media, almeno +0,015 sulle
      due configurazioni di riferimento
    & 0,6728, +0,023 e +0,026: confermata \\
    & 2. Sempre A sulle known-item sotto 0,65
    & 0,5589: confermata \\
    & 3. Sempre B: SciFact entro 0,02 dal solo vettore, NFCorpus fra 0,30 e
      0,34
    & 0,6483 e 0,3325: confermata \\
    & 4. Oracolo almeno 0,03 sopra sempre A su SciFact
    & +0,021: smentita \\
    & 5. Classificatore entro 0,005 dalla regola
    & $-$0,0012: confermata \\
  \addlinespace
  \texttt{costo-vettori}
    & 1. Allocazioni per ricerca giù di almeno il 90\%
    & in numero $-$9\% e $-$17\%: smentita \\
    & 2. Unione almeno tre volte più veloce, re-ranking almeno due
    & 2,40$\times$ e 2,33$\times$: a metà \\
    & 3. Memoria viva giù di almeno il 60\%
    & $-$65\%: confermata \\
    & 4. Documentale guadagna la parte del motore dei 24 ms in più
    & nessuna differenza fra i due binari dal vivo: i 23~ms erano il vettore
      della query: smentita$^{*}$ \\
    & 5. (seconda modifica) Re-ranking non oltre 18 ms, unione non oltre 30
    & 24,62 e 39,60: smentita \\
    & 6. (seconda modifica) Heap vivo cresce di meno di 1 MB
    & +0,4 MB: confermata \\
  \addlinespace
  \texttt{configurazione-\allowbreak consigliata}
    & 1. Impostazioni dell'indice: campi liberi, numeri esatti, i 21 articoli
      elisi
    & confermata \\
    & 2. Known-item: MRR@10 entro 0,005 da 0,950, atto primo in 275
      $\pm$ 3, entro 10 in almeno 298
    & 0,9498, 275, 299: confermata \\
    & 3. Nessuna delle 24 query perde un documento, almeno una ne guadagna
    & 0 perse, 7 guadagnano: confermata \\
  \addlinespace
  Diario, contesto di Ollama
    & 1. Con \texttt{context} a zero la richiesta resta identica e il nome del
      modello non cambia
    & test verde: confermata$^{*}$ \\
    & 2. Con \texttt{context} 8.192 un testo oltre 2.048 token ha un vettore
      diverso
    & vettori identici: smentita \\
    & 3. \texttt{TestBaselineRankingIsFrozen} passa senza modifiche
    & confermata$^{*}$ \\
    & 4. (dopo la 2) Con \texttt{num\_ctx} e \texttt{num\_batch} Ollama legge
      8.192 token
    & confermata \\
    & 5. (dopo la 2) A zero la richiesta resta \texttt{model} e
      \texttt{input}
    & test verde: confermata \\
  \addlinespace
  \texttt{contesto-ollama}
    & 1. Con il contesto predefinito Ollama legge al più 2.048 token$^{\dagger}$
    & confermata \\
    & 2. Con \texttt{num\_ctx} 8.192 legge fino a 8.192 token
    & 2.048 come senza: smentita \\
    & 3. Fra il 30\% e il 70\% dei testi integrali di Crispiano supera 2.048
      token
    & 334 su 563 (59\%): confermata \\
    & 4. Meno del 5\% dei testi integrali arriva a 8.192 token
    & 3 su 563: confermata \\
    & 5. Nessuna scheda supera 2.048 token
    & la più lunga 486: confermata \\
    & 6. (aggiunta) Meno dell'1\% dei documenti di SciFact e NFCorpus supera
      2.048 token
    & 3 e 2 (0,06\%): confermata \\
  \addlinespace
  \texttt{carico}
    & 1. Una alla volta, p50 di Koskidex nel container sotto quello di
      Elasticsearch come l'app
    & 2,12 contro 6,32~ms: confermata \\
    & 2. Da 1.000 risultati, Elasticsearch come l'app almeno il doppio che
      senza il testo
    & 27,5 contro 22,4~ms, su 5 richieste: smentita \\
    & 3. Elasticsearch senza il testo e Koskidex entro un fattore 2
    & rapporto 3,2: smentita \\
    & 4. Con la cache, Koskidex sotto 1~ms di p50
    & 0,29~ms: confermata \\
    & 5. Capacità di Koskidex almeno 1,5 volte Elasticsearch
    & 1,14 volte: smentita \\
    & 6. Sotto carico Elasticsearch fra 1,3 e 1,8~GB, Koskidex al più 400~MB
    & 1.566 e 355~MB: confermata \\
    & 7. A 16 client p99 di Koskidex sotto 100~ms
    & 105~ms: smentita \\
    & 8. Il nativo al più il 30\% sotto il container
    & 16\%: confermata \\
    & 9. A 100.000 documenti memoria di Koskidex almeno $\times$5 ed
      Elasticsearch sotto $\times$1,5; p50 di Koskidex almeno $\times$3,
      Elasticsearch sotto $\times$3
    & memoria $\times$7,85 e $\times$1,00; p50 $\times$2,57 e $\times$1,10: a
      metà \\
    & 10. Fra le prime cinque funzioni per CPU l'espansione dei refusi e la
      scrittura del JSON
    & la prima sì, la seconda no: a metà \\
  \addlinespace
  \texttt{prestazioni}
    & 1. Nessun risultato cambia: impronte a 10.018 e 100.000 documenti,
      metriche di \texttt{evaluate}
    & 429 e 429 impronte uguali, 9 valutazioni su 9: confermata \\
    & 2. Allocazioni per ricerca almeno il 60\% in meno
    & da 2,57 a 0,44~MB: confermata \\
    & 3. Una alla volta nel container p50 al più 1,5~ms e p95 al più 6
    & 1,12 e 2,74~ms: confermata \\
    & 4. Almeno 900 ricerche al secondo nel container
    & 2.614: confermata \\
    & 5. A 32 client p99 nel container al più 110~ms
    & 43~ms: confermata \\
    & 6. A 100.000 documenti p95 al più 40~ms e p99 al più 60
    & 10,4 e 15,5~ms: confermata \\
    & 7. Memoria a riposo entro il 10\%, sotto carico al più 300~MB
    & $-$0,2\% e 272~MB: confermata \\
    & 8. Con un client CPU del container al più 100\%
    & 83\%: confermata \\
  \addlinespace
  \texttt{giudice-llm}
    & 1. Kappa pesato fra modello e annotatore umano di almeno 0,60
    & 0,370: smentita \\
    & 2. Accordo semplice di almeno il 70\%
    & 54,5\%: smentita \\
    & 3. Meno del 10\% delle righe con un disaccordo fra 0 e 2
    & 11,2\%: smentita \\
    & 4. Il modello più generoso dell'annotatore umano
    & più severo: smentita \\
    & 5. c19 con l'accordo più alto delle quattro query
    & 80\%: confermata \\
  \addlinespace
  \texttt{ordine-fisso}
    & 1. Spento, il baseline non varia in cinque ripetizioni
    & nessuna query varia: confermata \\
    & 2. Spento, con BM25 al più 5 query variano e le medie coincidono alla
      quarta cifra
    & nessuna query varia: confermata \\
    & 3. Acceso, nessuna query varia
    & confermata \\
    & 4. Acceso contro spento, nDCG@10 entro 0,001, baseline uguale
    & SciFact BM25 \texttt{any} $-$0,0024: smentita \\
    & 5. Documentale, le 429 impronte non cambiano
    & 429 su 429: confermata \\
    & 6. Mediana del tempo per query entro il 5\%
    & fra $-$1,7\% e +3,2\%: confermata \\
    & 7. I test passano; con l'impostazione il test del percorso veloce
      confronta tutte le query
    & 4.200 su 4.200: confermata \\
  \addlinespace
  \texttt{date-importi}
    & 1. Tokenizer standard: data nel formato dell'atto almeno 90\% entro
      dieci, in un altro formato al più 10\%
    & 98-100\% e 0-8\%: confermata \\
    & 2. Tokenizer di Koskidex: barra e punto almeno 90\%, col mese in
      lettere al più 10\%
    & 94\% e 0-8\%: confermata \\
    & 3. Acceso: ogni coppia almeno 90\%, atto primo entro 5 punti dal
      controllo
    & 98-100\%, uguale al controllo: confermata \\
    & 4. Importi: spento fra formati al più 10\%, acceso almeno 90\%
    & 0\% e 100\%: confermata \\
    & 5. Known-item automatiche, MRR@10 entro 0,005
    & da 0,9750 a 0,9767 e 0,9772: confermata \\
    & 6. Numeri da 1 a 31: candidati mai in aumento, in calo almeno per metà
    & in calo per 4 su 12: smentita \\
    & 7. Indicizzazione accesa al più il 20\% più lenta
    & +41-46\%: smentita \\
    & 8. I test passano
    & confermata \\
    & 9. (dopo l'esito) Saltando i testi senza il formato, nessun risultato
      cambia
    & 12 configurazioni su 12: confermata \\
    & 10. (dopo l'esito) Indicizzazione accesa al più il 15\% più lenta
    & +29-34\%: smentita \\
  \addlinespace
  \texttt{espansioni-sinonimo}
    & 1. Con \texttt{synonym} le metriche non dipendono dall'ordine
    & 8 configurazioni su 8: confermata \\
    & 2. Stessi candidati di \texttt{blended}
    & 8 su 8: confermata \\
    & 3. SciFact \texttt{any}, nDCG@10 fra $-$0,005 e +0,02
    & $-$0,0028 e $-$0,0023: confermata, ma nel verso opposto al previsto \\
    & 4. NFCorpus \texttt{any} entro 0,01
    & $-$0,0007 e $-$0,0016: confermata \\
    & 5. Known-item automatiche, MRR@10 entro 0,01
    & 0 e +0,0025: confermata \\
    & 6. Mediana del tempo per query entro il 5\%
    & fra $-$1,4\% e +2,8\%: confermata \\
  \addlinespace
  \texttt{allocazioni}
    & 1. Nessun risultato cambia: impronte a 10.018 e 100.000 documenti,
      metriche di \texttt{evaluate}
    & 429 e 429 impronte uguali, 9 valutazioni su 9: confermata \\
    & 2. Allocazioni per ricerca almeno il 60\% in meno
    & da 0,435 a 0,081~MB: confermata \\
    & 3. La raccolta dei candidati con refuso sotto il 10\% dei byte allocati
    & dall'86\% al 30\%: smentita \\
    & 4. Capacità nativa almeno +5\%
    & +56\%: confermata \\
    & 5. Capacità nel container almeno +5\%
    & +44\%: confermata \\
    & 6. A 100.000 documenti p95 e p99 almeno $-$15\%
    & $-$4\% e $-$7\%: smentita \\
    & 7. Una alla volta nel container, p50 non peggiore di oltre il 5\%
    & $-$19\%: confermata \\
  \addlinespace
  \texttt{date-app}
    & 1. Dall'app, spento: data nel formato dell'atto almeno 90\% entro dieci,
      in un altro al più 10\%
    & 98-100\% e 0-8\%: confermata \\
    & 2. Acceso: ogni coppia di formati almeno 90\%
    & 98-100\%: confermata \\
    & 3. Importi fra formati: acceso almeno 90\%, spento al più 10\%
    & 100\% acceso, ma 100\% anche spento: a metà \\
    & 4. Known-item automatiche, MRR@10 entro 0,005
    & +0,0053: smentita \\
    & 5. Known-item umane: nessuna entra o esce dai primi dieci, MRR@10 entro
      0,01
    & nessuna cambia: confermata \\
    & 6. Le 24: senza cifre stessi atti, con l'anno risultati solo in calo
    & 21 su 21; zero risultati in tutti e due: confermata \\
    & 7. Indicizzazione dall'app accesa al più il 10\% più lenta
    & +6,4\%: confermata \\
    & 8. I test di Documentale passano
    & confermata \\
  \addlinespace
  \texttt{importi-esatti}
    & 1. Importi senza refusi: fra formati al più 10\% entro dieci
    & 0\%: confermata \\
    & 2. Con anche la normalizzazione: fra formati almeno 90\%, MRR@10 degli
      importi almeno 0,95
    & 100\% e 0,990: confermata \\
    & 3. Risultati sulle query degli importi almeno il 20\% in meno
    & da 175 a 106, $-$39\%: confermata \\
    & 4. Date, MRR@10 entro 0,01
    & +0,018: smentita, in meglio \\
    & 5. Known-item automatiche identiche query per query
    & 300 su 300: confermata \\
    & 6. Known-item umane identiche query per query
    & 76 su 76: confermata \\
    & 7. Le 21 del confronto senza cifre identiche
    & confermata \\
    & 8. I test passano
    & confermata \\
  \addlinespace
  \texttt{accenti}
    & 1. Nessun risultato cambia: impronte, \texttt{evaluate}, test
    & 429 e 429, 9 su 9: confermata \\
    & 2. Allocazioni per ricerca almeno $-$25\%
    & da 0,078 a 0,052~MB, $-$33\%: confermata \\
    & 3. La catena degli accenti sotto 1~KB a ricerca
    & da 27,3 a 0,014~KB: confermata \\
    & 4. Capacità nativa fra $-$3\% e +15\%
    & +1,7\%: confermata \\
    & 5. Una alla volta nel container, p50 non peggiore di oltre il 5\%
    & $-$4\%: confermata \\
  \addlinespace
  \texttt{reranker}
    & 1. SciFact, nDCG@10 riordinato almeno 0,72
    & da 0,676 a 0,723: confermata, di poco \\
    & 2. NFCorpus, nDCG@10 riordinato almeno 0,33
    & da 0,306 a 0,331: confermata, di poco \\
    & 3. Known-item automatiche: LA riordinato resta sotto LT
    & 0,953 contro 0,960: confermata, ma di 0,007 \\
    & 6. Mediana del riordino almeno 2~s a query
    & da 5,1 a 12,0~s: confermata \\
  \bottomrule
\end{longtable}
}

In tutto le previsioni sono 193: 137 confermate, 45 smentite, 10 a metà, e una
senza un esito registrato (la settima della voce E1). Delle 192 con un esito,
più di una su quattro non ha tenuto del tutto. La prima di \texttt{contesto-ollama},
segnata con $\dagger$, è scritta dopo una prova di sviluppo non archiviata che
già la indicava, e il README lo dichiara. Gli esperimenti \texttt{cause-divergenza},
\texttt{numero-atto}, \texttt{costo-sottostringa} e \texttt{innesto-parita} non
hanno previsioni numeriche: il primo registrava due ipotesi sulle cause della
divergenza residua, scritte senza soglia e sbagliate tutte e due. Le quattro
ipotesi scritte nel diario il 23 settembre sul confronto fra scheda e testo
intero aspettano ancora la misura (sezione~\ref{sec:scheda-testo}). Aspettano
la misura, e quindi mancano dalla tabella, anche le previsioni già scritte e
committate degli esperimenti che girano a raccolta delle known-item umane
chiusa: \texttt{known-item-umane} (7), \texttt{italiano-umane} (4),
\texttt{scelta-umane} (4), \texttt{scheda-testo} (5) e la parte umana di
\texttt{reranker} (2); entreranno con il loro esito.

\chapter{Le impostazioni aggiunte a Koskidex}\label{chapter:appendice-impostazioni}

Il 23 settembre 2026 le impostazioni di un indice di Koskidex erano otto: i campi
cercabili e mostrati, le regole di ordinamento, le stopword, i sinonimi, la
tolleranza ai refusi, i pesi dei campi e la mappa del sito. La
tabella~\ref{tab:impostazioni} elenca quelle aggiunte per questa tesi, con il
nome con cui compaiono nell'API. Ognuna ha come valore predefinito il
comportamento di partenza, che la tabella scrive per primo: un indice che non le
imposta si comporta come quel giorno, e \texttt{TestBaselineRankingIsFrozen} lo
controlla: fa parte di \texttt{make check}, che ho lanciato prima di ogni
commit (è una regola di lavoro, non un controllo automatico).

{\footnotesize\renewcommand{\arraystretch}{0.92}
\begin{longtable}{>{\raggedright\arraybackslash}p{4.1cm}>{\raggedright\arraybackslash}p{2.5cm}>{\raggedright\arraybackslash}p{6cm}}
  \caption{Le impostazioni aggiunte a Koskidex, con il valore predefinito per
  primo.}\label{tab:impostazioni}\\
  \toprule
  impostazione & valori & cosa fa \\
  \midrule
  \endfirsthead
  \toprule
  impostazione & valori & cosa fa \\
  \midrule
  \endhead
  \bottomrule
  \endfoot
  \multicolumn{3}{l}{\emph{Recupero e punteggio lessicale}
  (capitolo~\ref{chapter:lessicale})} \\
  \texttt{retrieval\_mode} & \texttt{all}, \texttt{any} & recupero congiuntivo
  o disgiuntivo (sezione~\ref{sec:difetto0}) \\
  \texttt{scoring\_mode} & \texttt{legacy}, \texttt{bm25} & punteggio euristico
  o BM25 (sezione~\ref{sec:difetto1}) \\
  \texttt{bm25\_k1}, \texttt{bm25\_b} & 1,2 e 0,75 & i parametri di BM25, quelli
  predefiniti di Lucene; valgono per un indice creato con le impostazioni
  predefinite, mentre un aggiornamento dall'API che non li passa li porta a
  zero (Documentale non usa BM25 e non ne è toccato) \\
  \texttt{bm25\_expansion} & vuoto, \texttt{blended}, \texttt{synonym} &
  frequenza documentale di un'espansione: del termine trovato, o la più alta fra
  i termini trovati dallo stesso termine della query; con \texttt{synonym}
  anche la frequenza nel documento è la somma delle espansioni, come la
  \emph{SynonymQuery} di Lucene (sezione~\ref{sec:espansioni}) \\
  \texttt{stable\_term\_order} & \texttt{false}, \texttt{true} & i termini
  trovati da una parola della query si visitano in ordine lessicografico, non in
  quello della mappa del vocabolario (sezioni~\ref{sec:espansioni}
  e~\ref{sec:latenza}) \\
  \texttt{minimum\_should\_match} & vuoto, specifica & in recupero disgiuntivo,
  quanti termini deve contenere un documento, con la sintassi di Elasticsearch
  (sezione~\ref{sec:nessuna-unica}) \\
  \texttt{coordination} & \texttt{false}, \texttt{true} & moltiplica il
  punteggio per la quota di termini della query presenti, come il \emph{coord} di
  Lucene prima della versione~7 \\
  \midrule
  \multicolumn{3}{l}{\emph{Analisi lessicale}
  (sezioni~\ref{sec:stopword} e~\ref{sec:italiana})} \\
  \texttt{stemmer} & vuoto, \texttt{porter}, \texttt{italian\_light} & nessuno
  stemmer, Porter per l'inglese, lo stemmer leggero di Savoy per l'italiano \\
  \texttt{elision\_articles} & vuoto, lista & gli articoli elisi da togliere
  davanti a un termine, come l'\texttt{ElisionFilter} di Lucene; la lista
  italiana è quella di Lucene \\
  \texttt{tokenizer} & vuoto, \texttt{standard} & spezzare su tutto ciò che non è
  lettera o cifra, o tenere la punteggiatura interna alle parole come il
  tokenizer standard di Elasticsearch \\
  \midrule
  \multicolumn{3}{l}{\emph{Compatibilità con Elasticsearch}
  (sezioni~\ref{sec:innesto} e~\ref{sec:differenze})} \\
  \texttt{all\_terms\_in\_one\_field} & \texttt{false}, \texttt{true} & con
  tutti i termini obbligatori, un campo del documento deve contenerli tutti, come
  \texttt{best\_fields} con \texttt{operator: and} \\
  \texttt{disable\_prefix\_search} & \texttt{false}, \texttt{true} & toglie la
  ricerca per prefisso: restano le corrispondenze esatte e con refusi \\
  \texttt{prefix\_length} & 0, intero & quanti caratteri iniziali di una
  corrispondenza con refusi devono essere esatti \\
  \texttt{substring\_match} & \texttt{false}, \texttt{true} & aggiunge i
  documenti con un termine che contiene l'intera query, come il
  \emph{wildcard} \texttt{*termine*} della ricerca delle cartelle \\
  \texttt{typo\_tolerance.} \texttt{disable\_on\_numbers} & \texttt{false},
  \texttt{true} & nessun refuso sui termini della query fatti di sole cifre, come
  in Meilisearch (sezione~\ref{sec:numero-comune}) \\
  \midrule
  \multicolumn{3}{l}{\emph{Date e importi} (sezione~\ref{sec:date})} \\
  \texttt{normalize\_dates} & \texttt{false}, \texttt{true} & le date valide in
  quattro formati, con i mesi in italiano, diventano un numero solo,
  \texttt{aaaammgg}, negli atti e nelle query \\
  \texttt{normalize\_amounts} & \texttt{false}, \texttt{true} & agli importi si
  tolgono i punti delle migliaia \\
  \texttt{typo\_tolerance.} \texttt{disable\_on\_amounts} & \texttt{false},
  \texttt{true} & nessun refuso sui termini della query fatti di cifre, punti e
  virgole, come \texttt{1.234,56} \\
  \midrule
  \multicolumn{3}{l}{\emph{Vettori e fusione} (capitolo~\ref{chapter:ibrido})} \\
  \texttt{embedder} & vuoto, \texttt{ollama} con modello e campi & il server
  calcola i vettori dei documenti e delle query con un modello locale
  (sezione~\ref{sec:embedding}) \\
  \texttt{embedder.} \texttt{context} & 0, intero & quanti token del testo legge
  il modello; con 0 decide Ollama, che ne legge 2.048
  (sezione~\ref{sec:embedding}) \\
  \texttt{hybrid\_mode} & vuoto, \texttt{union}, \texttt{vector} & chi entra fra i
  candidati: i lessicali, riordinati dal vettore; i lessicali più i più vicini;
  solo i più vicini (sezione~\ref{sec:difetto2}) \\
  \texttt{vector\_top\_k} & 100, intero & quanti documenti porta il vettore \\
  \texttt{fusion\_mode} & vuoto, \texttt{rrf}, \texttt{convex} & somma,
  Reciprocal Rank Fusion con $k = 60$, o combinazione convessa di punteggi
  normalizzati (sezione~\ref{sec:difetto3}) \\
  \texttt{vector\_weight} & 20, numero & la costante della somma \\
  \texttt{fusion\_alpha} & 0,5, numero & il peso del lessicale nella
  combinazione convessa \\
\end{longtable}
}

Due aggiunte non sono impostazioni ma liste da passare a \texttt{stop\_words},
che esisteva già: le stopword inglesi di Lucene e quelle italiane di Snowball.


%%%%%%%%%%%%%%%%%%%%%%%%%%%%%%%% Bibliografia
\pagestyle{plain}
\begin{thebibliography}{99}

\bibitem{DaCompilare}
Voce di esempio, da sostituire.

\end{thebibliography}


%%%%%%%%%%%%%%%%%%%%%%%%%%%%%%%% Ringraziamenti
\chapter*{Ringraziamenti}\label{chapter:ringraziamenti}
\addcontentsline{toc}{chapter}{Ringraziamenti}
\markboth{Ringraziamenti}{}

Questa tesi chiude un percorso che non ho fatto da solo, e ci tengo a ringraziare
chi lo ha reso più leggero.

Ai miei compagni di università, Leopoldo, Luigi Z., Luigi B., Aurora, Enrico
ed Elia, per aver condiviso con me lezioni, esami e preoccupazioni.

A Dieffetech, la mia azienda, per aver accompagnato questi anni di studio con il
lavoro, e in modo speciale a Francesco Desiderio, che la guida, per la fiducia
che mi ha dato.

Ai miei amici, Giovanni, Cris, Kleana, Federica, Yasmine, Vasco, Anastasija,
Tommaso e Imane, per esserci sempre stati, dentro e fuori dall'università.

A Giulia, la mia ragazza, per il sostegno di ogni giorno.

Alla mia famiglia, a mamma, papà e mio fratello, che mi hanno sempre sostenuto.

A Endri, il mio migliore amico, che c'è sempre stato: nei giochi, nella vita e in
ogni esperienza, anche quella organizzata all'ultimo minuto.

\newpage
\thispagestyle{empty}
\null\vspace{\stretch{1}}
\vspace{\stretch{3}}\null
\newpage

\blankpage
\end{document}
